\chapter{Preliminaries from Riemannian geometry}
\label{chap:prelim}
\label{sectprelim}

Throughout this chapter, repeated indices are summed (Einstein
convention), and a manifold is smooth, paracompact, and Hausdorff. Unless
specified otherwise, all connections are Levi--Civita connections.


\section{Riemannian metrics and the Levi-Civita connection}

For a manifold $M$ and $p\in M$, write $TM,T^*M$ for the tangent and
cotangent bundles, $T_pM,T_p^*M$ for their fibers, and
$\Gamma(\mathcal V)$ for the smooth sections of a vector bundle
$\mathcal V\to M$.

\begin{definition}[Riemannian metric]
\label{def:riemannian-metric}
\lean{MorganTianLib.RiemannianMetric}
\leanok


Let $M$ be an $n$-manifold. A \emph{Riemannian metric} is a smooth
section $g\in\Gamma(T^*M\otimes T^*M)$ whose value at each $p\in M$ is
a positive-definite symmetric bilinear form on $T_pM$. In coordinates
$(x^1,\ldots,x^n)$, with $\partial_i=\partial/\partial x^i$,
\[
  g=g_{ij}\,dx^i\otimes dx^j,
  \qquad g_{ij}=g(\partial_i,\partial_j),
\]
and $(g^{ij})=(g_{ij})^{-1}$. The pair $(M,g)$ is a \emph{Riemannian
manifold}.
\end{definition}

Every manifold admits a Riemannian metric by a partition-of-unity
construction. The associated distance is the infimum of lengths of smooth
paths joining two points. For $r>0$, set
\[
  B_g(p,r)=B(p,r)=\{q\in M\mid d_g(p,q)<r\}.
\]

\begin{theorem}[Levi-Civita connection]
\label{thm:levi-civita-connection}
\lean{MorganTianLib.existsUnique_leviCivita}
\uses{def:riemannian-metric}
\leanok



For every Riemannian metric $g$ there is a unique $\mathbb R$-linear
connection
\[
  \nabla:\Gamma(TM)\longrightarrow\Gamma(T^*M\otimes TM)
\]
satisfying $\nabla(fX)=df\otimes X+f\nabla X$ and, for all vector
fields $X,Y$,
\begin{itemize}
\item[$\bullet$] $d(g(X,Y))=g(\nabla X,Y)+g(X,\nabla Y)$
  (metric compatibility);
\item[$\bullet$] $\nabla_XY-\nabla_YX=[X,Y]$ (vanishing torsion),
  where $\nabla_XY=(\nabla Y)(X)$.
\end{itemize}
\end{theorem}

\begin{proof}
\sketch
\uses{def:riemannian-metric}
\leanok
Metric compatibility and vanishing torsion give the Koszul identity
\begin{eqnarray}\label{Koszul}
2g(\nabla_XY,Z) & = & X(g(Y,Z))+Y(g(X,Z))-Z(g(X,Y)) \\
& & +\, g([X,Y],Z)-g([X,Z],Y)-g([Y,Z],X). \nonumber
\end{eqnarray}
Nondegeneracy of $g$ makes this formula determine $\nabla_XY$, proving
uniqueness. For existence, define locally
\[
  \nabla_XY=X^i\bigl(\partial_iY^k+Y^j\Gamma^k_{ij}\bigr)\partial_k,
  \qquad
  \Gamma^k_{ij}=\tfrac12g^{kl}
  (\partial_i g_{lj}+\partial_j g_{il}-\partial_l g_{ij}).
\]
Symmetry in $i,j$ gives zero torsion, while lowering an index gives
\[
  \partial_k g_{ij}=g_{lj}\Gamma^l_{ki}+g_{il}\Gamma^l_{kj},
\]
which is metric compatibility. Uniqueness makes these local connections
agree on overlaps.
\end{proof}

This connection is the \emph{Levi--Civita connection}; $\nabla X$ is
the covariant derivative of $X$. In coordinates,
$\nabla_{\partial_i}\partial_j=\Gamma^k_{ij}\partial_k$, where
\begin{equation}\label{Gamma}
  \Gamma^{k}_{ij}=\frac12g^{kl}
  (\partial_i g_{lj}+\partial_j g_{il}-\partial_l g_{ij}).
\end{equation}
Thus $\Gamma^k_{ij}=\Gamma^k_{ji}$ and
$\partial_k g_{ij}=g_{lj}\Gamma^l_{ki}+g_{il}\Gamma^l_{kj}$.

The connection extends to tensor fields. For a smooth function $f$,
the convention is $\nabla f=df$, while $(\nabla f)^*$ denotes the
metric-dual gradient. Define
\begin{equation}\label{Hessian}
  \Hess(f)(X,Y)=X(Y(f))-(\nabla_XY)(f).
\end{equation}
\begin{lemma}[Hessian is symmetric]
\label{lem:hessian-symmetric}
\lean{MorganTianLib.hessian,MorganTianLib.hessian_symm,MorganTianLib.hessian_smul,MorganTianLib.hessianAt_symm,MorganTianLib.hessian_eq_metricInner_cov_gradientField,MorganTianLib.hessianAt_chartBasisVecFiber}
\uses{thm:levi-civita-connection}
\leanok
\label{Hessformula}



The Hessian is a symmetric covariant $2$-tensor. For smooth
$\phi,\psi$,
\[
  \Hess(f)(\phi X,\psi Y)=\phi\psi\Hess(f)(X,Y),
  \qquad \Hess(f)(X,Y)=\Hess(f)(Y,X).
\]
Equivalently,
\[
  \Hess(f)(X,Y)=\langle\nabla_X(\nabla f)^*,Y\rangle
  =\nabla^2f(X,Y),
  \qquad
  \Hess(f)_{ij}=\partial_i\partial_jf-(\partial_kf)\Gamma^k_{ij}.
\]
\end{lemma}

\begin{proof}
\sketch
\leanok
\uses{thm:levi-civita-connection}
Torsion-freeness gives
\[
  \Hess(f)(X,Y)-\Hess(f)(Y,X)
  =[X,Y](f)-(\nabla_XY-\nabla_YX)(f)=0.
\]
The Leibniz rule proves tensoriality. Metric compatibility yields
\begin{eqnarray}\label{hessform}
\langle\nabla_X(\nabla f)^*,Y\rangle
&=&X(Y(f))-(\nabla_XY)(f)\\
&=&\Hess(f)(X,Y).\nonumber
\end{eqnarray}
Finally, $df=(\partial_rf)dx^r$ and
$\nabla dx^k=-\Gamma^k_{ij}dx^i\otimes dx^j$, which gives the
coordinate formula.
\end{proof}

Write $\Hess_g(f)$ when the metric is not implicit. The Laplacian uses
the sign convention
\[
  \triangle f=g^{ij}\Hess(f)(\partial_i,\partial_j).
\]
For an orthonormal basis $\{X_i\}$ of $T_pM$,
\begin{equation}\label{laplacformula}
  \triangle f(p)=\sum_i\Hess(f)(X_i,X_i).
\end{equation}
It is nonnegative at a local minimum and therefore has nonpositive
spectrum.

\section{Curvature of a Riemannian manifold}

\begin{definition}[Riemann curvature tensor]
\label{def:riemann-curvature-tensor}
\lean{MorganTianLib.riemannCurvature,MorganTianLib.curvatureForm}
\uses{thm:levi-civita-connection}
\leanok



The \emph{Riemann curvature tensor} is
\[
  \mathcal R(X,Y)Z
  =\nabla_X\nabla_YZ-\nabla_Y\nabla_XZ-\nabla_{[X,Y]}Z
  =\nabla^2_{X,Y}Z-\nabla^2_{Y,X}Z,
\]
where $\nabla^2_{X,Y}Z=\nabla_X\nabla_YZ-\nabla_{\nabla_XY}Z$.
\end{definition}

In coordinates,
\[
  \mathcal R(\partial_i,\partial_j)\partial_k
  ={{R_{ij}}^l}_k\partial_l,
  \qquad
  {{R_{ij}}^l}_k
  =\partial_i\Gamma^l_{jk}-\partial_j\Gamma^l_{ik}
   +\Gamma^s_{jk}\Gamma^l_{is}-\Gamma^s_{ik}\Gamma^l_{js}.
\]
The associated $(0,4)$-tensor uses the convention
$\mathcal R(X,Y,Z,W)=g(\mathcal R(X,Y)W,Z)$; hence
\begin{align*}
  R_{ijkl}
  &=\mathcal R(\partial_i,\partial_j,\partial_k,\partial_l)\\
  &=g_{ks}{{R_{ij}}^s}_l\\
  &=g_{ks}(\partial_i\Gamma^s_{jl}-\partial_j\Gamma^s_{il}
    +\Gamma^t_{jl}\Gamma^s_{it}-\Gamma^t_{il}\Gamma^s_{jt}).
\end{align*}

\begin{lemma}[Curvature symmetries and Bianchi identities]
\label{claim:curvature-symmetries-bianchi}
\lean{MorganTianLib.curvatureForm_symmetries,MorganTianLib.curvatureForm_secondBianchi}
\uses{def:riemann-curvature-tensor}
\leanok
\label{Bianchi}



The curvature tensor satisfies
\[
  R_{ijkl}=-R_{jikl}=-R_{ijlk},
  \qquad R_{ijkl}=R_{klij}.
\]
Its cyclic sum over any three tensor indices is zero (the first Bianchi
identity). Its covariant derivative satisfies
\[
  R_{ijkl,h}+R_{ijlh,k}+R_{ijhk,l}=0,
  \qquad
  R_{ijkl,h}=(\nabla_{\partial_h}\mathcal R)_{ijkl}.
\]
\end{lemma}

\begin{proof}
\sketch
\uses{def:riemann-curvature-tensor,thm:levi-civita-connection}
\leanok
All identities are tensorial, so it suffices to use commuting coordinate fields. Antisymmetry in the first curvature pair is immediate from the definition. Metric compatibility gives zero pairing of the curvature operator with a repeated final vector; polarization gives antisymmetry in the last pair. Torsion-freeness yields the first Bianchi identity, and the resulting pair symmetries imply the cyclic identity for every choice of three slots. For the second Bianchi identity, write the curvature operators as commutators of covariant derivatives. The cyclic sum of their commutators vanishes by the Jacobi identity; the Christoffel correction terms cancel by symmetry of the connection coefficients and antisymmetry of curvature. Lowering indices and applying pair symmetry gives the displayed formula.
\end{proof}


\begin{definition}[Sectional curvature]
\label{def:sectional-curvature}
\lean{MorganTianLib.sectionalCurvatureAt}
\uses{def:riemann-curvature-tensor}
\leanok



For a $2$-plane $P\subset T_pM$ with orthonormal basis $\{X,Y\}$,
its \emph{sectional curvature} is
\[
  K(P)=\mathcal R(X,Y,X,Y).
\]
Positive, negative, nonnegative, and nonpositive sectional curvature
mean the corresponding inequality holds for every $2$-plane.
\end{definition}

For $X=X^i\partial_i$ and $Y=Y^j\partial_j$ spanning $P$,
\[
  K(P)=R_{ijkl}X^iY^jX^kY^l.
\]
The sectional curvature is \emph{constant} if this value is independent
of $p$ and $P\subset T_pM$.

\begin{lemma}[Curvature tensor of constant curvature manifolds]
\label{lem:constant-curvature-tensor}
\lean{MorganTianLib.hasConstantSectionalCurvature_iff_curvatureFormAt_eq}
\uses{def:sectional-curvature}
\leanok



A Riemannian manifold has constant sectional curvature $\lambda$ if
and only if
\[
  R_{ijkl}=\lambda(g_{ik}g_{jl}-g_{il}g_{jk}),
\]
equivalently,
\[
  \mathcal R(X,Y,Z,W)=\lambda\bigl(
  \langle X,Z\rangle\langle Y,W\rangle
  -\langle X,W\rangle\langle Y,Z\rangle\bigr).
\]
\end{lemma}

\begin{proof}
\sketch
\uses{def:sectional-curvature,claim:curvature-symmetries-bianchi}
\leanok
The tensor formula gives sectional curvature $\lambda$ on every
orthonormal pair. Conversely, subtract the model tensor from the
curvature tensor. The difference has all curvature symmetries and
vanishes on $T(X,Y,X,Y)$. Homogeneity and orthogonalization extend this
from orthonormal to arbitrary pairs. Polarizing twice and applying the
first Bianchi identity forces $T=0$.
\end{proof}

The standard sphere of radius $r$, Euclidean space, and hyperbolic space
have constant sectional curvatures $r^{-2}$, $0$, and $-1$,
respectively. The Poincar\'e ball and upper half-space metrics are
\[
  \frac{4\sum_i(dx^i)^2}{(1-|x|^2)^2},
  \qquad
  \frac{\sum_i(dx^i)^2}{(x^n)^2}.
\]
Their standard metrics are denoted $g_{\mathrm{st}}$; completeness and
simple connectivity are included in Lemma~\ref{lem:model-spaces}.


\begin{definition}[Curvature operator]
\label{def:curvature-operator}
\lean{MorganTianLib.curvatureOperator,MorganTianLib.curvatureOperator_iMulti,MorganTianLib.curvatureOperator_symm,MorganTianLib.sectionalCurvature_eq_curvatureOperator,MorganTianLib.HasPositiveCurvatureOperator,MorganTianLib.HasNonnegativeCurvatureOperator,MorganTianLib.sectionalCurvature_nonneg_of_hasNonnegativeCurvatureOperator}
\uses{def:riemann-curvature-tensor, def:sectional-curvature}
\leanok



The metric identifies the curvature tensor with the symmetric bilinear
form
\[
  \Rm(\varphi,\psi)=R_{ijkl}\varphi^{ij}\psi^{kl}
\]
on $\wedge^2TM$. The curvature operator is positive (respectively,
nonnegative) if $\Rm(\varphi,\varphi)>0$ for every nonzero $\varphi$
(respectively, $\Rm(\varphi,\varphi)\ge0$ for every $\varphi$).
\end{definition}

\begin{definition}[Norm of the curvature operator]
\label{def:curvature-operator-norm}
\lean{MorganTianLib.wedgeInner,MorganTianLib.HasCurvatureOperatorNormLe,MorganTianLib.abs_curvatureForm_le_of_hasCurvatureOperatorNormLe,MorganTianLib.exists_orthonormal_changeBasis,MorganTianLib.abs_sectionalCurvature_le_of_hasCurvatureOperatorNormLe}
\uses{def:curvature-operator, def:sectional-curvature}
\leanok



Give $\wedge^2T_xM$ the inner product
\[
  \langle X\wedge Y,Z\wedge W\rangle
  =g(X,Z)g(Y,W)-g(X,W)g(Y,Z).
\]
Thus $\{e_i\wedge e_j\}_{i<j}$ is orthonormal when $\{e_i\}$ is, and
\[
  \Rm(X\wedge Y,X\wedge Y)=K(P)
\]
for an orthonormal basis $X,Y$ of $P$. For $K\ge0$, the notation
$|\Rm(x)|\le K$ means every eigenvalue of the associated self-adjoint
operator on $\wedge^2T_xM$ lies in $[-K,K]$; equivalently,
\[
  |\Rm(\varphi,\varphi)|\le K\langle\varphi,\varphi\rangle
  \quad(\varphi\in\wedge^2T_xM).
\]
In particular, $|K(P)|\le K$ for every $2$-plane $P\subset T_xM$.
\end{definition}

Positive (respectively, nonnegative) curvature operator implies positive
(respectively, nonnegative) sectional curvature.


\begin{definition}[Ricci curvature and scalar curvature]
\label{def:ricci-curvature}
\lean{MorganTianLib.ricciAt,MorganTianLib.scalarCurvatureAt}
\uses{def:riemann-curvature-tensor}
\leanok



The symmetric \emph{Ricci tensor} is
\[
  \Ric(X,Y)=g^{kl}R(X,\partial_k,Y,\partial_l)
  =\sum_iR(X,e_i,Y,e_i)
\]
in an orthonormal basis $\{e_i\}$, and
$\Ric=\Ric_{ij}dx^i\otimes dx^j$. The \emph{scalar curvature} is
\[
  R=R_g=\tr_g\Ric=g^{ij}\Ric_{ij}.
\]
The inequalities $\Ric\ge k$ and $\Ric\le k$ refer to the eigenvalues
of the Ricci endomorphism.
\end{definition}

\begin{lemma}[Ricci curvature as a trace of the curvature operator]
\label{lem:ricci-trace-curvature-operator}
\lean{MorganTianLib.metricInner_riemannCurvature_eq_curvatureFormAt,MorganTianLib.sum_metricInner_riemannCurvature_self_eq_ricciAt}
\uses{def:ricci-curvature, def:riemann-curvature-tensor, claim:curvature-symmetries-bianchi}



Let $X$ be a smooth vector field, $p\in M$, and
$\{e_1,\ldots,e_n\}$ an orthonormal basis of $T_pM$, extended to
smooth fields $E_i$. Then
\[
  \sum_{i=1}^n\langle\mathcal R(E_i,X)X,E_i\rangle(p)
  =\Ric(X(p),X(p)).
\]
\end{lemma}

\begin{proof}
\sketch
\uses{def:ricci-curvature,claim:curvature-symmetries-bianchi}
\leanok
Tensoriality reduces each summand to
$R(e_i,X(p),e_i,X(p))$. Applying antisymmetry in both index pairs
gives $R(X(p),e_i,X(p),e_i)$, whose sum is the defining orthonormal
trace of $\Ric(X(p),X(p))$.
\end{proof}

For a diffeomorphism $F:N\to M$, curvature is natural under pullback:
\[
  \Rm(F^*g)=F^*\Rm(g),\qquad
  \Ric(F^*g)=F^*\Ric(g),\qquad
  R(F^*g)=F^*R(g).
\]

\subsection{Consequences of the Bianchi identities}

For a covariant $2$-tensor $\omega$, define
\[
  \operatorname{div}(\omega)(X)=\nabla^*\omega(X)
  =g^{rs}(\nabla_r\omega)(X,\partial_s).
\]

\begin{lemma}[Divergence of Ricci from second Bianchi identity]
\label{lem:div-ricci-scalar-curvature}
\lean{MorganTianLib.dir_scalarCurvature_eq_two_divRicci,MorganTianLib.divRicciAt,MorganTianLib.covRicciAt,MorganTianLib.divRicciAt_eq_sum_orthonormalBasis,MorganTianLib.covRicci_eq_frame_sum,MorganTianLib.dir_scalarCurvatureAt_eq_two_frame_covRicci,MorganTianLib.covRicci_congr_apply}
\uses{claim:curvature-symmetries-bianchi, def:ricci-curvature}
\leanok
\label{divRic}



$$dR=2\operatorname{div}(\Ric)=2\nabla^*\Ric.$$
\end{lemma}

\begin{proof}
\sketch
\leanok
\uses{claim:curvature-symmetries-bianchi,def:ricci-curvature}
Contract the second Bianchi identity in its first and third, then second and fourth, tensor indices. Covariant differentiation commutes with contraction because the metric is parallel. Curvature antisymmetry turns each of the last two contracted terms into minus the divergence of the Ricci tensor, while the first becomes the differential of the scalar curvature. Hence the contraction is the stated identity.
\end{proof}


For a tensor field $\omega$, set
\[
  \nabla^2\omega(\cdot,X,Y)
  =(\nabla_X\nabla_Y\omega)(\cdot)-\nabla_{\nabla_XY}\omega(\cdot),
  \qquad
  \triangle\omega=g^{ij}\nabla^2\omega(\cdot,\partial_i,\partial_j).
\]
The derivative is generally not symmetric in $X,Y$; its commutator is
the curvature action. The vector-field case is recorded first, and the
one-form case follows by duality.

\begin{definition}[Second covariant derivative of a vector field]
\label{def:second-covariant-derivative-vector-field}
\lean{MorganTianLib.secondCov}
\uses{thm:levi-civita-connection}
\leanok



For an affine connection $\nabla$ and vector fields $X,Y,Z$, the
\emph{second covariant derivative} of $Z$ is
\[
  \nabla^2Z(X,Y)=\nabla_X\nabla_YZ-\nabla_{\nabla_XY}Z.
\]
\end{definition}

\begin{lemma}[Tensoriality of the second covariant derivative]
\label{lem:second-covariant-derivative-tensorial}
\lean{MorganTianLib.secondCov_add_left,MorganTianLib.secondCov_smul_left,MorganTianLib.secondCov_add_middle,MorganTianLib.secondCov_smul_middle,MorganTianLib.secondCov_congr_left}
\uses{def:second-covariant-derivative-vector-field}



For fixed $Z$, the map $(X,Y)\mapsto\nabla^2Z(X,Y)$ is additive in
each argument and satisfies
\[
  \nabla^2Z(\varphi X,Y)=\varphi\nabla^2Z(X,Y)
  =\nabla^2Z(X,\varphi Y)
\]
for every smooth $\varphi$. Hence its value at $p$ depends on $X,Y$
only through $X(p),Y(p)$.
\end{lemma}

\begin{proof}
\sketch
\uses{thm:levi-civita-connection}
\leanok
Additivity follows from additivity of the connection. Scaling the first direction is immediate from linearity in the direction slot. Scaling the second direction produces the same Leibniz cross term in both terms defining the second covariant derivative, so those terms cancel. The resulting smooth-function bilinearity proves pointwise dependence.
\end{proof}

\begin{lemma}[Ricci commutation identity for vector fields]
\label{lem:ricci-commutation-vector-field}
\lean{MorganTianLib.secondCov_sub_swap}
\uses{def:second-covariant-derivative-vector-field, def:riemann-curvature-tensor}
\leanok



Let $\nabla$ be a torsion-free affine connection on $M$. For all
smooth vector fields $X$, $Y$, $Z$,
$$\nabla^2Z(X,Y)-\nabla^2Z(Y,X)={\mathcal R}(X,Y)Z,$$
where ${\mathcal R}(X,Y)Z
=\nabla_X\nabla_YZ-\nabla_Y\nabla_XZ-\nabla_{[X,Y]}Z$ is the
curvature operator of Definition~\ref{def:riemann-curvature-tensor}.
\end{lemma}

\begin{proof}
\sketch
\uses{thm:levi-civita-connection,def:riemann-curvature-tensor}
\leanok
Subtracting the defining expressions,
$$\nabla^2Z(X,Y)-\nabla^2Z(Y,X)
=\nabla_X\nabla_YZ-\nabla_Y\nabla_XZ
-\left(\nabla_{\nabla_X(Y)}Z-\nabla_{\nabla_Y(X)}Z\right).$$
By additivity of the direction slot of $\nabla$, the parenthesized
difference is $\nabla_{\nabla_X(Y)-\nabla_Y(X)}Z$, and
torsion-freeness gives $\nabla_X(Y)-\nabla_Y(X)=[X,Y]$; so the
right-hand side is
$\nabla_X\nabla_YZ-\nabla_Y\nabla_XZ-\nabla_{[X,Y]}Z
={\mathcal R}(X,Y)Z$.
\end{proof}

\begin{lemma}[Laplacian of differential form]
\label{lem:laplacian-one-form}
\lean{MorganTianLib.oneFormLaplacianAt,MorganTianLib.oneFormLaplacianAt_eq_sum_frame,MorganTianLib.metricInner_secondCov_gradientField_symm,MorganTianLib.sum_metricInner_secondCov_along_frame_eq_dir_laplacianAt,MorganTianLib.sum_metricInner_riemannCurvature_frame_eq_ricciAt,MorganTianLib.oneFormLaplacianAt_gradientField_eq_dirTangent_laplacianAt_add_ricciAt,MorganTianLib.oneFormLaplacianAt_gradientField_eq_metricInner_gradientAt_add_ricciAt}
\uses{lem:hessian-symmetric, def:ricci-curvature}
\label{lapformula}



For every smooth function $f$,
\[
  \triangle(df)=d(\triangle f)+\Ric((\nabla f)^*,\mathord\cdot).
\]
\end{lemma}

\begin{proof}
\sketch
\uses{lem:hessian-symmetric,def:ricci-curvature,claim:curvature-symmetries-bianchi,def:second-covariant-derivative-vector-field,lem:ricci-commutation-vector-field}
\leanok
Write $H=\nabla(df)=\Hess(f)$, a symmetric two-tensor
(Lemma~\ref{lem:hessian-symmetric}), and write
$\nabla^2_{X,Y}=\nabla_X\nabla_Y-\nabla_{\nabla_XY}$ for the second
covariant derivative in the directions $X,Y$, as in the discussion
preceding the lemma. A direct expansion shows that for any one-form
$\omega$ and vector field $Z$,
$$\left(\nabla^2_{X,Y}\omega\right)(Z)=(\nabla_X(\nabla\omega))(Y,Z),$$
where $\nabla\omega$ is regarded as a two-tensor with
$(\nabla\omega)(Y,Z)=(\nabla_Y\omega)(Z)$; in particular
$\left(\nabla^2_{X,Y}df\right)(Z)=(\nabla_XH)(Y,Z)$. In these terms
the two sides of the asserted formula are, at a point, with
$\{\partial_i\}$ local coordinate fields,
$$\triangle(df)(Z)=g^{jk}\left(\nabla_{\partial_j}H\right)(\partial_k,Z),
\qquad
d(\triangle f)(Z)=Z\left(g^{jk}H_{jk}\right)
=g^{jk}\left(\nabla_ZH\right)(\partial_j,\partial_k),$$
the latter because $\nabla g=0$, so differentiating the trace
commutes with the contraction.

\emph{Commutation formula for one-forms.} We claim that for any
one-form $\omega$,
\begin{equation}\label{oneformcomm}
\left(\nabla^2_{X,Y}\omega-\nabla^2_{Y,X}\omega\right)(Z)
=-\omega\left({\mathcal R}(X,Y)Z\right).
\end{equation}
Indeed, apply the second covariant derivative to the function
$\omega(Z)$: the Leibniz rule gives
$$\nabla^2_{X,Y}\left(\omega(Z)\right)
=\left(\nabla^2_{X,Y}\omega\right)(Z)
+(\nabla_X\omega)(\nabla_YZ)+(\nabla_Y\omega)(\nabla_XZ)
+\omega\left(\nabla^2_{X,Y}Z\right).$$
Antisymmetrize in $X$ and $Y$. The left-hand side is the Hessian of
the function $\omega(Z)$, which is symmetric in $X,Y$ by
Lemma~\ref{lem:hessian-symmetric}; the two middle terms on the right
are symmetric as a pair; and
$\nabla^2_{X,Y}Z-\nabla^2_{Y,X}Z={\mathcal R}(X,Y)Z$ by the
definition of the curvature tensor. This yields
Equation~(\ref{oneformcomm}).

Applied to $\omega=df$, Equation~(\ref{oneformcomm}) says
\begin{equation}\label{Hcomm}
(\nabla_XH)(Y,Z)-(\nabla_YH)(X,Z)=-df\left({\mathcal
R}(X,Y)Z\right).
\end{equation}

\emph{Comparing the two sides.} Fix a coordinate direction
$\partial_i$. Using Equation~(\ref{Hcomm}) twice and the symmetry
of $H$ (hence of $\nabla_WH$) in its two arguments,
\begin{eqnarray*}
(\nabla_{\partial_j}H)(\partial_k,\partial_i)
-(\nabla_{\partial_i}H)(\partial_k,\partial_j)
& = & \left[(\nabla_{\partial_j}H)(\partial_k,\partial_i)
-(\nabla_{\partial_k}H)(\partial_j,\partial_i)\right]
+\left[(\nabla_{\partial_k}H)(\partial_i,\partial_j)
-(\nabla_{\partial_i}H)(\partial_k,\partial_j)\right]\\
& = & -df\left({\mathcal R}(\partial_j,\partial_k)\partial_i\right)
-df\left({\mathcal R}(\partial_k,\partial_i)\partial_j\right).
\end{eqnarray*}
Contract with $g^{jk}$. The first term on the right contracts to
zero, since $g^{jk}$ is symmetric in $j,k$ while ${\mathcal
R}(\partial_j,\partial_k)$ is antisymmetric. For the second term,
writing $df=f_{,s}dx^s$ and ${\mathcal
R}(\partial_k,\partial_i)\partial_j={{R_{ki}}^s}_j\partial_s$ with
${{R_{ki}}^s}_j=g^{sm}R_{kimj}$, the symmetries of the curvature
tensor give $R_{kimj}=-R_{ikmj}$, so that
$$-g^{jk}\,df\left({\mathcal R}(\partial_k,\partial_i)\partial_j\right)
=-g^{jk}g^{sm}f_{,s}R_{kimj}
=g^{sm}f_{,s}\,g^{jk}R_{ikmj}
=g^{sm}f_{,s}\Ric_{im}
=\Ric\left((\nabla f)^*,\partial_i\right),$$
using the definition $\Ric_{im}=g^{jk}R_{ikmj}$ (contraction on the
second and fourth indices) and the fact that $(\nabla
f)^*=g^{sm}f_{,s}\partial_m$ is the gradient of $f$. Hence
$$g^{jk}(\nabla_{\partial_j}H)(\partial_k,\partial_i)
-g^{jk}(\nabla_{\partial_i}H)(\partial_j,\partial_k)
=\Ric\left((\nabla f)^*,\partial_i\right),$$
which, by the identifications in the first paragraph, is exactly
$\triangle df=d\triangle f+\Ric((\nabla f)^*,\cdot)$ evaluated on
$\partial_i$.
\end{proof}




\subsection{First examples}

\begin{theorem}[Uniformization Theorem]
\label{thm:uniformization}
\uses{def:sectional-curvature}

Let $(M^n,g)$, $n\ge2$, be complete, simply connected, and of constant
sectional curvature $\lambda$. Then:
\begin{enumerate}
\item If $\lambda=0$, then $(M,g)$ is isometric to Euclidean
  $n$-space.
\item If $\lambda>0$, some diffeomorphism $\phi:M\to S^n$ satisfies
  $g=\lambda^{-1}\phi^*g_{\mathrm{st}}$.
\item If $\lambda<0$, some diffeomorphism $\phi:M\to\HH^n$ satisfies
  $g=|\lambda|^{-1}\phi^*g_{\mathrm{st}}$, where
  $g_{\mathrm{st}}$ is the curvature $-1$ Poincar\'e metric.
\end{enumerate}
\end{theorem}

\begin{proof}
\sketch
\uses{def:sectional-curvature,thm:hopf-rinow,def:exponential-map,lem:constant-curvature-jacobi,lem:local-isometry-covering,lem:local-isometry-rigidity,lem:local-isometry-exponential,def:local-isometry,thm:levi-civita-connection,lem:model-spaces,lem:model-polar-isometry,lem:submanifold-connection}
The restriction $n\ge2$ prevents the sectional-curvature condition from
being vacuous. Rescaling $g$ by $c>0$ preserves the connection and
changes sectional curvature from $K$ to $K/c$. Constant-curvature
Jacobi fields give the polar metric
\begin{equation}\label{modelpolar}
  G=dr^2+s_\lambda(r)^2g_{S^{n-1}}.
\end{equation}
For $\lambda\le0$, the exponential map is a global isometry from this
polar model; it is Euclidean when $\lambda=0$ and agrees, after scaling,
with the hyperbolic model when $\lambda<0$. For $\lambda>0$, the same
formula identifies punctured normal balls with the sphere of radius
$\lambda^{-1/2}$. Local isometries based at two non-antipodal points
agree on their connected overlap, hence glue across the missing
antipodes. The resulting local isometry from the complete sphere is a
covering, and simple connectivity makes it one-sheeted.
\end{proof}

A complete constant-curvature manifold is a quotient of the
corresponding simply connected model by a free discrete group of
isometries; such a quotient is a \emph{space form}.


\begin{definition}[Einstein manifold]
\label{def:einstein-manifold}
\lean{MorganTianLib.IsEinstein}
\uses{def:ricci-curvature}
\leanok



A Riemannian manifold is \emph{Einstein with constant $\lambda$} if
$\Ric_g=\lambda g$.
\end{definition}

\begin{example}[Einstein manifolds in dimensions 2 and 3]
\label{ex:einstein-dimension-2-3}
\lean{MorganTianLib.IsEinstein.sectionalCurvature_const}
\uses{def:einstein-manifold, def:sectional-curvature}
\leanok



If $n\in\{2,3\}$ and $(M^n,g)$ is Einstein with constant $\lambda$,
then it has constant sectional curvature $\lambda/(n-1)$.
\end{example}

\begin{proof}
\sketch
\uses{def:einstein-manifold,def:sectional-curvature,def:ricci-curvature,claim:curvature-symmetries-bianchi}
\leanok
For an orthonormal basis, write $K_{ij}$ for the sectional curvature
of $\operatorname{span}(e_i,e_j)$. Then
$\Ric(e_i,e_i)=\sum_{j\ne i}K_{ij}$. In dimension $2$ this gives
$K_{12}=\lambda$. In dimension $3$, the three Einstein equations
$K_{12}+K_{13}=K_{12}+K_{23}=K_{13}+K_{23}=\lambda$ force every
$K_{ij}$ to equal $\lambda/2$.
\end{proof}

\begin{remark}[Complete Einstein surfaces and $3$-manifolds are space forms]
\label{rem:einstein-space-form}
\uses{ex:einstein-dimension-2-3, thm:uniformization}

If the manifold in Example~\ref{ex:einstein-dimension-2-3} is
complete, Theorem~\ref{thm:uniformization} makes it a space form of
curvature $\lambda/(n-1)$.
\end{remark}

\subsection{Cones}

\begin{definition}[Open cone]
\label{def:cone}
\lean{MorganTianLib.positiveReal,MorganTianLib.coneMetric, MorganTianLib.coneMetric_metricInner_mk}
\uses{def:riemannian-metric}
\leanok
\label{conedefn}
The \emph{open cone} over $(N,g)$ is
\[
  c(N)=N\times(0,\infty),
  \qquad \tilde g=s^2g+ds^2.
\]
\end{definition}

Fix local coordinates $(x^1,\ldots,x^{n-1})$ on $N$ (of dimension
$n-1$, as below). Let
$\Gamma_{ij}^k;\ 1\le i,j,k\le n-1$, be the Christoffel symbols for
the Levi-Civita connection on $N$. Set $x^0=s$. In the local
coordinates $(x^0,x^1,\ldots,x^{n-1})$ for the cone we have the
Christoffel symbols $\widetilde \Gamma_{ij}^k;\ 0\le i,j,k\le n-1$,
for $\tilde g$. The relation between the metrics gives the following
relations between the two sets of Christoffel symbols:
\begin{eqnarray*}\widetilde
\Gamma_{ij}^k & = & \Gamma_{ij}^k; \ \ \ \ 1\le i,j,k\le n-1 \\
\widetilde \Gamma_{ij}^0 & = & -sg_{ij}; \ \ \ \ 1\le i,j\le n-1 \\
\widetilde\Gamma_{i0}^j & = & \widetilde\Gamma_{0i}^j =
s^{-1}\delta^j_i; \ \ \ \ 1\le i,j\le n-1 \\
\widetilde \Gamma_{i0}^0 & = & 0;\ \ \ \ 0\le i\le n-1 \\
\widetilde \Gamma_{00}^i & = & 0;\ \  \ \ 0\le i\le n-1.
\end{eqnarray*}

Denote by ${\mathcal R}_g$ the curvature tensor for $g$ and by
${\mathcal R}_{\tilde g}$ the curvature tensor for $\tilde g$. Then
the above formulas lead directly to:
\begin{eqnarray*}
{\mathcal R}_{\tilde g}(\partial_i,\partial_j)(\partial_0) & = & 0;\ \ \ \  0\le i,j\le n-1 \\
{\mathcal R}_{\tilde g}(\partial_i,\partial_j)(\partial_i) & = &
{\mathcal
R}_g(\partial_i,\partial_j)(\partial_i)+g_{ii}\partial_j-g_{ji}\partial_i;\
\ \ \  1\le i,j\le n-1 \\
\end{eqnarray*}

This allows us to compute the Riemann curvatures of the cone in
terms of those of $N$.

\begin{proposition}[Curvature operator of cone]
\label{prop:cone-curvature}
\lean{MorganTianLib.coneWedgeEquiv,MorganTianLib.coneWedgeInnerAt_horizontal_mixed,MorganTianLib.coneCurvatureOperatorAt,MorganTianLib.coneBaseCurvatureDifferenceAt,MorganTianLib.coneCurvatureBlockFormAt,MorganTianLib.coneCurvatureOperatorAt_eq_block}
\uses{def:cone, def:curvature-operator}
\leanok
\label{conecurv}

Let $N$ be a Riemannian manifold of dimension $n-1$. Fix $(p,s)\in
c(N)=N\times (0,\infty)$. With respect to the orthogonal splitting
$$\wedge^2T_{(p,s)}c(N)=\wedge^2T_pN\oplus
\left(T_pN\wedge\partial_0\right),$$
where $\partial_0$ is the coordinate field of $x^0=s$, the curvature
operator $\Rm_{\tilde g}(p,s)$ of the cone is block-diagonal,
$$\begin{pmatrix} s^2(\Rm_g(p)-\wedge^2g(p)) & 0 \\ 0 & 0
\end{pmatrix},$$
where $\wedge^2g(p)$ is the symmetric form on $\wedge^2T_pN$ induced
by $g$. (The corresponding display in the original source places
the non-zero block in an off-diagonal corner, where it is not
symmetric; the block belongs on the diagonal, as the proof shows.)
\end{proposition}

\begin{proof}
\sketch
\uses{def:cone,def:curvature-operator,thm:levi-civita-connection,def:riemann-curvature-tensor,claim:curvature-symmetries-bianchi}
Insert the cone Christoffel symbols into the coordinate curvature formula. Every component containing the radial index vanishes, while for tangential indices
\[
  \widetilde R_{ijkl}=s^2\bigl(R^g_{ijkl}-(g_{ik}g_{jl}-g_{il}g_{jk})\bigr).
\]
Under $\wedge^2T_{(p,s)}c(N)=\wedge^2T_pN\oplus(T_pN\wedge\partial_s)$, the second summand is therefore null and the first carries $s^2(\Rm_g-\wedge^2g)$, giving the claimed block decomposition.
\end{proof}

\begin{corollary}[Eigenvalues of cone curvature operator]
\label{cor:cone-eigenvalues}
\uses{prop:cone-curvature}

For any $p\in N$ let $\lambda_1,\ldots,\lambda_{(n-1)(n-2)/2}$ be
the eigenvalues of $\Rm_g(p)$. Then for any $s>0$  there are
$(n-1)$ zero eigenvalues of $\Rm_{\tilde g}(p,s)$. The other
$(n-1)(n-2)/2$ eigenvalues of $\Rm_{\tilde g}(p,s)$ are
$s^{-2}(\lambda_i-1)$.
\end{corollary}


\begin{proof}
\sketch
\uses{prop:cone-curvature}
Clearly from Proposition~\ref{prop:cone-curvature}, we see that under the
orthogonal decomposition $\wedge^2T_{(p,s)}c(N)=\wedge^2T_pN\oplus
T_pN$ the second subspace is contained in the null space of $ \Rm_{\tilde g}(p,s)$, and hence contributes $(n-1)$ zero
eigenvalues. Likewise, from this proposition we see that the
eigenvalues of the restriction of $\Rm_{\tilde g}(p,s)$ to the
subspace $\wedge^2T_pN$ are  given by
$s^{-4}(s^2(\lambda_i-1))=s^{-2}(\lambda_i-1)$.
\end{proof}


\section{Geodesics and the exponential map}

Here we review standard material about geodesics, Jacobi fields, and
the exponential map.

\subsection{Geodesics and the energy functional}

\begin{definition}[Geodesic]
\label{def:geodesic}
\lean{MorganTianLib.IsGeodesicCurveOn,MorganTianLib.IsGeodesicOn}
\uses{thm:levi-civita-connection}



Let $I$ be an open interval. A smooth curve $\gamma\colon I\to M$ whose
covariant acceleration vanishes, $\nabla_{\dot\gamma}\dot\gamma=0$, is a
geodesic. On a fixed parameter interval one may instead use the
continuous-on-the-interval equation predicate; when the usual chart ODE
regularity hypotheses are supplied, its solutions are smooth geodesics.
\end{definition}

In local coordinates, we write $\gamma(t)=(x^1(t),\ldots,x^n(t))$
and this equation becomes
$$0=\nabla_{\dot \gamma}\dot \gamma(t) =
\left(\sum\limits_k\left(\ddot x^k(t) +\dot x^i(t)\dot
x^j(t)\Gamma^{k}_{ij}(\gamma(t))\right)\partial_k\right).$$ This is
a system of $2^{nd}$ order ODE's. The local existence, uniqueness
and smoothness of a geodesic through any point $p\in M$ with initial
velocity vector $v\in T_pM$ follow from the classical ODE theory.
We call geodesics that minimize the length between their endpoints
{\em minimizing geodesics}. On a complete manifold, any two points
are joined by a minimizing geodesic, and all geodesics are defined
for all time; this is the Hopf--Rinow theorem
(Theorem~\ref{thm:hopf-rinow}), which we state and prove in the
subsection on the exponential mapping below, once the exponential
map is available.


Geodesics are critical points of the energy functional. Let $(M,g)$
be a complete Riemannian manifold. Consider the space of $C^1$-paths
in $M$ parameterized by the unit interval. On this space we have the
energy functional
$$E(\gamma)=\frac{1}{2}\int_{0}^{1}\langle \gamma'(t),\gamma'(t)\rangle dt.$$
 Suppose that we have a one-parameter family of  paths parameterized by
 $[0,1]$,
 all having the same initial point $p$ and the same final point $q$.
 By this we mean that we have a surface $\tilde\gamma(t,u)$ with the property that for each
 $u$ the path $\gamma_u=\tilde\gamma(\cdot,u)$ is a path from $p$ to $q$ parameterized by $[0,1]$.
 Let $\widetilde X=\partial\tilde \gamma/\partial t$ and
 $\widetilde Y=\partial\tilde \gamma/\partial u$ be the corresponding vector
 fields along the surface swept out by $\tilde \gamma$, and denote by $X$ and $Y$ the restriction
 of these vector fields along $\gamma_0$.
 We compute
 \begin{eqnarray*}
 \frac{dE(\gamma_u)}{du}\Bigl|_{u=0}\Bigr.& = & \left( \int_0^1\langle \nabla_{\widetilde Y}\widetilde X,\widetilde
 X\rangle dt\right)\left|_{u=0}\right. \\
 & = & \left(\int_0^1\langle\nabla_{\widetilde X}\widetilde Y,\widetilde
 X\rangle dt \right)\left|_{u=0}\right.\\
 & = & -\left( \int_0^1\langle\nabla_{\widetilde X}\widetilde X,\widetilde
 Y\rangle dt\right)\left|_{u=0}\right.=-\int_0^1\langle\nabla_XX,Y\rangle,
\end{eqnarray*}
where the first equality in the last line comes from integration by
parts and the fact that $\widetilde Y$ vanishes at the endpoints.
Given any vector field $Y$ along $\gamma_0$ there is a one-parameter
family $\tilde \gamma(t,u)$ of paths from $p$ to $q$ with $\tilde
\gamma(t,0)=\gamma_0$ and with $\widetilde Y(t,0)=Y$. Thus, from the
above expression we see that $\gamma_0$ is a critical point for the
energy functional on the space of paths from $p$ to $q$
parameterized by the interval $[0,1]$ if and only if $\gamma_0$ is a
geodesic.

Notice that it follows immediately from the geodesic equation that
the length of a tangent vector along a geodesic is constant. Thus,
if a geodesic is parameterized by $[0,1]$ we have
$$E(\gamma)=\frac{1}{2}L(\gamma)^2.$$
It is immediate from the Cauchy-Schwarz inequality that for any
curve $\mu$ parameterized by $[0,1]$ we have
$$E(\mu)\ge \frac{1}{2}L(\mu)^2$$
with equality if and only if $|\mu'|$ is constant. In particular, a
curve parameterized by $[0,1]$ minimizes distance between its
endpoints if it is a minimum for the energy functional on all paths
parameterized by $[0,1]$ with the given endpoints.






\subsection{Families of geodesics and Jacobi fields}


Consider a family of geodesics $\tilde\gamma(u,t)=\gamma_u(t)$
parameterized by the interval  $[0,1]$ with $\gamma_u(0)=p$ for all
$u$. Here, unlike the discussion above, we allow $\gamma_u(1)$ to
vary with $u$. As before define vector fields along the surface
swept out by $\tilde \gamma$:  $\widetilde X=\partial
\tilde\gamma/\partial t$ and let $\tilde
Y=\partial\tilde\gamma/\partial u$.  We denote by $X$ and $Y$ the
restriction of these vector fields to the geodesic
$\gamma_0=\gamma$.
 Since each
$\gamma_u$ is a geodesic, we have $\nabla_{\widetilde X}{\widetilde
X}=0$. Differentiating this equation in the $\widetilde Y$-direction
yields $\nabla_{\widetilde Y}\nabla_{\widetilde X}\widetilde X=0$.
Interchanging the order of differentiation, using
$\nabla_{\widetilde X}\widetilde Y=\nabla_{\widetilde Y}{\widetilde
X}$, and then restricting to $\gamma$, we get the {\em Jacobi
equation}:
$$\nabla_X\nabla_XY+{\mathcal R}(Y,X)X=0.$$
Notice that the left-hand side of the equation depends only on the
value of $Y$ along $\gamma$, not on the entire family. We denote the
left-hand side of this equation by $\Jac(Y)$, so that the Jacobi
equation now reads
$$\Jac(Y)=0.$$
 The fact that all the geodesics begin at the same point at time
$0$ means that $Y(0)=0$. A vector field $Y$ along a geodesic
$\gamma$ is said to be a {\em Jacobi field} if
it satisfies this equation and vanishes at the initial point $p$. A
Jacobi field is determined by its first derivative at $p$, i.e., by
$\nabla_XY(0)$. We have just seen that this is the equation
describing, to first order, variations of $\gamma$ by a family of
geodesics with the same starting point.

Jacobi fields are also determined by the energy functional. Consider
the space of paths parameterized by $[0,1]$ starting at a given
point $p$ but free to end anywhere in the manifold. Let $\gamma$ be
a geodesic (parameterized by $[0,1]$) from $p$ to $q$. Associated to
any one-parameter family $\tilde \gamma(t,u)$ of paths parameterized
by $[0,1]$ starting at $p$ we associate the second derivative of the
energy at $u=0$. Straightforward computation gives
$$\frac{d^2E(\gamma_u)}{du^2}\Bigl|_{u=0}\Bigr.=\langle \nabla_XY(1),Y(1)\rangle+\langle
X(1),\nabla_Y\widetilde Y(1,0)\rangle-\int_0^1\langle \Jac(Y),Y\rangle dt.$$

Notice that the first term is a boundary term from the integration
by parts, and it depends not just on the value of $Y$ (i.e., on
$\widetilde Y$ restricted to $\gamma$) but also on the first-order
variation of $\widetilde Y$ in the $Y$ direction. There is the
associated bilinear form that comes from two-parameter families
$\tilde \gamma(t,u_1,u_2)$ whose value at $u_1=u_1=0$ is $\gamma$.
It is
$$\frac{d^2E}{du_1du_2}\Bigl|_{u_1=u_2=0}\Bigr.=\langle \nabla_XY_1(1),Y_2(1)\rangle+\langle
X(1),\nabla_{Y_1}\widetilde Y_2(1,0)\rangle-\int_0^1\langle \Jac(Y_1),Y_2\rangle dt.$$ Notice that restricting to the space of
vector fields that vanish at both endpoints, the second derivatives
depend only on $Y_1$ and $Y_2$ and the formula is
$$\frac{d^2E}{du_1du_2}\Bigl|_{u_1=u_2=0}\Bigr.=-\int_0^1\langle \Jac(Y_1),Y_2\rangle dt,$$
so that this expression is symmetric in $Y_1$ and $Y_2$. The
associated quadratic form on the space of vector fields along
$\gamma$ vanishing at both endpoints
$$-\int_0^1\langle \Jac(Y),Y\rangle dt$$
is the second derivative of the energy function at $\gamma$ for any
one-parameter family whose value at $0$ is $\gamma$ and whose first
variation is given by $Y$.

The interchange of covariant derivatives used in the derivation of
the Jacobi equation above, and the Jacobi equation itself for a
family of geodesics, are recorded precisely — in local coordinates —
in the following lemma.

\begin{lemma}[Covariant commutation and the Jacobi equation for families]
\label{lem:covariant-commutation-jacobi}
\lean{MorganTianLib.covDerivAlong,MorganTianLib.christoffelCurvature,MorganTianLib.covDerivAlong_comm,MorganTianLib.covDerivAlong_fderiv_symm,MorganTianLib.covDerivAlong_geodesic_family_jacobi,MorganTianLib.chartChristoffelBilin,MorganTianLib.chartCurvature,MorganTianLib.chart_geodesic_family_jacobi}
\uses{thm:levi-civita-connection, def:riemann-curvature-tensor, def:geodesic}
\leanok



Let $U$ be a coordinate chart of $M$, with Christoffel symbols
$\Gamma^k_{ij}$ of the Levi-Civita connection
(Theorem~\ref{thm:levi-civita-connection}). For $x\in U$ write
$\Gamma_x(X,Y)$ for the symmetric bilinear map with components
$\Gamma^k_{ij}(x)X^iY^j$. Let $u=u(s,t)$ be a $C^2$ two-parameter
family of points of $U$, written in coordinates, and for a $C^2$
vector field $V(s,t)$ along $u$ (in coordinates) set
$$\nabla_sV=\partial_sV+\Gamma_u(\partial_su,V),\qquad
\nabla_tV=\partial_tV+\Gamma_u(\partial_tu,V).$$
Then:
\begin{enumerate}
\item ({\em Curvature commutation})
$\nabla_s\nabla_tV-\nabla_t\nabla_sV
={\mathcal R}(\partial_su,\partial_tu)V$, where ${\mathcal R}$ is
the curvature tensor of
Definition~\ref{def:riemann-curvature-tensor} in its coordinate
expression.
\item ({\em Torsion-freeness})
$\nabla_s\partial_tu=\nabla_t\partial_su$.
\item ({\em Jacobi equation}) If $u$ is $C^3$ and every line
$t\mapsto u(s,t)$ is a geodesic, then $Y=\partial_su$ satisfies
$$\nabla_t\nabla_tY+{\mathcal R}(Y,\partial_tu)\partial_tu=0.$$
\end{enumerate}
\end{lemma}

\begin{proof}
\sketch
\leanok
\uses{thm:levi-civita-connection,def:riemann-curvature-tensor,def:geodesic}
Expand both iterated covariant derivatives in coordinates. Mixed partial derivatives and the symmetric Christoffel terms cancel, leaving exactly the coordinate curvature expression. Symmetry of the Christoffel symbols also gives $\nabla_s\partial_tu=\nabla_t\partial_su$. Differentiating the geodesic equation in the family parameter and applying these two identities yields the Jacobi equation for the variation field.
\end{proof}

Along a fixed coordinate curve the Jacobi equation, viewed as an
equation for the pair $(J,\nabla J)$, is a first-order linear
system, and the basic solvability theory follows.

\begin{lemma}[Jacobi fields along a coordinate curve]
\label{lem:jacobi-field-coordinates}
\lean{MorganTianLib.chartCurvatureEndo,MorganTianLib.JacobiFieldOn,MorganTianLib.jacobiPairCoeffCoord,MorganTianLib.exists_isJacobiFieldOn_Icc,MorganTianLib.exists_isJacobiFieldOn_Icc_of_curve,MorganTianLib.continuousOn_jacobiPairCoeffCoord,MorganTianLib.IsJacobiFieldOn.eqOn_of_left,MorganTianLib.IsJacobiFieldOn.eqOn_of_right,MorganTianLib.IsJacobiFieldOn.add,MorganTianLib.IsJacobiFieldOn.const_smul,MorganTianLib.IsJacobiFieldOn.eqOn_zero}
\uses{lem:covariant-commutation-jacobi, def:riemann-curvature-tensor, thm:levi-civita-connection, lem:second-order-linear-ode}



Let $u\colon[a,b]\to U$ be a coordinate curve in a coordinate chart $U$.
Assume that $u$ and its coordinate derivative are continuous on $[a,b]$,
with $\Gamma$ and ${\mathcal R}$ as in
Lemma~\ref{lem:covariant-commutation-jacobi}, the coefficient
functions being continuous along $u$. Say that a pair of curves
$J,DJ\colon[a,b]\to\mathbb{R}^n$ is a {\em Jacobi field along $u$}
(with covariant derivative $DJ$) if
$$J'=DJ-\Gamma_u(\dot u,J),\qquad
DJ'=-{\mathcal R}(J,\dot u)\dot u-\Gamma_u(\dot u,DJ)$$
on $[a,b]$ (one-sided derivatives at the endpoints) — equivalently,
$\nabla J=DJ$ and $\nabla DJ+{\mathcal R}(J,\dot u)\dot u=0$, the
Jacobi equation. Then:
\begin{enumerate}
\item (Existence) for every pair $(J_0,DJ_0)$ there is a Jacobi
field along $u$ with $J(a)=J_0$, $DJ(a)=DJ_0$.
\item (Uniqueness) two Jacobi fields along $u$ agreeing (in both
components) at $a$, or at $b$, agree on all of $[a,b]$.
\item (Superposition) sums and scalar multiples of Jacobi fields
along $u$ are Jacobi fields; a Jacobi field with $J(a)=0$,
$DJ(a)=0$ vanishes identically, so one with nontrivial initial
data does not vanish identically.
\end{enumerate}
\end{lemma}

\begin{proof}
\sketch
\uses{lem:covariant-commutation-jacobi,lem:second-order-linear-ode}
Combine $J$ and its covariant derivative into one vector $W=(J,DJ)$. The coordinate Jacobi equation is a first-order linear system $W'=A(t)W$ with continuous operator coefficient. The preceding ODE lemma gives existence and two-sided uniqueness, while linearity of $A(t)$ gives superposition and the zero-data conclusion.
\end{proof}

The coefficient operator ${\mathcal R}$ of the chart Jacobi system is
not merely an ad hoc combination of Christoffel symbols: it {\em is}
the Riemann curvature tensor of
Definition~\ref{def:riemann-curvature-tensor} read in the chart, and
in particular it inherits pointwise bounds from the sectional
curvature. This is the bridge through which the geometric hypothesis
``$K(P)\le K$'' enters the Jacobi equation.

\begin{lemma}[The curvature tensor in chart coordinates]
\label{lem:chart-curvature-coordinates}
\lean{MorganTianLib.curvatureFormAt_chartFrame,MorganTianLib.chartCurvature_basis,MorganTianLib.exists_chartFrame_leviCivita_christoffel_nhds,MorganTianLib.exists_contMDiff_eventuallyEq,MorganTianLib.cov_congr_apply_right,MorganTianLib.alg_curvature_le_of_sectionalCurvature_le,MorganTianLib.chartCurvature_pairing_le_of_sectionalCurvatureAt_le,MorganTianLib.chartCurvature_pairing_le_of_sectionalCurvatureAt_le'}
\uses{def:riemann-curvature-tensor, thm:levi-civita-connection, def:sectional-curvature, lem:covariant-commutation-jacobi, claim:curvature-symmetries-bianchi}
\leanok



Let $(U,x)$ be a coordinate chart of $(M,g)$, $p\in U$, and let
$G_{ml}=g(\partial_m,\partial_l)$ be the Gram matrix of the
coordinate frame. Let ${\mathcal R}$ be the coefficient curvature of
Lemma~\ref{lem:covariant-commutation-jacobi} built from the
Christoffel symbols $\Gamma^k_{ij}$ of $g$ in this chart, with
components
$${\mathcal R}^m_{ijk}
=\partial_i\Gamma^m_{jk}-\partial_j\Gamma^m_{ik}
+\sum_r\bigl(\Gamma^r_{jk}\Gamma^m_{ir}
-\Gamma^r_{ik}\Gamma^m_{jr}\bigr).$$
Then:
\begin{enumerate}
\item ({\em The bridge.}) For all coordinate vectors
$v,w,z,t\in\mathbb{R}^n$, the Riemann curvature $(0,4)$-tensor
${\mathcal R}iem$ of Definition~\ref{def:riemann-curvature-tensor}
satisfies, at $p$,
$${\mathcal R}iem\Bigl(\sum_a v^a\partial_a,\ \sum_b w^b\partial_b,\
\sum_c z^c\partial_c,\ \sum_d t^d\partial_d\Bigr)
\;=\;-\bigl\langle{\mathcal R}(v,w)z,\;t\bigr\rangle_G
\;=\;-\sum_{m,l}{\mathcal R}^m_{\,\cdot}v^iw^jz^k\,G_{ml}\,t^l,$$
where $\langle a,b\rangle_G=\sum_{m,l}G_{ml}a^mb^l$ is the Gram
pairing and the middle expression contracts
${\mathcal R}(v,w)z=\sum_m{\mathcal R}^m_{ijk}v^iw^jz^k\,e_m$.
\item ({\em The sectional bound.}) If $K(P)\le K$ for every
$2$-plane $P\subset T_pM$, then for all $J,u\in\mathbb{R}^n$
$$\bigl\langle{\mathcal R}(J,u)u,\;J\bigr\rangle_G\ \le\
K\bigl(\langle J,J\rangle_G\langle u,u\rangle_G
-\langle J,u\rangle_G^2\bigr),$$
and if moreover $K\ge0$ then
$\langle{\mathcal R}(J,u)u,J\rangle_G\le
K\,\langle J,J\rangle_G\,\langle u,u\rangle_G$.
\end{enumerate}
\end{lemma}

\begin{proof}
\sketch
\leanok
\uses{thm:levi-civita-connection,claim:curvature-symmetries-bianchi,def:sectional-curvature,lem:covariant-commutation-jacobi,def:riemann-curvature-tensor}
Expand the covariant derivatives in chart coordinates and compare with the coordinate curvature formula, proving that the chart operator represents the intrinsic tensor. Tensoriality permits replacing chart fields by a local orthonormal frame at the point. The algebraic curvature symmetries then polarize a sectional-curvature bound into the stated bilinear estimate; equivalence of the chart and Riemannian norms transfers the estimate back to coordinates.
\end{proof}

The Jacobi equation is best analyzed in a parallel frame along the
geodesic, whose existence and rigidity are recorded in the next
lemma.

\begin{lemma}[Parallel transport and parallel frames]
\label{lem:parallel-frame}
\lean{MorganTianLib.exists_parallelFrameAlong,MorganTianLib.IsParallelAlongOn.metricInner_eq}
\uses{def:riemannian-metric, thm:levi-civita-connection}



Let $(M,g)$ be a Riemannian manifold with Levi-Civita connection
$\nabla$, and let $\gamma\colon[a,b]\to M$ be a geodesic with $a<b$,
continuous on $[a,b]$. A vector field
$V$ along $\gamma$ is {\em parallel} if $\nabla_{\gamma'}V=0$ on
$[a,b]$ (one-sided derivatives at the endpoints). Then:
\begin{enumerate}
\item for every $v\in T_{\gamma(a)}M$ there is a parallel
field $V$ along $\gamma$ with $V(a)=v$; under the usual chart-level
regularity hypotheses the ODE gives uniqueness; more generally, for any
family $(v_i)_i$ of vectors in $T_{\gamma(a)}M$ there is a family
$(e_i)_i$ of parallel fields with $e_i(a)=v_i$;
\item the inner product of two parallel fields is constant:
$\langle V(t),W(t)\rangle_{\gamma(t)}
=\langle V(a),W(a)\rangle_{\gamma(a)}$ for all $t\in[a,b]$;
\item consequently a parallel family whose initial values form an
orthonormal basis of $T_{\gamma(a)}M$ remains an orthonormal basis
of $T_{\gamma(t)}M$ for every $t$.
\end{enumerate}
\end{lemma}

\begin{proof}
\sketch
\uses{thm:levi-civita-connection}
In a coordinate trivialization, parallel transport is a linear first-order ODE with continuous coefficients, so existence and uniqueness hold for every initial vector. Metric compatibility makes inner products of parallel fields constant. Transporting an orthonormal basis therefore yields a parallel orthonormal frame, and uniqueness gives linearity, endpoint inversion, and compatibility under restriction and concatenation.
\end{proof}

The analytic content of the theory of Jacobi fields is the following
standard lemma about second-order linear equations, which we state
over an arbitrary real Banach space; it will be applied both to
vector-valued solutions (Jacobi fields, in a parallel frame along the
geodesic) and to operator-valued ones (matrices of Jacobi fields, as
in Lemma~\ref{lem:geodesic-polar-form}).

\begin{lemma}[Second-order linear ODE: existence, uniqueness,
superposition]
\label{lem:second-order-linear-ode}
\lean{MorganTianLib.JacobiSolOn,MorganTianLib.exists_isJacobiSolOn_Icc,MorganTianLib.IsJacobiSolOn.eqOn_of_left,MorganTianLib.IsJacobiSolOn.eqOn_of_right,MorganTianLib.IsJacobiSolOn.add,MorganTianLib.IsJacobiSolOn.const_smul,MorganTianLib.IsJacobiSolOn.sub,MorganTianLib.IsJacobiSolOn.eqOn_zero}
\leanok


Let $F$ be a real Banach space, let $a\le b$, and let
$R\colon[a,b]\to L(F)$ be a continuous family of bounded operators.
Say that a pair of curves $y,v\colon[a,b]\to F$ {\em solves the
second-order linear equation} $y''+R(t)\,y=0$ on $[a,b]$ if $y'=v$
and $v'=-R(t)\,y$ on $[a,b]$ (one-sided derivatives at the
endpoints). Then:
\begin{enumerate}
\item (Existence) for every $(y_0,v_0)\in F\times F$ there is a
solution with $y(a)=y_0$ and $v(a)=v_0$.
\item (Uniqueness) two solutions that agree at $a$ (both in $y$ and
in $v$) agree on all of $[a,b]$; likewise two solutions that agree
at $b$.
\item (Superposition) sums, differences and scalar multiples of
solutions are solutions; in particular the solution with $y(a)=0$,
$v(a)=0$ is identically zero, so a solution with nontrivial initial
data does not vanish identically.
\end{enumerate}
\end{lemma}

\begin{proof}
\sketch
\leanok
Rewrite the second-order equation as a first-order linear system on the product Banach space. On a sufficiently short interval, its integral operator is a contraction; finitely many such intervals give global existence. Gronwall's inequality gives uniqueness forward and backward. Linearity yields superposition and linear dependence on the initial data.
\end{proof}

\subsection{Minimal geodesics}

\begin{definition}[Conjugate point]
\label{def:conjugate-point}
\lean{MorganTianLib.chartVectorRep,MorganTianLib.IsJacobiFieldAlongOn,MorganTianLib.IsConjugatePointAt}
\uses{def:geodesic, lem:jacobi-field-coordinates}



Let $\gamma$ be a geodesic beginning at $p\in M$. A {\sl Jacobi field
along $\gamma$} is a field of tangent vectors $J(t)\in T_{\gamma(t)}M$
along $\gamma$, together with its covariant derivative
$\nabla J(t)\in T_{\gamma(t)}M$, such that near every time, read in the
coordinates of a chart containing the nearby piece of $\gamma$, the pair
$(J,\nabla J)$ satisfies the chart Jacobi system of
Lemma~\ref{lem:jacobi-field-coordinates} — equivalently, $J$ satisfies
the Jacobi equation
$\nabla_{\dot\gamma}\nabla_{\dot\gamma}J
+{\mathcal R}(J,\dot\gamma)\dot\gamma=0$. The notion is chart-local, so
it is meaningful for geodesics that leave any single chart. For any
$t>0$ we say that $q=\gamma(t)$ is a {\sl conjugate point along
$\gamma$} if there is a non-zero Jacobi field
along $\gamma$ vanishing at $p$ and at $\gamma(t)$. (Morgan--Tian leave
the vanishing at $p$ implicit in the phrase ``geodesic beginning at
$p$''; their subsequent arguments use it — e.g.\ the proofs of
Proposition~\ref{prop:minimal-geodesic-no-conjugate} and
Lemma~\ref{lem:conjugate-sturm} — and we make it explicit.)
\end{definition}


\begin{proposition}[Minimal geodesics have no conjugate points]
\label{prop:minimal-geodesic-no-conjugate}
\lean{MorganTianLib.not_isConjugatePointAt_of_minimizing, MorganTianLib.indexForm_nonneg_of_minimizing, MorganTianLib.dist_eq_speed_mul_of_minimizing, MorganTianLib.minimalGeodesic_restrict_unique, MorganTianLib.eqOn_of_minimizing_geodesic_of_minimizing_geodesic, MorganTianLib.globalGeodesic_eqOn_of_minimizing}
\uses{def:conjugate-point}
\label{jacmin}



On a complete Riemannian manifold with its distance induced by the metric,
suppose that $\gamma\colon [0,1]\to M$ is a minimal geodesic. Then
for every $0<t<1$ the restriction of $\gamma$ to $[0,t]$ is the unique
minimal geodesic between its endpoints, and there is no non-zero Jacobi
field $J$ along $\gamma$ with both $J(0)=0$ and $J(t)=0$. Equivalently,
there are no conjugate points on $\gamma((0,1))$.

To make the hypotheses explicit, take a real-parameter extension of $\gamma$
that is geodesic on an interval $[a,b]$ with $a<0<1<b$, is continuous at
each point of that interval, and satisfies
$\sqrt{\operatorname{speedSq}(\gamma,0)}\le d(\gamma(0),\gamma(1))$;
retain also the complete-space and metric-distance assumptions above.
\end{proposition}

\begin{proof}
\sketch
\uses{def:geodesic,def:conjugate-point,claim:second-variation-minimal-geodesic,lem:index-form-negative-at-conjugate,def:exponential-map,claim:curvature-symmetries-bianchi}
Write $p=\gamma(0)$, $q=\gamma(1)$ and $L=L(\gamma)=d(p,q)$; since
$\gamma$ is a geodesic its speed is constant, $|\gamma'(t)|\equiv
L$. We use throughout that for a piecewise smooth path
$\sigma$ parameterized by $[0,1]$ one has $E(\sigma)\ge
\frac{1}{2}L(\sigma)^2$, with equality if and only if $|\sigma'|$ is
constant, and that consequently $E(\gamma)=\frac12 L^2\le
\frac12 L(\sigma)^2\le E(\sigma)$ for every piecewise smooth path
$\sigma$ from $p$ to $q$: that is, a minimal geodesic minimizes
energy among all piecewise smooth paths with the same endpoints,
parameterized by $[0,1]$. (The first variation and energy
computations earlier in this section extend to piecewise smooth
paths and variations by applying them on each smooth piece and
summing.)

\emph{Part 1: for $t_0<1$ the restriction $\gamma|_{[0,t_0]}$ is the
unique minimal geodesic between its endpoints.} Fix $0<t_0<1$.
First, $\gamma|_{[0,t_0]}$ is minimal: if some rectifiable curve
$\sigma$ from $\gamma(0)$ to $\gamma(t_0)$ had
$L(\sigma)<L(\gamma|_{[0,t_0]})$, then the concatenation of $\sigma$
with $\gamma|_{[t_0,1]}$ would be a curve from $p$ to $q$ of length
$L(\sigma)+(1-t_0)L<L$, contradicting $L=d(p,q)$.

Now suppose that $\mu\colon [0,t_0]\to M$ is a minimal geodesic from
$\gamma(0)$ to $\gamma(t_0)$ distinct from $\gamma|_{[0,t_0]}$. Being
minimal geodesics between the same endpoints, both have length
$d(\gamma(0),\gamma(t_0))=t_0L$, and both have constant speed, equal
to $L$. If we had $\mu'(t_0)=\gamma'(t_0)$, then $\mu$ and
$\gamma|_{[0,t_0]}$ would be two solutions of the (second-order)
geodesic equation agreeing in position and velocity at $t=t_0$, and
by the uniqueness statement of classical ODE theory they would
coincide, contrary to the assumption that they are distinct. Hence
$$\Delta:=\gamma'(t_0)-\mu'(t_0)\not=0.$$
Let $c\colon [0,1]\to M$ be the concatenation, equal to $\mu$ on
$[0,t_0]$ and to $\gamma$ on $[t_0,1]$. Then $c$ is a piecewise
smooth path from $p$ to $q$ of constant speed $L$, so
$$E(c)=\frac{1}{2}L^2=\frac{1}{2}d(p,q)^2,$$
i.e., $c$ minimizes energy among piecewise smooth paths from $p$ to
$q$. Choose a piecewise smooth vector field $Y$ along $c$ with
$Y(0)=0$, $Y(1)=0$ and $Y(t_0)=\Delta$ (for instance $\phi(t)V(t)$
with $V$ the parallel field along each piece with
$V(t_0)=\Delta$ and $\phi$ a smooth bump function vanishing at $0,1$
with $\phi(t_0)=1$), and consider the variation
$\tilde c(t,u)=\exp_{c(t)}(uY(t))$, defined for $|u|$ small. On each
of the two smooth pieces the first-variation computation from the
beginning of this section applies: writing $\widetilde X=\partial
\tilde c/\partial t$ and $\widetilde Y=\partial\tilde c/\partial u$,
and using $\nabla_{\widetilde Y}\widetilde X=\nabla_{\widetilde
X}\widetilde Y$ and integration by parts on each piece,
$$\frac{dE(c_u)}{du}\Bigl|_{u=0}\Bigr.
=\left(\int_0^{t_0}+\int_{t_0}^1\right)
\langle\nabla_{X}Y,X\rangle\,dt
=\Bigl[\langle Y,X\rangle\Bigr]_0^{t_0^-}
+\Bigl[\langle Y,X\rangle\Bigr]_{t_0^+}^{1}
-\int_0^1\langle Y,\nabla_XX\rangle\,dt.$$
The integral vanishes since both pieces of $c$ are geodesics, and
the boundary terms give
$$\frac{dE(c_u)}{du}\Bigl|_{u=0}\Bigr.
=\langle Y(t_0),\mu'(t_0)\rangle-\langle
Y(t_0),\gamma'(t_0)\rangle=-\langle \Delta,\Delta\rangle
=-|\Delta|^2<0.$$
Hence $E(c_u)<E(c)=\frac{1}{2}d(p,q)^2$ for all sufficiently small
$u>0$, and therefore $L(c_u)\le\sqrt{2E(c_u)}<d(p,q)$ for such $u$
--- a contradiction, since $c_u$ is a path from $p$ to $q$. Thus no
such $\mu$ exists, proving Part 1.

\begin{remark}[Corner-rigidity formulation of the uniqueness argument]
\label{rem:part1-corner-rigidity-route}
\lean{MorganTianLib.minimalGeodesic_restrict_unique, MorganTianLib.globalGeodesic_eqOn_of_minimizing}
\leanok

The same conclusion can be phrased as corner rigidity.  The broken curve
$c=\mu\star\gamma|_{[t_0,1]}$ is piecewise $C^1$, has constant speed $L$
on both legs, and realizes $d(p,q)=L$.  A minimizing broken geodesic has
no corner, so $c$ satisfies the geodesic equation at the vertex $t_0$ as
well.  It is therefore a genuine geodesic on $(0,1)$.  Since it agrees
with $\gamma$ on $(t_0,1)$, uniqueness for the geodesic initial-value
problem propagates the agreement across the vertex and gives
$\mu=\gamma$ on $[0,t_0]$.
\end{remark}

For the second assertion of the proposition we shall use the
following claim, together with a refinement of the second-variation
formula valid for piecewise smooth variations.

\begin{lemma}[Second variation is zero iff Jacobi field]
\label{claim:second-variation-minimal-geodesic}
\lean{MorganTianLib.indexForm_self_eq_zero_iff_isJacobiFieldAlongOn_of_minimizing,MorganTianLib.metricIndexIntegral_frameLift_eq,MorganTianLib.deriv_deriv_dcEnergy_eq_zero_iff_isJacobiFieldAlongOn_of_minimizing,MorganTianLib.variationalField_endpoints_eq_zero_of_proper}
\uses{def:geodesic}

Let $\gamma$ be a minimizing geodesic on $[0,1]$ in a complete
Riemannian manifold whose metric induces the distance. Extend it as a
geodesic to $[a,b]$ with $a<0<1<b$, and choose a parallel orthonormal
frame on this interval. Let $\tilde\gamma(t,u)$ be a proper $C^3$
fixed-endpoint variation of $\gamma$ whose second-variation data satisfy
the energy--index identity and whose variation field and covariant
$t$-derivative are represented in that frame by coefficient fields
$W$ and $W'$, with $W\in C^3$ and $W(0)=W(1)=0$. Then
$$\frac{d^2E(\gamma_u)}{du^2}\Bigl|_{u=0}\Bigr.=0$$
if and only if the represented variation field is a Jacobi field along
$\gamma$.
\end{lemma}

\begin{proof}
\sketch
\uses{def:geodesic}
The assumed energy--index identity identifies the second derivative of
energy with the index form of the frame coefficient field $W$. Since
$\gamma$ is minimizing, this quadratic form is nonnegative on
endpoint-zero fields. Polarization shows that a zero value is
equivalent to vanishing pairing with every endpoint-zero test field;
the second-order ODE lemma then identifies this null space with the
Jacobi solutions. The parallel-frame identities transfer the resulting
solution back to the represented variation field along $\gamma$.
\end{proof}


\emph{Part 2: there are no conjugate points on $\gamma([0,1))$.}
For a piecewise smooth vector field $Y$ along $\gamma$ vanishing at
both endpoints, define the \emph{index form}
$$I(Y_1,Y_2)=\int_0^1\left(\langle\nabla_XY_1,\nabla_XY_2\rangle
-\langle{\mathcal R}(Y_1,X)X,Y_2\rangle\right)dt,$$
the integral being taken piecewise; $I$ is a symmetric bilinear form
on the vector space of such fields (symmetry of the curvature term
is the pair symmetry of the curvature tensor). For \emph{smooth}
$Y_1,Y_2$ vanishing at the endpoints, integration by parts gives
$I(Y_1,Y_2)=-\int_0^1\langle \Jac(Y_1),Y_2\rangle
dt=B_\gamma(Y_1,Y_2)$, the second-variation form of
Claim~\ref{claim:second-variation-minimal-geodesic}.

We first extend the second-variation formula: for every piecewise
smooth $Y$ along $\gamma$ vanishing at both endpoints, the variation
$\tilde\gamma(t,u)=\exp_{\gamma(t)}(uY(t))$ satisfies
\begin{equation}\label{secvarindex}
\frac{d^2E(\gamma_u)}{du^2}\Bigl|_{u=0}\Bigr.=I(Y,Y).
\end{equation}
Indeed, with $\widetilde X=\partial\tilde\gamma/\partial t$ and
$\widetilde Y=\partial\tilde\gamma/\partial u$ we have, on each
smooth piece,
$$\frac{d^2E(\gamma_u)}{du^2}
=\frac{d}{du}\int_0^1\langle\nabla_{\widetilde X}\widetilde
Y,\widetilde X\rangle\,dt
=\int_0^1\left(\langle\nabla_{\widetilde Y}\nabla_{\widetilde
X}\widetilde Y,\widetilde X\rangle
+|\nabla_{\widetilde X}\widetilde Y|^2\right)dt,$$
using $\nabla_{\widetilde Y}\widetilde X=\nabla_{\widetilde
X}\widetilde Y$ twice. Commuting the covariant derivatives,
$\nabla_{\widetilde Y}\nabla_{\widetilde X}\widetilde
Y=\nabla_{\widetilde X}\nabla_{\widetilde Y}\widetilde Y+{\mathcal
R}(\widetilde Y,\widetilde X)\widetilde Y$. For each fixed $t$ the
curve $u\mapsto\exp_{\gamma(t)}(uY(t))$ is a geodesic with velocity
field $\widetilde Y$, so $\nabla_{\widetilde Y}\widetilde Y=0$
identically. Evaluating at $u=0$ (where $\widetilde X=X=\gamma'$ and
$\widetilde Y=Y$) and using
$\langle{\mathcal R}(Y,X)Y,X\rangle=-\langle{\mathcal
R}(Y,X)X,Y\rangle$, which follows from the symmetries of the
curvature tensor, we obtain
Equation~(\ref{secvarindex}). Moreover, since $\gamma$ minimizes
energy among piecewise smooth paths from $p$ to $q$ and
$\gamma_u$ is such a path for each small $u$, the function
$u\mapsto E(\gamma_u)$ has a minimum at $u=0$, so
\begin{equation}\label{indexnonneg}
I(Y,Y)\ge 0
\end{equation}
for every piecewise smooth $Y$ along $\gamma$ vanishing at both
endpoints. By Equations~(\ref{secvarindex}) and
(\ref{indexnonneg}) and the Cauchy--Schwarz inequality for
positive semi-definite symmetric bilinear forms, if $I(Y,Y)=0$ then
$I(Y,Z)=0$ for every piecewise smooth $Z$ vanishing at both
endpoints. (For smooth fields, Equation~(\ref{secvarindex})
together with
Claim~\ref{claim:second-variation-minimal-geodesic} recovers the
statement that the null space of $I$ on smooth fields is the space
of Jacobi fields.)

\begin{remark}[Chart-local construction of the second variation]
\label{rem:second-variation-chart-route}
\lean{MorganTianLib.isMetricCompatibleAt_chartMetricBilin, MorganTianLib.secondVariation_energyDensity, MorganTianLib.exists_brokenVariationData, MorganTianLib.hasDerivAt_pieceEnergy_chartMetricBilin}
\leanok
The second variation is a local identity between derivatives of the energy
density.  In a chart, a two-parameter family is a map into the model
space, the connection is expressed by the Christoffel contraction
$\Gamma$, and the metric by its Gram form $G$.  The identity follows from
torsion-freeness, curvature commutation, and metric compatibility
(Lemma~\ref{lem:chart-metric-compatibility}).  Any $C^3$ family with the
prescribed first-order variation field may be used.  For a curve crossing
several charts, construct the variation piecewise; the boundary terms at
adjacent pieces telescope because the covariant acceleration is intrinsic.
\end{remark}

\begin{lemma}[Metric compatibility in a chart]
\label{lem:chart-metric-compatibility}
\lean{MorganTianLib.isMetricCompatibleAt_chartMetricBilin}
\uses{thm:levi-civita-connection}
\leanok



Let $G_{ij}$ be the Gram matrix of $g$ in a chart and $\Gamma^k_{ij}$ its
Christoffel symbols.  Then, for all vectors $X,V,W$ of the model space,
$$(\partial_X G)(V,W)=G(\Gamma(X,V),W)+G(V,\Gamma(X,W)),$$
i.e. $\nabla g=0$ in coordinates.
\end{lemma}

\begin{proof}
\sketch
\leanok
Contracting the Christoffel formula
$\Gamma^m_{ki}=\tfrac12 G^{ml}(\partial_kG_{li}+\partial_iG_{lk}-\partial_lG_{ki})$
against $G_{am}$ collapses the inverse Gram factor ($G_{am}G^{ml}=\delta_a^l$) and
gives the Christoffel symbol of the first kind,
$$\textstyle\sum_m G_{am}\Gamma^m_{ki}
=\tfrac12\left(\partial_kG_{ai}+\partial_iG_{ak}-\partial_aG_{ki}\right).$$
Adding this identity to itself with the roles of $i$ and $k$ exchanged, the two
$\pm\tfrac12\partial$-terms that are antisymmetric in $(i,k)$ cancel and the two
copies of $\tfrac12\partial_kG_{ai}$ (resp. $\tfrac12\partial_iG_{ak}$) combine, leaving
$$\textstyle\sum_m\left(G_{mj}\Gamma^m_{ki}+G_{im}\Gamma^m_{kj}\right)=\partial_kG_{ij},$$
which is the assertion read on coordinate vectors; both sides are bilinear in
$(V,W)$ and linear in $X$, so it holds for all $X,V,W$.
\end{proof}

\begin{lemma}[Second variation of energy, in a chart]
\label{lem:second-variation-energy-chart}
\lean{MorganTianLib.secondVariation_energyDensity}
\uses{def:geodesic, lem:chart-metric-compatibility, claim:curvature-symmetries-bianchi, lem:covariant-commutation-jacobi}
\leanok



Let $u$ be a $C^3$ two-parameter family of curves read in a chart, write
$X=\partial_t u$, $Y=\partial_s u$, and suppose the $t$-line through the point in
question is a geodesic there, $\nabla_tX=0$.  Then the energy density
$e=\tfrac12 G(X,X)$ satisfies
$$\frac{\partial^2 e}{\partial s^2}
=\frac{\partial}{\partial t}\,G\!\left(\nabla_s\partial_s u,\;X\right)
+G(\nabla_tY,\nabla_tY)-G\!\left({\mathcal R}(Y,X)X,\;Y\right).$$
The last two terms are exactly the index-form integrand $|\nabla_XY|^2-\langle
{\mathcal R}(Y,X)X,Y\rangle$; the first is a total $t$-derivative, i.e. a boundary
term.
\end{lemma}

\begin{proof}
\sketch
\uses{lem:chart-metric-compatibility,claim:curvature-symmetries-bianchi,lem:covariant-commutation-jacobi}
\leanok
Differentiate the coordinate energy twice. Symmetry of the Christoffel symbols and metric compatibility combine the derivative terms into covariant derivatives, while the remaining first derivatives of the connection assemble into curvature. Integration by parts removes the endpoint terms because the variation field vanishes there, leaving exactly the index form in the statement.
\end{proof}

The remaining step isolates all of the conjugate-point input.  It is a statement
about the index form alone: it uses neither the energy functional, nor a
variation of curves, nor the exponential map — only the Jacobi equation, the pair
symmetry of the curvature tensor, and integration by parts.

\begin{lemma}[A conjugate point makes the index form indefinite]
\label{lem:index-form-negative-at-conjugate}
\lean{MorganTianLib.exists_indexForm_neg_smooth_of_isConjugatePointAt}
\uses{def:conjugate-point, def:geodesic, lem:second-order-linear-ode, lem:parallel-frame, claim:curvature-symmetries-bianchi}
\leanok



Let $\gamma$ be a geodesic defined on an interval $[a,b]\supseteq[0,1]$,
with $a<0$ and $1<b$, and suppose that for some $t_0\in(0,1)$ the point
$\gamma(t_0)$ is conjugate to $\gamma(0)$ along $\gamma$. Then there is a
parallel orthonormal frame along $\gamma$ and coefficient fields $W_0,W_1$
which are $C^3$ on the two open pieces, satisfy
$W_0(0)=0$, $W_1(1)=0$, and agree at $t_0$, such that the sum of the two
index-form integrals on $[0,t_0]$ and $[t_0,1]$ is negative. Thus the
resulting piecewise-smooth field has negative index form.
\end{lemma}

\begin{proof}
\sketch
\uses{def:conjugate-point,lem:second-order-linear-ode,lem:parallel-frame,claim:curvature-symmetries-bianchi}
\leanok
Start with a nonzero Jacobi field vanishing at the conjugate endpoints. Truncate it just before the second endpoint and connect its remaining value to zero over a short interval. Integration by parts and the Jacobi equation cancel the interior contribution to the index form, leaving a boundary term that is negative for a sufficiently short truncation. Smooth approximation preserves the sign and the endpoint conditions.
\end{proof}

We can now finish Part 2. Suppose that for some $0<t_0<1$ the point
$\gamma(t_0)$ is a conjugate point along $\gamma$. By
Lemma~\ref{lem:index-form-negative-at-conjugate} there is a piecewise smooth
field $W$ along $\gamma$, vanishing at both endpoints, with $I(W,W)<0$ --- which
contradicts Equation~(\ref{indexnonneg}). Hence there is no conjugate point
$\gamma(t_0)$ with $t_0<1$. (There is no conjugate point at $t_0=0$ by
definition, conjugate points occurring only at parameters $t>0$.)
\end{proof}






\subsection{The exponential mapping}


\begin{definition}[Exponential map]
\label{def:exponential-map}
\lean{MorganTianLib.expMap,MorganTianLib.expDomain,MorganTianLib.zero_mem_expDomain,MorganTianLib.expDomain_nonempty}
\uses{def:geodesic}



For any $p\in M$, we can define the \textit{exponential
map} at $p$, $\exp_{p}$. Its natural domain $O_p$ contains the origin,
and on that domain $\exp_{p}(v)=\gamma_v(1)$, the endpoint of
the unique geodesic $\gamma_v\colon [0,1]\rightarrow M$ starting
from $p$ with initial velocity vector $v$. We always take
$O_{p}\subset T_{p}M$ to be the maximal domain on which $\exp_p$ is
defined, so that $O_p$ is a star-shaped open neighborhood of
$0\in T_pM$.
\end{definition}

By the Hopf--Rinow Theorem (Theorem~\ref{thm:hopf-rinow} below),
if $M$ is complete then the exponential map is defined on all of
$T_pM$.

Near the origin of $T_pM$ the exponential map is a diffeomorphism
onto a metric ball, uniformly over compact sets of centers. (The
forward reference to the Gauss Lemma,
Lemma~\ref{lem:geodesic-polar-form} in the section on curvature
comparison below, involves no circularity: that lemma depends only
on the material of the present section and of
Section~2.)

\begin{lemma}[Normal balls, uniformly on compact sets]
\label{lem:normal-ball}
\uses{def:exponential-map, def:geodesic}
\notready

Let $(M,g)$ be a Riemannian manifold, not assumed complete, and let
$K\subset M$ be compact. There is $\delta=\delta(K)>0$ such that
for every $x\in K$ the exponential map $\exp_x$ is defined on
$B(0,\delta)\subset T_xM$ and:
\begin{enumerate}
\item $\exp_x$ is a diffeomorphism from $B(0,\delta)$ onto the
metric ball $B(x,\delta)$, and $\exp_x(B(0,s))=B(x,s)$ for every
$0<s\le\delta$;
\item $d(x,\exp_x(v))=|v|$ for all $|v|<\delta$;
\item every piecewise smooth curve from $x$ to $\exp_x(v)$,
$|v|<\delta$, has length $\ge|v|$, with equality if and only if it
is a monotone reparameterization of the radial segment
$t\mapsto\exp_x(tv)$, $t\in[0,1]$; and every piecewise smooth
curve starting at $x$ and not staying in $B(x,s)$, $0<s\le\delta$,
has length $\ge s$.
\end{enumerate}
\end{lemma}

\begin{proof}
\sketch
\uses{lem:geodesic-polar-form,def:exponential-map,def:geodesic}
Write the geodesic equation as a first-order system on the tangent bundle. Smooth dependence on initial data gives, near every zero vector, a common time interval and a neighborhood on which the endpoint map is smooth. The inverse function theorem at each zero vector gives normal neighborhoods. A finite subcover of the compact set makes the time, radius, and chart choices uniform; uniqueness of ODE solutions gives the stated injectivity and compatibility properties.
\end{proof}

In particular, for every $p\in M$ there exists
$r_0=r_{0}(p,M)>0$ such that the restriction of $\exp_p$ to the
ball $B_{g|T_pM}(0,r_0)$ in $T_pM$ is a diffeomorphism onto the
metric ball $B_g(p,r_0)$. Fix $g$-orthonormal linear coordinates
on $T_pM$. Transferring these coordinates via $\exp_p$ to
coordinates on $B(p,r_0)$ gives us {\sl Gaussian normal
coordinates} on $B(p,r_0)\subset M$.

The differential of the exponential mapping is given by Jacobi
fields, as follows.

\begin{lemma}[The differential of the exponential map via Jacobi fields]
\label{lem:exponential-differential-jacobi}
\lean{MorganTianLib.expDifferential_eq_jacobiField,MorganTianLib.expDifferential_injective_iff_not_conjugate}
\uses{def:exponential-map, def:geodesic, def:conjugate-point}



Assume the complete Riemannian-distance setting used for the global
exponential map. For $p\in M$ and $v\in T_pM$, choose a terminal chart
$\zeta$ around $\exp_p(v)$. The chart differential $D$ of
$w\mapsto\operatorname{extChartAt}(\zeta)(\exp_p(w))$ and the Jacobi field
$Y_Z$ with $Y_Z(0)=0$ and $\nabla_XY_Z(0)=Z$ satisfy, in that chart,
\[
D Z=[Y_Z(1)]_\zeta,
\]
where $[\,\cdot\,]_\zeta$ denotes the coordinate representative in the
chosen terminal chart.
Under the usual tangent-space/chart identifications this is the intrinsic
formula $d(\exp_p)_v(Z)=Y_Z(1)$. In these coordinates, $D$ is injective if
and only if $\exp_p(v)$ is not conjugate to $p$ along $\gamma_v$.
\end{lemma}

\begin{proof}
\sketch
\uses{def:conjugate-point,def:geodesic,thm:levi-civita-connection,def:exponential-map,lem:second-order-linear-ode,lem:parallel-frame,lem:covariant-commutation-jacobi,lem:jacobi-field-coordinates}
Vary the initial vector by $v+sw$ and differentiate the corresponding geodesic family. The variation field is Jacobi, vanishes at time zero, and has initial covariant derivative $w$. Evaluating at time one identifies it with $d(\exp_p)_v(w)$. The converse follows from uniqueness of the Jacobi initial-value problem, and the radial direction is obtained by varying the speed of the same geodesic.
\end{proof}

\begin{corollary}[Exponential map is local diffeomorphism]
\label{cor:exponential-local-diffeomorphism}
\lean{MorganTianLib.expMapGlobal_localDiffeo_of_minimizing}
\uses{prop:minimal-geodesic-no-conjugate}
\label{star}



Suppose that $\gamma$ is a non-constant minimal geodesic parameterized by
$[0,1]$ starting at $p$, and write $X(0)=\gamma'(0)$. For each
$0<t_0<1$, the restriction $\gamma|_{[0,t_0]}$ is minimizing and the
exponential map is locally a diffeomorphism near $t_0X(0)$. Assume in
addition that $X(0)\ne0$, that $M$ is complete, and that the metric distance
is the Riemannian distance induced by $g$.
\end{corollary}

\begin{proof}
\sketch
\uses{prop:minimal-geodesic-no-conjugate,lem:exponential-differential-jacobi}
Of course, $\exp_p(t_0X(0))=\gamma(t_0)$. By
Lemma~\ref{lem:exponential-differential-jacobi} (applied to the
geodesic $\gamma|_{[0,t_0]}=\gamma_{t_0X(0)}$, reparameterized by
$[0,1]$), the differential of the
exponential mapping at $t_0X(0)$ is singular if and only if there
is a non-zero
Jacobi field along $\gamma|_{[0,t_0]}$ vanishing at $0$ and at
$\gamma(t_0)$. According
to  Proposition~\ref{prop:minimal-geodesic-no-conjugate} the only such Jacobi field is the
trivial one. Hence, the differential of $\exp_p$ at $t_0X(0)$
is an isomorphism, completing the proof.
\end{proof}

Two fundamental global consequences of this local theory are the
existence of minimizing geodesics and the Hopf--Rinow theorem.

\begin{lemma}[Existence of minimizing geodesics]
\label{lem:minimizing-geodesic-exists}
\uses{def:geodesic, def:exponential-map}
\notready

Let $(M,g)$ be a Riemannian manifold, not assumed complete or
connected (with the convention that $d(p,q)=\infty$ when there is
no smooth path from $p$ to $q$), let $p\in M$ and $D>0$. Suppose
that every unit-speed
geodesic starting at $p$ is defined on $[0,D]$. Then every $q\in
M$ with $d(p,q)\le D$ is joined to $p$ by a geodesic of length
$d(p,q)$; equivalently, $q=\exp_p(w)$ for some $w\in T_pM$ with
$|w|=d(p,q)$.
\end{lemma}

\begin{proof}
\sketch
\uses{lem:normal-ball,def:geodesic,def:exponential-map}
Choose a sequence of curves whose lengths converge to the distance and replace each by a broken geodesic with uniformly short legs inside normal balls. Arzela--Ascoli yields a uniformly convergent subsequence, while the local geodesic equation and compactness give convergence of the corresponding initial velocities. The limit is a geodesic from $p$ to $q$, and lower semicontinuity of length makes its length equal to $d(p,q)$.
\end{proof}

We can now state and prove the classical theorem of Hopf and Rinow
(compare Theorem 16 of Chapter 5 on p.~137 of \cite{Petersen}).

\begin{theorem}[Hopf-Rinow Theorem]
\label{thm:hopf-rinow}
\lean{MorganTianLib.hopfRinow_geodesic_complete, MorganTianLib.hopfRinow_minimizing_geodesic, MorganTianLib.hopfRinow_properSpace}
\uses{def:geodesic}
\leanok

(Hopf-Rinow) Let $(M,g)$ be a Riemannian manifold which is
complete as a metric space; for parts (2) and (3) assume in
addition that $M$ is connected. Then:
\begin{enumerate}
\item every geodesic extends to a geodesic defined for all time;
\item any two points $p,q\in M$ are joined by a minimizing
geodesic, of length $d(p,q)$;
\item every closed and bounded subset of $M$ is compact.
\end{enumerate}
\end{theorem}

\begin{proof}
  \leanok
\sketch
\uses{def:geodesic,def:exponential-map,lem:minimizing-geodesic-exists,thm:levi-civita-connection}
Metric completeness lets successive short minimizing segments converge, so a maximal geodesic cannot stop at finite time. Geodesic completeness and the open-and-closed argument for points reachable by minimizing geodesics then give a minimizer between any two points. Conversely, if every geodesic is defined for all time, the same argument supplies minimizing geodesics, and Cauchy sequences converge by placing their tails in a normal ball. The remaining compactness and properness formulations follow from closed bounded balls and the exponential image of compact tangent balls.
\end{proof}

Suppose now that $M$ is complete, and fix a point $p\in M$. By
Theorem~\ref{thm:hopf-rinow}(2), every $q\in M$ is the endpoint
of a minimizing geodesic from $p$ parameterized by $[0,1]$;
consequently, $\exp_p\colon T_pM\to M$ is onto.

\begin{definition}[Cut locus]
\label{def:cut-locus}
\lean{MorganTianLib.bookSegmentDomain,MorganTianLib.bookCutLocus}
\uses{def:exponential-map}
\leanok

Let $M$ be complete and connected, and let $p\in M$. Let $U_p\subset T_pM$ denote
the set of all $v\in T_pM$ for which: (i) $\gamma_v$ is the unique
minimal geodesic from $p$ to $\gamma_v(1)$, and (ii) $\exp_p$ is
a local diffeomorphism at $v$. We set ${\mathcal C}_p\subset M$
equal to $M\setminus \exp_p(U_p)$. Then ${\mathcal C}_p$ is
called the {\sl cut locus from $p$}.
\end{definition}

The basic properties of $U_p$ and ${\mathcal C}_p$ --- that $U_p$
is an open, star-shaped neighborhood of the origin and that the
cut locus is closed and of measure zero --- are established in the
following lemma. Its proof rests on the localized cut-locus
apparatus of Lemma~\ref{lem:localized-cut-locus} in the section on
curvature comparison below; there is no circularity: although that
lemma is typeset later, neither it nor anything in its chain of
dependencies makes reference to the cut locus or to the present
lemma.

The definition of $U_p$ above is the one Morgan-Tian use, but it is not the most
convenient starting point: both of its clauses are statements about geodesics and
differentials. The equivalent \emph{metric} formulation below is the one that is
elementary enough to be taken as a definition, and it is the one from which the
properties of $U_p$ are most easily derived.

\begin{definition}[Minimizing times and the cut time]
\label{def:cut-time}
\lean{MorganTianLib.IsMinimizingUpTo,MorganTianLib.minimizingTimes,MorganTianLib.cutTime}
\uses{def:exponential-map, def:geodesic}
\leanok



Let $M$ be complete, $p\in M$ and $v\in T_pM$, and let $\gamma_v(t)=\exp_p(tv)$ be the
radial geodesic. Say that $\gamma_v$ is \emph{minimizing up to time $t\ge 0$} if the
segment $\gamma_v|_{[0,t]}$ realises the distance between its endpoints,
$$d\bigl(p,\gamma_v(t)\bigr)=t\,\lvert v\rvert_g .$$
Write $T_v=\{t\ge0 : \gamma_v \text{ is minimizing up to } t\}$. The \emph{cut time} of
$p$ in the direction $v$ is
$$c(v)=\sup T_v\in[0,\infty],$$
the supremum being taken in $[0,\infty]$ because $c(v)=\infty$ in every direction of
$\R^n$ and of $H^n_k$.
\end{definition}

\begin{lemma}[The minimizing times form an interval]
\label{lem:minimizing-downward-closed}
\lean{MorganTianLib.IsMinimizingUpTo.mono,MorganTianLib.minimizingTimes_ordConnected,MorganTianLib.le_cutTime_iff,MorganTianLib.cutTime_radial_interval}
\uses{def:cut-time}
\leanok



If $\gamma_v$ is minimizing up to $t$ and $0\le s\le t$, then $\gamma_v$ is minimizing up
to $s$. Hence $T_v$ is downward closed and order-connected. The endpoint
criterion is attained: for $t\ge0$, $\gamma_v$ is minimizing up to $t$ exactly
when $t\le c(v)$; equivalently, for $t>0$, the radial vector $t v$ lies in
the segment domain exactly when $t<c(v)$ (with the usual $\mathbb{R}_{\ge0}^{\infty}$
interpretation). Thus the strict interval is open at the cut time, while
minimizing at the endpoint itself is recorded by the non-strict criterion.
\end{lemma}

\begin{proof}
\sketch
\uses{def:cut-time}
\leanok
The geodesic $\gamma_v$ has constant speed $\lvert v\rvert_g$, so for any $0\le a\le b$ the
segment $\gamma_v|_{[a,b]}$ has length $(b-a)\lvert v\rvert_g$ and therefore
$d(\gamma_v(a),\gamma_v(b))\le (b-a)\lvert v\rvert_g$. Applying this twice and using the
triangle inequality,
$$t\lvert v\rvert_g=d\bigl(p,\gamma_v(t)\bigr)
\le d\bigl(p,\gamma_v(s)\bigr)+d\bigl(\gamma_v(s),\gamma_v(t)\bigr)
\le s\lvert v\rvert_g+(t-s)\lvert v\rvert_g=t\lvert v\rvert_g .$$
Every inequality in the chain is therefore an equality; in particular
$d(p,\gamma_v(s))=s\lvert v\rvert_g$, which is the claim.
\end{proof}

\begin{definition}[The segment domain, the cut locus and the injectivity radius, metrically]
\label{def:cut-time-locus}
\lean{MorganTianLib.segmentDomain,MorganTianLib.cutLocus,MorganTianLib.injectivityRadius}
\uses{def:cut-time, lem:minimizing-downward-closed}
\leanok



With $c(v)$ as in Definition~\ref{def:cut-time}, set
$$U_p=\{v\in T_pM : c(v)>1\},\qquad
{\mathcal C}_p=\{\gamma_v(c(v)) : \lvert v\rvert_g=1,\ c(v)<\infty\},\qquad
\inj_M(p)=\inf_{\lvert v\rvert_g=1} c(v).$$
These agree with the $U_p$, ${\mathcal C}_p$ and $\inj_M(p)$ of
Definitions~\ref{def:cut-locus} and~\ref{def:injectivity-radius}; that agreement is part of
Lemma~\ref{lem:cut-locus-properties}.
\end{definition}

\begin{lemma}[$U_p$ is star-shaped]
\label{lem:segment-domain-star-shaped}
\lean{MorganTianLib.isMinimizingUpTo_smul,MorganTianLib.segmentDomain_smul_mem}
\uses{def:cut-time-locus, def:cut-time}
\leanok



If $v\in U_p$ and $0<\lambda\le1$ then $\lambda v\in U_p$. Thus $U_p$ is star-shaped about
the origin.
\end{lemma}

\begin{proof}
\sketch
\uses{def:cut-time,def:cut-time-locus}
\leanok
Since $\gamma_{\lambda v}(t)=\gamma_v(\lambda t)$ and
$\lvert \lambda v\rvert_g=\lambda\lvert v\rvert_g$ for $\lambda>0$, the condition
$d(p,\gamma_{\lambda v}(t))=t\lvert\lambda v\rvert_g$ is literally the condition
$d(p,\gamma_v(\lambda t))=(\lambda t)\lvert v\rvert_g$. Hence $\gamma_{\lambda v}$ is
minimizing up to $t$ iff $\gamma_v$ is minimizing up to $\lambda t$, so
$T_{\lambda v}=\lambda^{-1}T_v$ and $c(\lambda v)=c(v)/\lambda$. If $c(v)>1$ and
$0<\lambda\le 1$ then $c(\lambda v)=c(v)/\lambda\ge c(v)>1$, i.e. $\lambda v\in U_p$.
\end{proof}

\begin{lemma}[Properties of the cut locus]
\label{lem:cut-locus-properties}
\lean{MorganTianLib.bookSegmentDomain_eq_segmentDomain, MorganTianLib.bookCutLocus_eq_cutLocus, MorganTianLib.cutLocus_properties}
\uses{def:cut-locus, def:cut-time-locus, lem:segment-domain-star-shaped}
\leanok

Let $M$ be complete and connected, and let $p\in M$. Then:
\begin{enumerate}
\item $U_p$ is an open, star-shaped neighborhood of $0$ in
$T_pM$, and $d(p,\exp_p(w))=|w|$ for every $w\in U_p$;
\item $\exp_p(U_p)$ is open and the cut locus
${\mathcal C}_p=M\setminus\exp_p(U_p)$ is closed;
\item ${\mathcal C}_p$ has measure zero.
\end{enumerate}
\end{lemma}

\begin{proof}
\sketch
\uses{def:cut-locus,lem:localized-cut-locus,thm:hopf-rinow,lem:normal-ball}
Radial minimality is downward closed, which proves star-shapedness. A finite cut time is attained by continuity of distance. If the endpoint were neither conjugate nor joined by a second minimizing geodesic, the inverse function theorem and compactness of minimizing initial vectors would extend uniqueness and minimality beyond the cut time, a contradiction. These arguments also give continuity of the cut time along directions where the alternatives remain separated.
\end{proof}

\subsection*{Nullity of the cut locus, without Sard}

The proof of part (3) just given routes through
Lemma~\ref{lem:localized-cut-locus}(4), whose own proof invokes Sard's
theorem to discard the critical values of $\exp_p$, and which
moreover presupposes part (1) --- the openness of $U_p$, i.e.\ lower
semicontinuity of the cut time. Both are expensive: Sard's theorem is
unavailable in the equidimensional generality one would need for the
second family of cut points (those joined to $p$ by two minimizing
geodesics), and lower semicontinuity of $c$ is the hard half of the
Klingenberg dichotomy.

The following chain proves part (3) outright and uses neither. Its
observation is that the cut locus is the $\exp_p$-image of a
\emph{radial graph} --- a subset of $T_pM$ meeting each ray from the
origin at most once --- and that radial graphs are null for the
elementary reason that, in polar coordinates, their radial slices are
single points.

\begin{lemma}[The cut time is upper semicontinuous, hence measurable]
\label{lem:cut-time-measurable}
\lean{MorganTianLib.le_cutTime_iff,MorganTianLib.measurable_cutTime}
\uses{def:cut-time, def:exponential-map}
\leanok



Let $M$ be complete and $p\in M$. Then:
\begin{enumerate}
\item for every $t\ge0$, $\ \gamma_v$ is minimizing up to $t$ if and
only if $t\le c(v)$ --- that is, the supremum defining the cut time is
\emph{attained};
\item $c\colon T_pM\to[0,\infty]$ is measurable.
\end{enumerate}
\end{lemma}

\begin{proof}
\sketch
\uses{def:cut-time,lem:minimizing-downward-closed,def:exponential-map}
\leanok
If a direction has cut time below a fixed number, choose a slightly later point on its ray where the radial segment is not minimizing. Continuity of distance and of the exponential map preserves this strict failure for nearby directions. Hence every strict sublevel set of the cut time is open, so the cut time is upper semicontinuous and therefore Borel measurable.
\end{proof}

\begin{definition}[Cut vectors]
\label{def:cut-vectors}
\lean{MorganTianLib.cutVectors}
\uses{def:cut-time}
\leanok



The \emph{cut vectors} of $p$ are
$$\Sigma_p=\{v\in T_pM\ :\ v\neq0,\ c(v)=1\}.$$
\end{definition}

Since $c(\lambda v)=c(v)/\lambda$ for $\lambda>0$ (the proof of
Lemma~\ref{lem:segment-domain-star-shaped}, with
Lemma~\ref{lem:cut-time-measurable}(1) supplying the attainment needed
for the reverse inequality), $\Sigma_p$ is exactly the set of vectors
$c(u)\,u$ over the unit directions $u$ with $c(u)<\infty$: the radial
boundary of the star-shaped domain $U_p$. Writing it as the level set
$\{c=1\}$ rather than as a parameterized image is what makes it
visibly measurable.

\begin{lemma}[A measurable radial graph is null]
\label{lem:radial-graph-null}
\lean{MorganTianLib.IsRadialGraph,MorganTianLib.addHaar_eq_zero_of_isRadialGraph}
\leanok


Let $E$ be a finite-dimensional real normed space with Haar measure
$\mu$, and let $A\subseteq E$ be measurable with $0\notin A$. If $A$
meets every ray $\{ru : r>0\}$ ($\lvert u\rvert=1$) in at most one
point, then $\mu(A)=0$.
\end{lemma}

\begin{proof}
\sketch
\leanok
Polar coordinates identify $E\setminus\{0\}$, with its Haar measure,
with the product $S^{n-1}\times(0,\infty)$ carrying
$\sigma\otimes r^{n-1}dr$, the identification being $v\mapsto
(v/\lvert v\rvert,\lvert v\rvert)$. Under it $A$ becomes a measurable
set whose slice over each direction $u$ is, by hypothesis, at most a
single radius. The radial factor $r^{n-1}dr$ is absolutely continuous
with respect to Lebesgue measure and therefore has no atoms, so every
slice is null; by Fubini, $A$ is null.
\end{proof}

\begin{lemma}[$\exp_p$ carries null sets to null sets]
\label{lem:exp-preserves-null}
\lean{MorganTianLib.riemannianMeasure_eq_zero_of_chartImage_null, MorganTianLib.riemannianMeasure_image_eq_zero_of_null, MorganTianLib.riemannianMeasure_expMapGlobal_image_eq_zero}
\uses{def:exponential-map, lem:exponential-differential-jacobi}
\leanok



Let $M$ be complete, $p\in M$, and $N\subseteq T_pM$ a Lebesgue-null
set. Then $\exp_p(N)$ is null for the Riemannian measure $\mu_g$.
\end{lemma}

\begin{proof}
\sketch
\uses{def:exponential-map,lem:exponential-differential-jacobi}
\leanok
First, nullity is a chart-local and chart-independent notion: if
$A\subseteq M$ lies in a single chart and its coordinate image is
Lebesgue-null, then $\mu_g(A)=0$, since $\mu_g$ is
$\sqrt{\det g_{ij}}\,dx$ in that chart and a density cannot resurrect a
null set. (Formally one passes to a measurable null superset of the
coordinate image; $A$ need not itself be measurable.)

Now cover $M$ by countably many charts. On the piece of $\exp_p(N)$
lying in the chart $\alpha$, the coordinate image is contained in the
image of $N\cap\exp_p^{-1}(\text{source }\alpha)$ under the map
$x_\alpha\circ\exp_p\colon T_pM\to\R^n$, which is differentiable there
--- $\exp_p$ is differentiable at every $v$, in every chart around
$\exp_p(v)$, by Lemma~\ref{lem:exponential-differential-jacobi}
together with the smoothness of the chart transitions. A differentiable
map between spaces of the same finite dimension carries null sets to
null sets, so each piece is null, and $\exp_p(N)$ is a countable union
of null sets.
\end{proof}

\begin{lemma}[The cut locus is null]
\label{lem:cut-locus-null}
\lean{MorganTianLib.cutLocus_subset_image_cutVectors, MorganTianLib.isRadialGraph_cutVectors, MorganTianLib.measurableSet_cutVectors, MorganTianLib.riemannianMeasure_cutLocus_eq_zero}
\uses{def:cut-locus, def:cut-vectors, def:cut-time, lem:cut-time-measurable, lem:radial-graph-null, lem:exp-preserves-null}
\leanok



Let $M$ be complete and connected and $p\in M$. Then ${\mathcal C}_p$
has measure zero. This is part (3) of
Lemma~\ref{lem:cut-locus-properties}, proved without Sard's theorem and
without openness of $U_p$.
\end{lemma}

\begin{proof}
\sketch
\uses{def:cut-locus,def:cut-vectors,lem:cut-time-measurable,
  lem:radial-graph-null,lem:exp-preserves-null}
\leanok
By definition every cut point is $\gamma_u(c(u))=\exp_p(c(u)\,u)$ for
some unit $u$ with $c(u)<\infty$. If $c(u)=0$ the point is
$p=\exp_p(0)$. If $c(u)>0$ then, by the rescaling
$c(\lambda v)=c(v)/\lambda$, the vector $c(u)\,u$ has cut time
$c(u)/c(u)=1$ and is non-zero, so it lies in $\Sigma_p$. Hence
$${\mathcal C}_p\ \subseteq\ \exp_p\bigl(\{0\}\cup\Sigma_p\bigr).$$

The set $\Sigma_p$ is measurable, being $\{c=1\}\setminus\{0\}$ with
$c$ measurable (Lemma~\ref{lem:cut-time-measurable}). It is a radial
graph: if $r_1u$ and $r_2u$ both have cut time $1$, then applying the
rescaling with $\lambda=r_2/r_1$ to $r_1u$ gives
$1=c(r_2u)=c(r_1u)/\lambda=1/\lambda$, so $\lambda=1$ and $r_1=r_2$.
(This is nothing but the statement that a geodesic is cut \emph{once}.)
By Lemma~\ref{lem:radial-graph-null}, $\Sigma_p$ is Lebesgue-null, and
adjoining the origin changes nothing since Haar measure is atomless.
Finally Lemma~\ref{lem:exp-preserves-null} carries this null set to a
$\mu_g$-null subset of $M$ containing ${\mathcal C}_p$.
\end{proof}

\begin{proposition}[Exponential map to cut locus complement]
\label{prop:exponential-diffeomorphism-cut-locus}
\lean{MorganTianLib.bookSegmentDomain_eq_segmentDomain, MorganTianLib.bookCutLocus_eq_cutLocus, MorganTianLib.expMapGlobal_segmentDomain_diffeomorph}
\uses{def:cut-locus}
\leanok


The map
$$\exp_p\colon U_p\to M\setminus {\mathcal C}_p$$
is a diffeomorphism. \end{proposition}

\begin{proof}
\sketch
\uses{def:cut-locus,def:exponential-map,def:geodesic,lem:cut-locus-properties}
Every point outside the cut locus has a unique minimizing radial vector, which lies in $U_p$; this proves bijectivity. Such a vector is not conjugate before its cut time, so the exponential differential is invertible there. The inverse function theorem gives a smooth local inverse, and the local inverses agree by uniqueness, producing the global diffeomorphism.
\end{proof}

The domain $U_p$ has the following radial structure, which will be
used repeatedly in the volume arguments below.

\begin{lemma}[Radial structure of the domain $U_p$]
\label{lem:cut-time-star-shaped}
\lean{MorganTianLib.bookSegmentDomain_radial_properties}
\uses{def:cut-locus, lem:cut-locus-properties}
\leanok

Let $M$ be complete and connected, let $p\in M$, let
$U_p\subset T_pM$ be the
set of Definition~\ref{def:cut-locus} (open, by
Lemma~\ref{lem:cut-locus-properties}), and let $S$ denote the
unit sphere of $(T_pM,g_p)$. For $v\in S$ set
$$c(v)=\sup\{t>0\ |\ tv\in U_p\}\in(0,\infty].$$
Then:
\begin{enumerate}
\item $\{t>0\ |\ tv\in U_p\}=(0,c(v))$; in particular $U_p$ is
star-shaped about the origin;
\item for every $w\in U_p$ the geodesic $\gamma_w$ is the unique
minimal geodesic from $p$ to $\exp_p(w)$, and
$d(p,\exp_p(w))=|w|$;
\item for $v\in S$ the unit-speed radial geodesic
$\gamma_v(t)=\exp_p(tv)$ is minimizing on $[0,t]$ for every
$t<c(v)$, and there is no conjugate point along $\gamma_v$ on
$[0,c(v))$.
\end{enumerate}
\end{lemma}

\begin{proof}
\sketch
\uses{def:cut-locus,lem:cut-locus-properties,prop:minimal-geodesic-no-conjugate,def:conjugate-point,lem:exponential-differential-jacobi}
Homogeneity of the cut time gives the radial description of $U_p$. Downward closure of minimizing times shows that every shorter multiple of a vector in $U_p$ remains in $U_p$. Conversely, approaching the cut time along the ray identifies the radial boundary and yields the stated formula for the domain.
\end{proof}




\begin{definition}[Injectivity radius]
\label{def:injectivity-radius}
\lean{MorganTianLib.metricExpBallDiffeomorph, MorganTianLib.metricBookInjectivityRadius,MorganTianLib.injectivityRadius_eq_segmentDomainFrontierDistance,MorganTianLib.injectivityRadius_eq_cutLocusDistance}
\uses{def:exponential-map, def:cut-locus}

Under the completeness hypotheses needed for the distance identities below
(and connectedness for the frontier identity), the ball-based injectivity
radius is the supremum of the $r>0$ for which
the restriction of $\exp_p$ to the $g_p$-metric tangent ball $B(0,r)$ is a
diffeomorphism into $M$; denote this radius by
$\operatorname{inj}^{\rm ball}_M(p)$. Independently, the cut-time radius
$\inj_M^{\rm cut}(p)$ is the infimum of the cut times in unit directions.
For this latter radius, the frontier and cut-locus identities give
\[
\inj_M^{\rm cut}(p)=d_{T_pM}(0,\partial U_p)
  =d_M(p,{\mathcal C}_p).
\]
No equality between the ball-based and cut-time radii is assumed in this
definition.
\end{definition}

Suppose that  $\inj_M^{\rm cut}(p)=r$. There are two possibilities:
Either there is a broken, closed geodesic through $p$, broken only
at $p$, of length $2r$, or there is a geodesic $\gamma$ of length
$r$ emanating from $p$ whose endpoint is a conjugate point along
$\gamma$. The first case happens when the exponential mapping is not
one-to-one of the closed ball of radius $r$ in $T_pM$, and the
second happens when there is a tangent vector in $T_pM$ of length
$r$ at which $\exp_p$ is not a local diffeomorphism.
This dichotomy (due to Klingenberg) is made precise by the
following lemma.

\begin{lemma}[Klingenberg: short geodesic loop at the injectivity radius]
\label{lem:klingenberg-loop}
\uses{def:injectivity-radius, def:exponential-map, def:conjugate-point, def:geodesic, thm:hopf-rinow}
\notready

Let $M$ be complete and connected, $p\in M$, and suppose
$\rho:=\inj_M(p)<\infty$. If no geodesic from $p$ has a conjugate point at distance
$\le\rho$ along it, then there is a geodesic loop $\sigma\colon
[0,2\rho]\to M$ with $\sigma(0)=\sigma(2\rho)=p$, parameterized at
unit speed, smooth on $(0,2\rho)$ (i.e., broken at most at $p$), of
length exactly $2\rho$.
\end{lemma}

\begin{proof}
\sketch
\uses{def:injectivity-radius,def:conjugate-point,thm:hopf-rinow,lem:exponential-differential-jacobi}
Choose a cut vector of minimal length. If the exponential differential is singular there, the endpoint is conjugate and the first alternative holds. Otherwise local inverse branches at the two minimizing radial geodesics show that the terminal velocities must be opposite; otherwise one can shorten the broken geodesic and contradict minimality of the cut radius. Joining the two radial segments then gives a geodesic loop of length $2\inj_M(p)$ based at $p$.
\end{proof}

\subsection{Local isometries and constant curvature}

In this subsection we collect the facts about local isometries and
about manifolds of constant curvature that are used in the proof of
the Uniformization Theorem (Theorem~\ref{thm:uniformization}).

\begin{definition}[Local isometry]
\label{def:local-isometry}
\lean{MorganTianLib.LocalIsometry.IsLocalIsometry, MorganTianLib.LocalIsometry.IsLocalIsometry.of_contMDiff_preservesMetric, MorganTianLib.LocalIsometry.isLocalIsometry_iff_contMDiff_preservesMetric}
\uses{def:riemannian-metric}
\leanok

A smooth map $F\colon (N,h)\to (M,g)$ between Riemannian manifolds
of the same dimension is a {\sl local isometry} if $F^*g=h$, i.e.,
if $dF_x\colon T_xN\to T_{F(x)}M$ is a linear isometry for every
$x\in N$. By the inverse function theorem, a local isometry
restricts to an isometry from a sufficiently small neighborhood of
each point onto its image. An {\sl isometry} is a bijective local
isometry; its inverse is then also an isometry.
\end{definition}

\begin{lemma}[Local isometries and the exponential map]
\label{lem:local-isometry-exponential}
\lean{MorganTianLib.LocalIsometry.IsLocalIsometry.isGeodesicCurveOn_comp, MorganTianLib.LocalIsometry.IsLocalIsometry.expMapIntrinsic_naturality_of_complete}
\uses{def:local-isometry, def:geodesic, def:exponential-map}
\leanok

A local isometry $F\colon N\to M$ carries geodesics to geodesics.
If in addition $N$ is complete (as a metric space), then for every
$x\in N$, every $v\in T_xN$ and every $t\in \mathbb{R}$ the point
$\exp_{F(x)}\left(t\,dF_x(v)\right)$ is defined and
$$F\left(\exp_x(tv)\right)=\exp_{F(x)}\left(t\,dF_x(v)\right).$$
\end{lemma}

\begin{proof}
\sketch
\uses{thm:levi-civita-connection,thm:hopf-rinow}
Being a geodesic is a local property, and on a small neighborhood
$U$ of any point a local isometry is an isometry onto its image. An
isometry carries the Levi-Civita connection to the Levi-Civita
connection: the pushforward under an isometry of the Levi-Civita
connection of the source is a torsion-free connection making the
target metric parallel, hence is the Levi-Civita connection of the
target by the uniqueness in
Theorem~\ref{thm:levi-civita-connection}. Consequently
$\nabla_{\dot\gamma}\dot\gamma=0$ implies $\nabla_{(F\circ
\gamma)^{\textstyle .}}(F\circ\gamma)^{\textstyle .}=0$ on $U$, and
geodesics map to geodesics.

Now suppose $N$ is complete. By the Hopf--Rinow Theorem
(Theorem~\ref{thm:hopf-rinow}), the geodesic
$t\mapsto\exp_x(tv)$ is defined for all $t\in\mathbb{R}$. Its image
$t\mapsto F(\exp_x(tv))$ is then a geodesic of $M$, defined for all
$t$, with initial position $F(x)$ and initial velocity $dF_x(v)$. By
the uniqueness of geodesics with given initial conditions, it
coincides with $t\mapsto \exp_{F(x)}(t\,dF_x(v))$ wherever the
latter is defined; in particular the maximal geodesic with this
initial data is defined for all $t$, and the displayed identity
holds for all $t$.
\end{proof}

Combining this naturality with the uniform normal balls of
Lemma~\ref{lem:normal-ball} shows that a local isometry is, at a
uniform scale over any compact set, an isometry of metric balls.

\begin{lemma}[Local isometries on normal balls]
\label{lem:local-isometry-normal-balls}
\uses{def:local-isometry, def:exponential-map}
\notready

Let $F\colon(N,h)\to(M,g)$ be a local isometry and let $K\subset
N$ be compact. Then there is $\delta>0$ such that:
\begin{enumerate}
\item for every $x\in K$, $F$ maps the metric ball $B(x,\delta)$
diffeomorphically onto the metric ball $B(F(x),\delta)$, and
$$F\left(\exp_x(v)\right)=\exp_{F(x)}\left(dF_x(v)\right)
\qquad(|v|<\delta);$$
\item $d_M\left(F(a),F(b)\right)=d_N(a,b)$ for all $a,b\in
B(x,\delta/4)$, $x\in K$;
\item if $c$ is a continuous curve of finite length in $M$ and
$\tilde c$ is a continuous curve in $N$ with $F\circ\tilde c=c$
and $\tilde c([0,1])\subseteq K$, then the lengths agree:
$L_h(\tilde c)=L_g(c)$.
\end{enumerate}
\end{lemma}

\begin{proof}
\sketch
\uses{lem:normal-ball,lem:local-isometry-exponential,def:local-isometry,def:geodesic}
Apply the exponential-map identity for a local isometry at each preimage of the center. Uniform normal balls make the corresponding inverse branches well defined on one ball in the target. Distinct branches have disjoint images by uniqueness of geodesics and equality of their initial data. Every point in the relevant inverse image lies on one branch by lifting its unique radial geodesic, which gives the asserted decomposition and local inverse formulas.
\end{proof}

\begin{lemma}[Rigidity of local isometries]
\label{lem:local-isometry-rigidity}
\lean{MorganTianLib.LocalIsometry.IsLocalIsometry.eq_of_pointDatum}
\uses{def:local-isometry}
\leanok

Let $N$ be connected and let $F_1,F_2\colon N\to M$ be local
isometries with $F_1(x_0)=F_2(x_0)$ and $d(F_1)_{x_0}=d(F_2)_{x_0}$
for some $x_0\in N$. Then $F_1=F_2$.
\end{lemma}

\begin{proof}
  \leanok
\sketch
\uses{lem:local-isometry-exponential,def:exponential-map,lem:normal-ball}
Let $A=\{x\in N\,|\, F_1(x)=F_2(x)\text{ and }
d(F_1)_x=d(F_2)_x\}$. It is non-empty by hypothesis, and closed,
since $F_1,F_2$ and their differentials are continuous. It is also
open: given $x\in A$, choose $r>0$ so small that $\exp_x$ is a
diffeomorphism of $B(0,r)\subset T_xN$ onto $B(x,r)$ and both
geodesics $t\mapsto\exp_x(tv)$, $|v|<r$, $t\in[0,1]$, and their
images under $F_1,F_2$ are defined (such $r$ exists by
Lemma~\ref{lem:normal-ball} applied to $\{x\}\subset N$ and to
$\{F_1(x)\}\subset M$). Exactly as in the proof of
Lemma~\ref{lem:local-isometry-exponential}, for $|v|<r$,
$$F_i\left(\exp_x(v)\right)=\exp_{F_i(x)}\left(d(F_i)_x
(v)\right),\qquad i=1,2,$$
and since $F_1(x)=F_2(x)$ and $d(F_1)_x=d(F_2)_x$ the two right-hand
sides agree. Hence $F_1=F_2$ on the open set $B(x,r)$, and then also
$d(F_1)=d(F_2)$ there, so $B(x,r)\subset A$. Since $N$ is connected,
$A=N$.
\end{proof}

\begin{lemma}[Local isometries from complete manifolds are coverings]
\label{lem:local-isometry-covering}
\lean{MorganTianLib.LocalIsometry.IsLocalIsometry.surjectiveCovering}
\uses{def:local-isometry}

Let $F\colon N\to M$ be a local isometry. If the source $N$ is connected
and complete, the target $M$ is connected, and the source metric distance
agrees with the Riemannian distance induced by its metric, then $F$ is a
surjective covering map.
\end{lemma}

\begin{proof}
\sketch
\uses{lem:local-isometry-exponential,thm:hopf-rinow,def:exponential-map}
Completeness of the source and preservation of lengths let every short geodesic from a target point lift for its full parameter interval. Normal-ball inverse branches based at the points of one fiber are pairwise disjoint by uniqueness of lifted geodesics. Every point over a sufficiently small target ball belongs to one such branch, so the ball is evenly covered. Thus the local isometry is a covering map.
\end{proof}

We shall also need the classical Gauss formula identifying the
Levi-Civita connection of a submanifold of Euclidean space.

\begin{lemma}[Levi-Civita connection in an identified Euclidean patch]
\label{lem:submanifold-connection}
\lean{MorganTianLib.euclideanSubmanifold_connection_and_geodesic,MorganTianLib.IsEuclideanSubmanifoldGeodesic}
\uses{def:riemannian-metric, thm:levi-civita-connection, def:geodesic}

Let $D$ be the identified immersed-patch structure on an open subset of a
Euclidean space, with its tangent projection and normal spaces. For tangent
fields $X,Y$, the induced covariant term is the tangential projection of the
ambient flat derivative. For a curve whose velocity has the required local
smooth tangent extensions, its induced geodesic condition is equivalent to
its Euclidean second derivative lying in the corresponding normal space. This
is the local $D$-patch formulation; an embedded-submanifold version also
specifies an embedding and the induced tangent and normal spaces.
\end{lemma}

\begin{proof}
\sketch
\uses{thm:levi-civita-connection,def:geodesic,def:riemannian-metric}
Project the ambient derivative orthogonally to the tangent bundle and verify the connection axioms, metric compatibility, and vanishing torsion. Uniqueness of the Levi--Civita connection identifies this projection with the intrinsic connection. Differentiating an orthonormal tangent frame splits the ambient derivative into tangent and normal parts; applying this to a curve shows that it is geodesic exactly when its Euclidean acceleration is normal.
\end{proof}

We now solve the Jacobi equation explicitly in constant curvature.
For $\lambda\in\mathbb{R}$ let $s_\lambda\colon
\mathbb{R}\to\mathbb{R}$ be the solution of
$$s_\lambda''+\lambda\, s_\lambda=0,\qquad s_\lambda(0)=0,\quad
s_\lambda'(0)=1,$$
that is, $s_\lambda(r)=\frac{1}{\sqrt\lambda}\sin(\sqrt\lambda\,r)$
for $\lambda>0$, $s_0(r)=r$, and
$s_\lambda(r)=\frac{1}{\sqrt{-\lambda}}\sinh(\sqrt{-\lambda}\,r)$
for $\lambda<0$. Set
$$D_\lambda=\begin{cases}\pi/\sqrt{\lambda}&\lambda>0\\
+\infty&\lambda\le 0,\end{cases}$$
so that $s_\lambda>0$ on $(0,D_\lambda)$.

\begin{lemma}[Jacobi fields and the exponential map in constant curvature]
\label{lem:constant-curvature-jacobi}
\uses{lem:constant-curvature-tensor, def:geodesic, def:exponential-map, thm:hopf-rinow}
\notready

Let $(M,g)$ be a complete Riemannian $n$-manifold of constant
sectional curvature $\lambda$ and let $p\in M$.
\begin{enumerate}
\item Let $\gamma$ be a unit-speed geodesic with $\gamma(0)=p$ and
$X=\gamma'$. The Jacobi fields along $\gamma$ vanishing at $0$ are
exactly the fields
$$Y(t)=a\,t\,X(t)+s_\lambda(t)E(t),$$
where $a\in \mathbb{R}$ and $E$ is a parallel field along $\gamma$
with $E(0)$ perpendicular to $X(0)$; here
$\nabla_XY(0)=aX(0)+E(0)$.
\item $\exp_p$ is defined on all of $T_pM$ and is non-singular at
every $v\in T_pM$ with $|v|<D_\lambda$.
\item On $B(0,D_\lambda)\setminus\{0\}\subset T_pM$, in polar
coordinates $(r,\theta)$ (where $\theta$ denotes coordinates on the
unit sphere of $(T_pM,g_p)$),
$$\exp_p^*g=dr^2+s_\lambda(r)^2\,g_{S^{n-1}},$$
where $g_{S^{n-1}}$ is the round metric on the unit sphere of
$(T_pM,g_p)$; moreover $\exp_p^*g$ is a smooth metric on all of
$B(0,D_\lambda)$.
\end{enumerate}
\end{lemma}

\begin{proof}
\sketch
\uses{lem:constant-curvature-tensor,claim:curvature-symmetries-bianchi,thm:hopf-rinow}
By completeness and the Hopf--Rinow Theorem
(Theorem~\ref{thm:hopf-rinow}), $\exp_p$ is defined on all of
$T_pM$.

(1) By Lemma~\ref{lem:constant-curvature-tensor}, for any vector
field $Y$ along the unit-speed geodesic $\gamma$,
$$\langle{\mathcal R}(Y,X)X,W\rangle={\mathcal R}(Y,X,W,X)
=\lambda\left(\langle Y,W\rangle\langle X,X\rangle
-\langle Y,X\rangle\langle X,W\rangle\right)$$
for every $W$, so ${\mathcal R}(Y,X)X=\lambda\left(Y-\langle
Y,X\rangle X\right)=\lambda Y^\perp$, where $Y^\perp$ denotes the
component of $Y$ perpendicular to $X$. The Jacobi equation
$\nabla_X\nabla_XY+{\mathcal R}(Y,X)X=0$ therefore reads
$\nabla_X\nabla_XY=-\lambda Y^\perp$. Write $u(t)=\langle
Y,X\rangle$; since $X$ is parallel along $\gamma$,
$u''=\langle\nabla_X\nabla_XY,X\rangle=-\lambda\langle
Y^\perp,X\rangle=0$, so $u(t)=u(0)+u'(0)t$. Choosing a parallel
orthonormal frame $E_1=X,E_2,\ldots,E_n$ along $\gamma$ and writing
$Y=uX+\sum_{i\ge 2}y_iE_i$, the perpendicular components satisfy
$y_i''+\lambda y_i=0$. For a Jacobi field with $Y(0)=0$ we get
$u(t)=at$ with $a=\langle \nabla_XY(0),X(0)\rangle$ and
$y_i(t)=y_i'(0)s_\lambda(t)$, which is the asserted form with
$E=\sum_{i\ge2}y_i'(0)E_i$. Conversely each such field solves the
Jacobi equation and vanishes at $0$.

(2) As recalled in the discussion of the exponential mapping above,
for $v,w\in T_pM$ the differential of $\exp_p$ at $v$
applied to $w$ equals $Y_w(1)$, where $Y_w$ is the Jacobi field
along the geodesic $t\mapsto\exp_p(tv)$ with $Y_w(0)=0$ and
$\nabla Y_w(0)=w$. Fix $v$ with $0<r:=|v|<D_\lambda$ and let
$\gamma(s)=\exp_p(s\,v/r)$ be the unit-speed reparameterization. If
$\widetilde Y$ is a Jacobi field along $\gamma$ then
$t\mapsto\widetilde Y(rt)$ is a Jacobi field along
$t\mapsto\exp_p(tv)$, and every Jacobi field arises this way; under
this correspondence $\nabla Y(0)=r\,\nabla\widetilde Y(0)$.
Writing $w=cv/r+w^\perp$ with $w^\perp\perp v$, part (1) gives
\begin{equation}\label{dexpformula}
d(\exp_p)_v(w)=Y_w(1)=\widetilde Y(r)
=c\,\gamma'(r)+\frac{s_\lambda(r)}{r}\,P_\gamma(w^\perp),
\end{equation}
where $P_\gamma$ denotes parallel transport along $\gamma$ from $p$
to $\gamma(r)$ (extended linearly). Since $P_\gamma(w^\perp)$ is
perpendicular to $\gamma'(r)$ (parallel transport preserves inner
products and $w^\perp\perp\gamma'(0)$), the right side vanishes only
if $c=0$ and $s_\lambda(r)w^\perp=0$; as $s_\lambda(r)>0$ for
$0<r<D_\lambda$, the differential is injective. At $v=0$,
$d(\exp_p)_0=\mathrm{id}$.

(3) Fix a unit vector $v$ and $0<r<D_\lambda$. In polar coordinates,
the coordinate field $\partial_r$ at the point $rv$ corresponds to
the vector $v\in T_pM$, and the tangent space to the sphere of
radius $r$ at $rv$ consists of the vectors $w\perp v$. By
Equation~(\ref{dexpformula}) (with $c=1,w^\perp=0$, respectively
$c=0$),
$$d(\exp_p)_{rv}(v)=\gamma'(r),\qquad
d(\exp_p)_{rv}(w)=\frac{s_\lambda(r)}{r}P_\gamma(w)\quad (w\perp
v).$$
Hence, using that parallel transport is a linear isometry and
$\gamma'(r)\perp P_\gamma(w)$:
$$(\exp_p^*g)(v,v)=|\gamma'(r)|^2=1,\qquad
(\exp_p^*g)(v,w)=0,\qquad
(\exp_p^*g)(w_1,w_2)
=\frac{s_\lambda(r)^2}{r^2}\langle w_1,w_2\rangle$$
for $w,w_1,w_2\perp v$. A tangent vector to the unit sphere
$\theta\mapsto$ direction, scaled to the sphere of radius $r$, has
Euclidean length $r$ times its unit-sphere length; the displayed
formulas say exactly that
$\exp_p^*g=dr^2+s_\lambda(r)^2g_{S^{n-1}}$. Finally $\exp_p^*g$ is
smooth on $B(0,D_\lambda)$, being the pullback of a smooth metric
under a smooth map, and it is a metric (positive definite) there
because $d\exp_p$ is injective at each point of $B(0,D_\lambda)$ by
part (2).
\end{proof}

The following consequence packages the case $\lambda\le0$ of the
preceding lemma for repeated use.

\begin{lemma}[Polar model of nonpositive constant curvature]
\label{lem:model-polar-isometry}
\uses{def:sectional-curvature, def:exponential-map, def:local-isometry, lem:constant-curvature-jacobi}
\notready

Let $(M,g)$ be a complete, simply-connected Riemannian $n$-manifold
of constant sectional curvature $\lambda\le0$ and let $p\in M$.
Then $G:=\exp_p^*g$ is a smooth metric on $T_pM$, given in
polar coordinates on $T_pM\setminus\{0\}$ by
$$G=dr^2+s_\lambda(r)^2\,g_{S^{n-1}},$$
and $\exp_p\colon(T_pM,G)\to(M,g)$ is an isometry.
\end{lemma}

\begin{proof}
\sketch
\uses{lem:constant-curvature-jacobi,lem:local-isometry-covering,def:local-isometry}
Constant-curvature Jacobi fields give the polar formula for $\exp_p^*g$ and show that $\exp_p$ is nonsingular for all radii when $\lambda\le0$. The same formula makes the tangent-space polar metric complete. The exponential map is therefore a local isometry from a complete connected source, hence a covering; simple connectivity of the target makes it a global isometry.
\end{proof}

We can now verify the properties of the three model spaces asserted
after Definition~\ref{def:sectional-curvature}.

\begin{lemma}[The model spaces]
\label{lem:model-spaces}
\uses{def:sectional-curvature, def:geodesic, lem:constant-curvature-tensor}
\notready

Let $n\ge2$.
\begin{enumerate}
\item $(\Ar^n,g_{\mathrm{euc}})$ is complete, simply connected, of
constant sectional curvature $0$.
\item For $\rho>0$ the sphere $S^n_\rho\subset\Ar^{n+1}$ of radius
$\rho$, with the induced metric $\rho^2g_{\mathrm{st}}$, is
complete, simply connected, of constant sectional curvature
$1/\rho^2$.
\item The upper half-space $\{x\in\Ar^n\,|\,x^n>0\}$ with the
metric $g_{ij}=(x^n)^{-2}\delta_{ij}$ is complete, simply
connected, of constant sectional curvature $-1$. (This is the
upper half-space presentation of $(\HH^n,g_{\mathrm{st}})$
recalled after Definition~\ref{def:sectional-curvature}; the
Poincar\'e disk presentation is carried to it by the classical
inversion map, as verified at the end of the proof.)
\end{enumerate}
\end{lemma}

\begin{proof}
\sketch
\uses{lem:constant-curvature-tensor,prop:cone-curvature,def:cone,def:curvature-operator,thm:levi-civita-connection,claim:curvature-symmetries-bianchi,def:riemann-curvature-tensor}
We use throughout that multiplying a metric by a constant $c>0$
leaves the Christoffel symbols (Equation~(\ref{Gamma})), hence the
$(1,3)$ curvature tensor, unchanged, multiplies the $(0,4)$-tensor
by $c$, and multiplies each sectional curvature by $1/c$
(evaluate on a rescaled orthonormal basis); it also multiplies
distances by $\sqrt c$, so preserves completeness.

(1) In Cartesian coordinates $g_{ij}=\delta_{ij}$, so
$\Gamma^k_{ij}=0$ and $R_{ijkl}=0$; hence every sectional
curvature vanishes. Geodesics are straight lines, defined for all
time; $(\Ar^n,d_{\mathrm{euc}})$ is a complete metric space, and
$\Ar^n$ is convex, hence simply connected.

(2) By the scaling rule it suffices to treat $\rho=1$. In polar
coordinates the Euclidean metric on $\Ar^{n+1}\setminus\{0\}$ is
$dr^2+r^2g_{S^n}$, which is exactly the cone metric
(Definition~\ref{def:cone}) over $(S^n,g_{\mathrm{st}})$. By (1)
its curvature operator vanishes, so
Proposition~\ref{prop:cone-curvature} (applied with $N=S^n$) gives
$$0=s^2\left(\Rm_{g_{\mathrm{st}}}-\wedge^2g_{\mathrm{st}}\right)
\qquad\text{on }\wedge^2T_qS^n,$$
i.e., $\Rm_{g_{\mathrm{st}}}=\wedge^2g_{\mathrm{st}}$ as bilinear
forms on $\wedge^2TS^n$. A $(0,4)$-tensor antisymmetric in its
first two and in its last two indices (as the curvature tensor is,
by Claim~\ref{claim:curvature-symmetries-bianchi}) is determined by
the bilinear form it induces on $\wedge^2$, so
$R_{ijkl}=g_{ik}g_{jl}-g_{il}g_{jk}$ on $S^n$, and
Lemma~\ref{lem:constant-curvature-tensor} gives constant sectional
curvature $1$. The sphere is compact, hence complete as a metric
space, and it is simply connected for $n\ge2$ (classical
topology).

(3) Write $g=e^{2\phi}\delta$ with $\phi=-\log x^n$, so
$\partial_i\phi=-\delta_{in}/x^n$. From
Equation~(\ref{Gamma}),
$$\Gamma^k_{ij}=\delta_{kj}\,\partial_i\phi
+\delta_{ik}\,\partial_j\phi-\delta_{ij}\,\partial_k\phi
=\frac{-\delta_{kj}\delta_{in}-\delta_{ik}\delta_{jn}
+\delta_{ij}\delta_{kn}}{x^n},$$
so that
$$\nabla_{\partial_i}\partial_j=\frac{-\delta_{jn}\partial_i
-\delta_{in}\partial_j+\delta_{ij}\partial_n}{x^n},\qquad
\nabla_{\partial_i}\partial_n=-\frac{\partial_i}{x^n}.$$
A direct computation of ${\mathcal
R}(\partial_i,\partial_j)\partial_k
=\nabla_{\partial_i}\nabla_{\partial_j}\partial_k
-\nabla_{\partial_j}\nabla_{\partial_i}\partial_k$ from these
formulas (expanding
$\nabla_{\partial_i}\left(\nabla_{\partial_j}\partial_k\right)$ by
the Leibniz rule, with $\partial_i(1/x^n)=-\delta_{in}/(x^n)^2$;
all terms carrying two indices equal to $n$ cancel in the
antisymmetrization) leaves
$${\mathcal R}(\partial_i,\partial_j)\partial_k
=\frac{\delta_{ik}\,\partial_j-\delta_{jk}\,\partial_i}{(x^n)^2}.$$
Hence, using $g_{ab}=\delta_{ab}/(x^n)^2$,
$$R_{ijkl}=g\left({\mathcal
R}(\partial_i,\partial_j)\partial_l,\partial_k\right)
=\frac{\delta_{il}\delta_{jk}-\delta_{jl}\delta_{ik}}{(x^n)^4}
=-\left(g_{ik}g_{jl}-g_{il}g_{jk}\right),$$
and Lemma~\ref{lem:constant-curvature-tensor} gives constant
sectional curvature $-1$. For completeness: for any piecewise
smooth curve $c$,
$$L_g(c)=\int\frac{|c'|_{\mathrm{euc}}}{x^n(c)}\,dt
\ \ge\ \int\left|\frac{d}{dt}\log x^n(c(t))\right|dt
\ \ge\ \left|\log\frac{x^n(c(1))}{x^n(c(0))}\right|,$$
so points of a $d_g$-bounded set have $x^n$ bounded above and away
from $0$. Let $(x_k)$ be a $d_g$-Cauchy sequence; it lies in a
slab $S_{a,b}=\{a\le x^n\le b\}$ with $0<a\le b<\infty$. Applying
the displayed estimate to initial segments shows that every
piecewise smooth curve of $g$-length $<1$ starting in $S_{a,b}$
stays in the enlarged slab $S_{a/e,\,be}$, on which
$(be)^{-2}\delta\le g\le(a/e)^{-2}\delta$. Hence, for $k,l$ large,
curves from $x_k$ to $x_l$ of $g$-length
$<\min\left(2\,d_g(x_k,x_l),1\right)$ stay in $S_{a/e,\,be}$ and
give $|x_k-x_l|_{\mathrm{euc}}\le 2be\,d_g(x_k,x_l)$: the sequence
is Euclidean-Cauchy, so it converges to some $x_\infty$ in the
closed slab $S_{a,b}$. The straight segments from $x_k$ to
$x_\infty$ lie in the convex slab $S_{a,b}$, where $g\le
a^{-2}\delta$, so $d_g(x_k,x_\infty)\le
a^{-1}|x_k-x_\infty|_{\mathrm{euc}}\to0$. Thus $(M,g)$ is
complete. The half-space is convex, hence simply
connected.

Finally we verify that the Poincar\'e disk presentation is carried
to the upper half-space presentation by the classical inversion.
Let
$$\iota(x)=\frac{2\,(x+e_n)}{|x+e_n|^2}-e_n,$$
the inversion in the sphere of radius $\sqrt2$ centered at $-e_n$;
it is an involution of $\Ar^n\setminus\{-e_n\}$. From the
elementary identity
$\left|\frac{a}{|a|^2}-\frac{b}{|b|^2}\right|=\frac{|a-b|}{|a|\,|b|}$
(expand both sides), applied with $a=x+e_n$, $b=x'+e_n$,
\begin{equation}\label{invconf}
|\iota(x)-\iota(x')|
=\frac{2\,|x-x'|}{|x+e_n|\,|x'+e_n|}.
\end{equation}
Direct computation of the $n$-th coordinate gives, using
$|x+e_n|^2=|x|^2+2x^n+1$,
$$\iota(x)^n=\frac{2(x^n+1)}{|x+e_n|^2}-1
=\frac{1-|x|^2}{|x+e_n|^2},$$
so $\iota$ maps the open unit disk into the upper half-space; and
since $\iota(e_n)=0$, Equation~(\ref{invconf}) with $x'=e_n$ gives
$|\iota(x)|=|x-e_n|/|x+e_n|$, which is $<1$ exactly when $x^n>0$,
so $\iota$ maps the upper half-space into the disk. Being an
involution, $\iota$ therefore restricts to mutually inverse
diffeomorphisms between the disk and the half-space. By
Equation~(\ref{invconf}), $d\iota_x$ multiplies all lengths by the
factor $2/|x+e_n|^2$, so $\iota^*\delta=
\left(2/|x+e_n|^2\right)^2\delta$ as quadratic forms; hence
$$\iota^*\left(\frac{\delta}{(y^n)^2}\right)
=\frac{\left(\frac{2}{|x+e_n|^2}\right)^2\,\delta}
{\left(\frac{1-|x|^2}{|x+e_n|^2}\right)^2}
=\frac{4\,\delta}{\left(1-|x|^2\right)^2},$$
which is precisely the Poincar\'e metric of the disk presentation.
\end{proof}

\section{Computations in Gaussian normal coordinates}

In this section we compute the metric and the Laplacian (on
functions) in local Gaussian coordinates. A direct computation shows
that in Gaussian normal coordinates on a metric ball about $p\in M$
the metric takes the form
\begin{eqnarray}\label{metricexp}
g_{ij}(x)& = & \delta_{ij} + \frac{1}{3}R_{iklj}x^kx^l +
\frac{1}{6}R_{iklj,s}x^kx^lx^s \\
& & +(\frac{1}{20}R_{iklj,st} +\frac{2}{45}\sum_m
R_{iklm}R_{jstm})x^kx^lx^sx^t + O(r^{5}), \nonumber
\end{eqnarray}
where $r$ is the distance from $p$. (See, for example Proposition
3.1 on page 41 of~\cite{Sakai}, with the understanding that, with
the conventions there, the quantity $R_{ijkl}$ there differs by sign
from ours.) The fourth-order coefficients in
Equation~(\ref{metricexp}) are a classical computation which we do
not reproduce here, citing~\cite{Sakai} for them; the second- and
third-order terms, which are the ones used in the sequel, we now
derive from the Jacobi equation.

\begin{lemma}[Leading terms of the metric in normal coordinates]
\label{lem:normal-coordinates-expansion}
\uses{def:exponential-map, def:riemann-curvature-tensor}
\notready

In Gaussian normal coordinates $(x^1,\ldots,x^n)$ centered at $p$,
$$g_{ij}(x)=\delta_{ij}+\frac13R_{iklj}x^kx^l
+\frac16R_{iklj,s}x^kx^lx^s+O(|x|^4),$$
where the curvature tensor and its covariant derivative are
evaluated at $p$ in the orthonormal frame defining the coordinates.
\end{lemma}

\begin{proof}
\sketch
\uses{lem:constant-curvature-jacobi,claim:curvature-symmetries-bianchi,def:geodesic}
Radial geodesics are straight coordinate rays, and the Gauss lemma gives the exact radial identities for the metric. The Christoffel symbols vanish at the center, so first derivatives of the metric vanish there. Differentiating the Jacobi equation for coordinate variations gives the quadratic curvature term, yielding the stated Taylor expansion and the corresponding inverse-metric and volume-density expansions.
\end{proof}

Let $\gamma$ be a geodesic in $M$ emanating from $p$ in the direction $v$.
Choose local coordinates $\theta^1,\ldots,\theta^{n-1}$ on the unit sphere in
$T_pM$ in a neighborhood of $v/|v|$. Then $(r, \theta^{1}, ... ,\theta^{n-1})$
are local coordinates at any point of the ray emanating from the origin in the
$v$ direction (except at $p$). Transferring these via $\exp_p$ produces
local coordinates $(r,\theta^{1}, ... ,\theta^{n-1})$ along $\gamma$. Using
Gauss's lemma (Lemma 12 of Chapter 5 on p. 133 of \cite{Petersen}), we can
write the metric locally as
\[ g=dr^{2}+
r^{2}h_{ij}(r,\theta)d\theta^{i}\otimes d\theta^{j}. \] Then the
volume form
\begin{align*}
dV &= \sqrt{\det(g_{ij})}dr\wedge d\theta^{1}
\wedge\cdots\wedge
d\theta^{n-1}\\
& = r^{n-1}\sqrt{\det(h_{ij})}dr\wedge d\theta^{1}
\wedge\cdots\wedge d\theta^{n-1}.
\end{align*}



\begin{lemma}[Laplacian in local coordinates, Christoffel form]
\label{lem:laplacian-christoffel-formula}
\lean{MorganTianLib.laplacianAt_eq_chart_formula,MorganTianLib.laplacianAt_eq_invGram_sum}
\uses{lem:hessian-symmetric, thm:levi-civita-connection}
\leanok



In any coordinate chart, the Laplacian of a smooth function $f$
satisfies
\[ \triangle f
   = g^{ij}\left(\partial_i\partial_jf-\Gamma^k_{ij}\partial_kf\right)
   = g^{ij}\partial_i\partial_jf+b^k\partial_kf,
   \qquad b^k:=-g^{ij}\Gamma^k_{ij}, \]
where $(g^{ij})$ is the inverse of the Gram matrix $(g_{ij})$ of the
metric in the chart and $\Gamma^k_{ij}$ are the Christoffel symbols of
the Levi-Civita connection.
\end{lemma}

\begin{proof}
\sketch
\uses{lem:hessian-symmetric,thm:levi-civita-connection}
\leanok
Trace the coordinate Hessian formula $\Hess(f)_{ij}=\partial_i\partial_jf-\Gamma^k_{ij}\partial_kf$ with $g^{ij}$. Metric compatibility and differentiation of the inverse matrix convert the contracted Christoffel term into the stated derivatives of $g^{ij}$ and $\det g$, giving the divergence-form expression.
\end{proof}

\begin{lemma}[Laplacian in local coordinates]
\label{lem:laplacian-local-formula}
\lean{MorganTianLib.laplacianAt_eq_chart_divergence}
\uses{lem:laplacian-christoffel-formula}
\leanok



The \textit{Laplacian operator} acting on scalar
functions on $M$ is given in local coordinates by
\[ \triangle  = \frac{1}{\sqrt{\det(g)}}\partial_i
\left(g^{ij}\sqrt{\det(g)}\partial_j\right). \]
\end{lemma}

\begin{proof}
\sketch
\uses{lem:hessian-symmetric,lem:laplacian-christoffel-formula}
\leanok
Let us compute the derivative at a point $p$. We have
$$\frac{1}{\sqrt{\det(g)}}\partial_i
\left(g^{ij}\sqrt{\det(g)}\partial_j\right)f=g^{ij}\partial_i\partial_jf+\partial_ig^{ij}\partial_jf+
\frac{1}{2}g^{ij}\partial_iTr(\tilde g)\partial_jf,$$ where $\tilde
g=g(p)^{-1}g$. On the other hand, by the Christoffel form of the
coordinate Laplacian (Lemma~\ref{lem:laplacian-christoffel-formula}),
$$\triangle
f=g^{ij}\partial_i\partial_jf-g^{ij}\Gamma_{ij}^k\partial_kf.$$ Thus,
to prove the claim it suffices to show that
$$g^{ij}\Gamma_{ij}^k=-(\partial_ig^{ik}+\frac{1}{2}g^{ik}\Tr(\partial_i\tilde g)).$$
From the definition of the Christoffel symbols we have
$$g^{ij}\Gamma_{ij}^k=\frac{1}{2}g^{ij}g^{kl}(\partial_ig_{jl}+\partial_jg_{il}-\partial_lg_{ij}).$$
Of course, $g^{ij}\partial_ig_{jl}=-\partial_ig^{ij}g_{jl}$, so that
 $g^{ij}g^{kl}\partial_ig_{jl}=-\partial_ig^{ik}$. It follows by symmetry that $g^{ij}g^{jl}\partial_jg_{il}
 =-\partial_ig^{ik}$. The last term is clearly $-\frac{1}{2}g^{ik}\Tr(\partial_i\tilde g).$
\end{proof}

Using Gaussian local coordinates near $p$, we have
\begin{align*}
\triangle r &=
\frac{1}{r^{n-1}\sqrt{\det(h)}}\partial_r\left(r^{n-1}\sqrt{\det(h)}\right)\\
&= \frac{n-1}{r}
 + \partial_r\log\left(\sqrt{\det(h)}\right).
\end{align*}
From this one computes directly that
 \[ \triangle r
 = \frac{n-1}{r} - \frac{r}{3}\Ric(v,v)+ O(r^{2}), \]
 where $v = \dot r(0)$, cf,  p.265-268 of \cite{Petersen}. So
 $$\triangle r \leq
 \frac{n-1}{r}\ \ \text{when } r \ll 1\ \ \text{and } \Ric > 0.$$
This local computation has a global analogue, in which the
inequality holds across the cut locus in the following weak sense.

\begin{remark}[Laplacian in weak sense]
\label{rem:laplacian-weak-sense}
\lean{MorganTianLib.weakLaplacianLE_iff_test_integral}
\uses{lem:laplacian-local-formula, def:weak-laplacian}
\label{calabi}

Let $f$ be a locally Lipschitz function on $M$ and let $h$ be a
locally integrable function on $M$. The statement that $\triangle
f\leq h$ in {\sl the sense of distributions} (or equivalently {\sl
in the weak sense}) means that for any non-negative test function
$\phi$, that is to say for any non-negative compactly supported
$C^\infty$-function $\phi$, we have
$$\int_Mf\triangle \phi\, d\vol\le \int_M h\,\phi\, d\vol.$$
By Rademacher's theorem a locally Lipschitz function is
differentiable almost everywhere, its almost-everywhere defined
gradient $\nabla f$ is locally bounded, and it coincides with the
distributional gradient; in particular the restriction of $\nabla
f$ to any compact subset of $M$ is an $L^2$ one-form. Moreover, for
every test function $\phi$,
$$\int_Mf\triangle
\phi\, d\vol=-\int_M\langle \nabla f,\nabla\phi\rangle\, d\vol.$$
Indeed, covering the compact support of $\phi$ by finitely many
charts and using a subordinate partition of unity, it suffices to
prove this in a single chart; there, by the divergence form of the
Laplacian (Lemma~\ref{lem:laplacian-local-formula}),
$$\sqrt{\det g}\ \triangle\phi
=\partial_i\left(\sqrt{\det g}\,g^{ij}\,\partial_j\phi\right),$$
and the assertion becomes the integration-by-parts identity for the
locally Lipschitz function $f$ against the compactly supported
smooth vector field with components $\sqrt{\det
g}\,g^{ij}\partial_j\phi$, which is the statement that the
almost-everywhere gradient of $f$ represents its distributional
gradient (standard real analysis).
\end{remark}

\begin{example}[Distance function Laplacian bound]
\label{ex:distance-function-laplacian}
\uses{lem:laplacian-local-formula, def:ricci-curvature, rem:laplacian-weak-sense}

Let $(M,g)$ be a complete connected Riemannian manifold of
dimension $n\ge2$, let $p\in M$, and let $f(x)=d(p,x)$ be the
distance function from $p$. If $(M, g)$ has
$\Ric \geq 0$, then
\[ \triangle f \leq \frac{n-1}{f} \]
 in the sense of distributions (the meaning of which is explained
in Remark~\ref{rem:laplacian-weak-sense} above).
\end{example}

\begin{proof}
  \uses{def:cut-locus, lem:cut-locus-properties, prop:exponential-diffeomorphism-cut-locus, lem:cut-time-star-shaped, lem:volume-element-comparison, lem:geodesic-polar-form, def:ricci-curvature, rem:laplacian-weak-sense, rem:laplacian-weak-green-identity}
\sketch

Use polar coordinates about $p$ off the cut locus. The distance function has unit radial gradient and its Hessian on tangent directions is the second fundamental form of the distance sphere. The assumed mean-curvature bound therefore gives the pointwise Laplacian inequality. Integrate over the domain with a small ball and the cut locus removed, apply the divergence theorem, and let the auxiliary boundaries shrink; their contributions vanish by the local normal-coordinate estimates and nullity of the cut locus, giving the distributional inequality.
\end{proof}
\noindent [Compare \cite{Petersen}, p. 284 Lemma 42].

\medskip

Since $\abs{\nabla f} = 1$ and $\triangle f$ is the mean curvature
of the geodesic sphere $\partial B(x,r)$, $\Ric(v,v)$ measures
the difference of the mean curvature between the standard Euclidean
sphere and the geodesic sphere in the direction $v$. Another
important geometric object  is the shape operator associated to $f$,
denoted $S$. By definition it is the Hessian of $f$; i.e.,
$S=\nabla^2f=\Hess(f)$.

\section{Basic curvature comparison results}
In this section we will recall some of the basic curvature
comparison results in Riemannian geometry. The reader can refer to
\cite{Petersen}, Section 1 of Chapter 9 for details.

\smallskip

We fix a point $p\in M$. For any real number $k\ge 0$ let $H^n_k$ denote the
simply connected, complete Riemannian $n$-manifold of constant sectional
curvature $-k$. Fix a point $q_k\in H^n_k$, and consider the exponential map
$\exp_{q_k}\colon T_{q_k}(H^n_k)\to H^n_k$. This map is a global
diffeomorphism. Let us consider the pullback, $\tilde h_k$, of the Riemannian
metric on $H^n_k$ to $T_{q_k}H^n_k$. A formula for this tensor is easily given
in polar coordinates on $T_{q_k}(H^n_k)$ in terms of the following function.

\begin{definition}[Sinh comparison function]
\label{def:sinh-comparison-function}
\lean{MorganTianLib.snK}
\leanok


We define a function $\sn_k$ as follows:$$\sn_k(r)=\begin{cases}
r \ \  & \text{if }k=0 \\
\frac{1}{\sqrt{k}}\sinh(\sqrt{k}r) \ \ & \text{if } k>0.
\end{cases}$$
\end{definition}

The function $\sn_{k}(r)$ is the solution to the equation
\begin{align*}
 \varphi'' - k\varphi &= 0,\\
\varphi(0) &= 0,\\
\varphi'(0) &= 1.
\end{align*}
We define
 $\ct_{k}(r) = \frac{\sn_{k}^{'}(r)}{\sqrt{k}\sn_{k}(r)}$.

Before stating the comparison theorems we assemble the technical
apparatus on which all of them rest: the structure of the metric in
geodesic polar coordinates, the Riccati equation satisfied by the
shape operator of the geodesic spheres, a scalar Riccati comparison
lemma, and the elementary Sturm comparison bounding the location of
conjugate points from below.

The Gauss Lemma — the assertion that the polar coordinate spheres meet
the radial geodesics orthogonally — rests on a single fact about Jacobi
fields, which we isolate first because it is proved entirely along one
geodesic, with no reference to the exponential map or to polar
coordinates.

\begin{lemma}[Radial Jacobi fields stay orthogonal to the geodesic]
\label{lem:radial-jacobi-orthogonal}
\lean{MorganTianLib.chartCurvature_inner_velocity_eq_zero,MorganTianLib.innerVelocity,MorganTianLib.innerVelocity_eq_chartMetricInner_of_geodesicAt,MorganTianLib.IsJacobiFieldOn.hasDerivAt_chartInner_fst_velocity,MorganTianLib.IsJacobiFieldOn.hasDerivAt_chartInner_snd_velocity,MorganTianLib.IsJacobiFieldAlongOn.innerVelocity_snd_eq,MorganTianLib.IsJacobiFieldAlongOn.innerVelocity_fst_eq,MorganTianLib.IsJacobiFieldAlongOn.innerVelocity_fst_eq_zero,MorganTianLib.IsJacobiFieldAlongOn.innerVelocity_snd_eq_zero}
\uses{def:geodesic, def:conjugate-point, def:riemann-curvature-tensor, claim:curvature-symmetries-bianchi}
\leanok



Let $\gamma$ be a geodesic on $[a,b]$ in a Riemannian manifold $(M,g)$,
and let $J$ be a Jacobi field along $\gamma$, with covariant derivative
$\nabla J$. Then:
\begin{enumerate}
\item $\langle\nabla J,\dot\gamma\rangle$ is constant on $[a,b]$;
\item $\langle J,\dot\gamma\rangle$ is an affine function of $t$:
$$\langle J(t),\dot\gamma(t)\rangle
=\langle J(a),\dot\gamma(a)\rangle
+(t-a)\,\langle\nabla J(a),\dot\gamma(a)\rangle;$$
\item consequently, if $J$ is {\sl radial}, i.e.\ $J(a)=0$ and
$\langle\nabla J(a),\dot\gamma(a)\rangle=0$, then both
$\langle J(t),\dot\gamma(t)\rangle=0$ and
$\langle\nabla J(t),\dot\gamma(t)\rangle=0$ for every $t\in[a,b]$.
\end{enumerate}
\end{lemma}

\begin{proof}
\sketch
\uses{claim:curvature-symmetries-bianchi,thm:levi-civita-connection,def:geodesic}
\leanok
For a Jacobi field $J$ along a geodesic, differentiate $\langle J,\gamma'\rangle$ twice. The geodesic equation removes derivatives of $\gamma'$, and curvature antisymmetry makes the second derivative zero. The initial conditions force both the value and first derivative to vanish, so $J$ remains orthogonal to $\gamma'$.
\end{proof}

A second consequence of the curvature symmetries, also needed below, is
that the operator appearing in the Jacobi equation is self-adjoint.

\begin{lemma}[The Jacobi operator is self-adjoint]
\label{lem:jacobi-operator-symmetric}
\lean{MorganTianLib.curvatureFormAt_jacobi_symm,MorganTianLib.chartCurvature_inner_symm,MorganTianLib.chartCurvatureEndo_inner_symm}
\uses{def:riemann-curvature-tensor, claim:curvature-symmetries-bianchi}
\leanok



Let $u\in T_qM$. The {\sl Jacobi operator} $R_u\colon T_qM\to T_qM$,
$R_u(X)={\mathcal R}(X,u)u$, is self-adjoint:
$$\langle R_u(X),Y\rangle=\langle X,R_u(Y)\rangle
\qquad\text{for all }X,Y\in T_qM.$$
\end{lemma}

\begin{proof}
\sketch
\uses{claim:curvature-symmetries-bianchi}
\leanok
Unwinding the definitions, $\langle R_u(X),Y\rangle={\mathcal
R}(X,u,u,Y)$ and $\langle X,R_u(Y)\rangle={\mathcal R}(Y,u,u,X)$, so
the claim is that ${\mathcal R}(X,u,u,Y)={\mathcal R}(Y,u,u,X)$. By
antisymmetry in the last two arguments
(Claim~\ref{claim:curvature-symmetries-bianchi},
$R_{ijkl}=-R_{ijlk}$),
$${\mathcal R}(X,u,u,Y)=-{\mathcal R}(X,u,Y,u),\qquad
{\mathcal R}(Y,u,u,X)=-{\mathcal R}(Y,u,X,u),$$
and by the pair symmetry $R_{ijkl}=R_{klij}$ of the same claim,
${\mathcal R}(X,u,Y,u)={\mathcal R}(Y,u,X,u)$. Combining the three
identities gives ${\mathcal R}(X,u,u,Y)={\mathcal R}(Y,u,u,X)$.
\end{proof}

\begin{lemma}[Radial geometry in geodesic polar coordinates]
\label{lem:geodesic-polar-form}
\lean{MorganTianLib.expDifferential_gauss,MorganTianLib.expDifferential_gauss_radial_pairing,MorganTianLib.expDifferential_gauss_radial_pairing_fixed_chart,MorganTianLib.geodesic_polar_frame_form}
\uses{def:exponential-map, def:geodesic, lem:hessian-symmetric, def:riemann-curvature-tensor, def:sectional-curvature, lem:radial-jacobi-orthogonal, lem:jacobi-operator-symmetric}

Let $p\in M$, let $V$ be an open set of unit vectors in $T_pM$, let
$0<r_0\le\infty$, and suppose that $\exp_p$ is defined on the cone
$C=\{tw\,|\,w\in V,\ 0\le t<r_0\}$ and is non-singular at every
point of $C$. Let $\tilde g=\exp_p^*g$ on $C$, let
$(r,\theta^1,\ldots,\theta^{n-1})$ be polar coordinates on $C$
(where $\theta$ are local coordinates on $V$), and let $\hat
g=\hat g_{ij}(\theta)d\theta^i\otimes d\theta^j$ denote the round
metric of the unit sphere of $T_pM$ in these coordinates. Then:
\begin{enumerate}
\item (Gauss Lemma) $\tilde g=dr^2+g_{ij}(r,\theta)\,
d\theta^i\otimes d\theta^j$; the coordinate fields
$\partial_{\theta^i}$ restrict along each radial geodesic
$\gamma_w(t)=\exp_p(tw)$ to Jacobi fields vanishing at $t=0$, and
the radial field $\partial_r$ restricts to $\gamma_w'$.
\item Let $f=r$ be the radial coordinate function and
$S=\Hess(f)$. Then the gradient of $f$ is the radial field
$\partial_r$, $S(\partial_r,\cdot)=0$,
$$S_{ij}:=S(\partial_{\theta^i},\partial_{\theta^j})
=\tfrac{1}{2}\partial_rg_{ij},$$
and the associated $g$-symmetric operator $A$ on the tangent space
of the coordinate spheres, defined by $\langle
A(X),Y\rangle=S(X,Y)$, equals $X\mapsto\nabla_X\partial_r$ and
satisfies along each radial geodesic the Riccati equation
$$\nabla_{\partial_r}A+A^2+R_{\partial_r}=0,\qquad
R_{\partial_r}(X):={\mathcal R}(X,\partial_r)\partial_r,$$
where for unit $X\perp\partial_r$ one has $\langle
R_{\partial_r}(X),X\rangle=K(X\wedge\partial_r)$, the sectional
curvature of the plane spanned by $X$ and $\partial_r$.
\item As $r\to0$, locally uniformly in $\theta$,
$$g_{ij}(r,\theta)=r^2\left(\hat g_{ij}(\theta)+O(r^2)\right),
\qquad A(r)=\frac{1}{r}\,\mathrm{Id}+O(r).$$
\item Setting $\lambda(r,\theta)
=\sqrt{\det\left(g_{ij}(r,\theta)\right)/\det\left(\hat
g_{ij}(\theta)\right)}$, the volume form of $\tilde g$ is
$\lambda\, dr\wedge d\mathrm{vol}_{S^{n-1}}$ and
$$\partial_r\log\lambda=\Tr(A),$$
where $\Tr$ is the trace of the operator $A$. Moreover
$\lambda(r,\theta)/r^{n-1}\to 1$ as $r\to 0$, locally uniformly in
$\theta$.
\end{enumerate}
\end{lemma}

\begin{proof}
\sketch
\uses{def:conjugate-point,claim:curvature-symmetries-bianchi,lem:hessian-symmetric,thm:levi-civita-connection,lem:exponential-differential-jacobi,lem:second-order-linear-ode,lem:jacobi-small-time,lem:parallel-frame,lem:jacobi-matrix-inverse,lem:radial-jacobi-orthogonal}
Differentiate the family $r\mapsto\exp_p(ru)$ in angular directions. The resulting fields are Jacobi fields with initial data $J(0)=0$ and $\nabla_rJ(0)$ equal to the angular vector. The Gauss lemma makes them orthogonal to the radial field, so the pullback metric splits into radial and angular parts. A parallel frame identifies the angular differential with the matrix Jacobi solution; its determinant gives the radial volume density, and differentiating that determinant yields the trace of the shape operator.
\end{proof}

\begin{lemma}[Small-time asymptotics of Jacobi-type solutions]
\label{lem:jacobi-small-time}
\lean{MorganTianLib.IsJacobiSolOn.max_norm_le,MorganTianLib.IsJacobiSolOn.norm_snd_le,MorganTianLib.IsJacobiSolOn.norm_fst_le,MorganTianLib.IsJacobiSolOn.norm_snd_sub_le,MorganTianLib.IsJacobiSolOn.norm_fst_sub_le}
\uses{lem:second-order-linear-ode}
\leanok



Let $F$ be a real Banach space and let $R\colon[0,b]\to L(F)$ be
continuous with $\|R(t)\|\le C$ for all $t\in[0,b]$; set
$K=\max(1,C)$. Let $(y,v)$ solve the second-order linear equation
$y''+R(t)\,y=0$ on $[0,b]$ in the sense of
Lemma~\ref{lem:second-order-linear-ode}, with $y(0)=0$ and
$v(0)=v_0$, and put $M=\|v_0\|\,e^{Kb}$. Then for all $t\in[0,b]$:
$$\max\bigl(\|y(t)\|,\|v(t)\|\bigr)\le\|v_0\|\,e^{Kt},\qquad
\|y(t)\|\le M\,t,$$
$$\|v(t)-v_0\|\le\tfrac12\,CM\,t^2,\qquad
\|y(t)-t\,v_0\|\le\tfrac16\,CM\,t^3.$$
More generally, without assuming $y(0)=0$, any solution on $[a,b]$
satisfies the a-priori bound
$\max(\|y(t)\|,\|v(t)\|)\le
\max(\|y(a)\|,\|v(a)\|)\,e^{K(t-a)}$.
\end{lemma}

\begin{proof}
\sketch
\uses{lem:second-order-linear-ode}
\leanok
The a-priori bound is Gr\"onwall's inequality for the pair
$W=(y,v)$: with $A(t)(y,v)=(v,-R(t)\,y)$ as in the proof of
Lemma~\ref{lem:second-order-linear-ode} we have
$\|W'\|=\|A(t)\,W\|\le K\,\|W\|$ on $[a,b]$, whence
$\|W(t)\|\le\|W(a)\|\,e^{K(t-a)}$ in the sup norm
$\|(y,v)\|=\max(\|y\|,\|v\|)$. With $a=0$ and $y(0)=0$ this reads
$\max(\|y(t)\|,\|v(t)\|)\le\|v_0\|\,e^{Kt}$; in particular
$\|v\|\le M$ on $[0,b]$.

By the fundamental theorem of calculus,
$y(t)=\int_0^tv(s)\,ds$, so $\|y(t)\|\le\int_0^t\|v\|\le M\,t$.
Next, $v(t)-v_0=-\int_0^tR(s)\,y(s)\,ds$, so
$$\|v(t)-v_0\|\le\int_0^t\|R(s)\|\,\|y(s)\|\,ds
\le\int_0^tCMs\,ds=\tfrac12\,CM\,t^2.$$
Finally $y(t)-t\,v_0=\int_0^t\bigl(v(s)-v_0\bigr)\,ds$, so
$$\|y(t)-t\,v_0\|\le\int_0^t\tfrac12\,CMs^2\,ds
=\tfrac16\,CM\,t^3.$$
\end{proof}

\begin{lemma}[Inverse asymptotics of the matrix Jacobi solution]
\label{lem:jacobi-matrix-inverse}
\lean{MorganTianLib.IsJacobiSolOn.norm_one_sub_inv_smul_fst_le,MorganTianLib.IsJacobiSolOn.isUnit_fst,MorganTianLib.IsJacobiSolOn.norm_snd_mul_inverse_fst_sub_le}
\uses{lem:second-order-linear-ode, lem:jacobi-small-time}
\leanok



Let $F$ be a unital real Banach algebra and let
$R\colon[0,b]\to L(F)$ be continuous with $\|R\|\le C$ on $[0,b]$;
set $K=\max(1,C)$ and $M=e^{Kb}$. Let $(y,v)$ solve $y''+R(t)\,y=0$
on $[0,b]$ in the sense of Lemma~\ref{lem:second-order-linear-ode},
with $y(0)=0$ and $v(0)=1$ (the unit of $F$). Then for
$t\in(0,b]$:
\begin{enumerate}
\item $\|1-t^{-1}y(t)\|\le\tfrac16\,CM\,t^2$; in particular, if
$CMt^2<6$ then $y(t)$ is invertible in $F$;
\item if $CMt^2\le3$, then
$$\bigl\|v(t)\,y(t)^{-1}-t^{-1}\cdot1\bigr\|\le 2\,CM\,t.$$
\end{enumerate}
Applied to the matrix Jacobi field ${\mathcal J}$ of
Lemma~\ref{lem:geodesic-polar-form}(3) (with $F$ the algebra of
operators on the orthogonal complement of the geodesic's direction),
this is the assertion that ${\mathcal J}(r)$ is invertible for small
$r>0$ with $A(r)={\mathcal J}'(r){\mathcal J}(r)^{-1}
=\tfrac1r\,\mathrm{Id}+O(r)$.
\end{lemma}

\begin{proof}
\sketch
\uses{lem:jacobi-small-time}
\leanok
(1) By Lemma~\ref{lem:jacobi-small-time},
$\|y(t)-t\cdot1\|\le\tfrac16\,CM\,t^3$; dividing by $t>0$ gives the
displayed bound for $z:=t^{-1}y(t)$. If $CMt^2<6$ then
$\|1-z\|<1$, so the geometric series
$\sum_{n\ge0}(1-z)^n$ converges and inverts
$z=1-(1-z)$; hence $z$, and with it $y(t)=t\,z$, is invertible.

(2) Write $w=1-z$, so $\|w\|\le\tfrac16CMt^2\le\tfrac12$ by (1) and
the hypothesis. The geometric series gives
$\|z^{-1}\|\le(1-\|w\|)^{-1}\le2$. By
Lemma~\ref{lem:jacobi-small-time} again,
$\|v(t)-1\|\le\tfrac12CMt^2$, so
$$\|v(t)-z\|\le\|v(t)-1\|+\|1-z\|
\le\tfrac12CMt^2+\tfrac16CMt^2=\tfrac23\,CMt^2.$$
Since $y(t)^{-1}=t^{-1}z^{-1}$ and $1=z\,z^{-1}$,
$$v(t)\,y(t)^{-1}-t^{-1}\cdot1
=t^{-1}\bigl(v(t)-z\bigr)z^{-1},$$
whence
$\|v(t)y(t)^{-1}-t^{-1}\cdot1\|
\le t^{-1}\cdot\tfrac23CMt^2\cdot2=\tfrac43\,CMt\le2\,CMt$.
\end{proof}


\begin{lemma}[No conjugate point $\Rightarrow$ the matrix Jacobi field is invertible]
\label{lem:jacobi-matrix-nonsingular}
\lean{MorganTianLib.isUnit_of_not_isConjugatePointAt,MorganTianLib.exists_isJacobiFieldAlongOn_mem,MorganTianLib.IsJacobiFieldAlongOn.eqOn_zero_of_mem,MorganTianLib.frameVec_frameLift}
\uses{def:conjugate-point, lem:second-order-linear-ode, lem:parallel-frame}
\leanok



Let $\gamma$ be a geodesic, let $E_1,\ldots,E_n$ be a parallel $g$-orthonormal
frame along $\gamma$, and let ${\mathcal J}$ be the matrix Jacobi field in that
frame, i.e.\ the solution of ${\mathcal J}''+{\mathcal R}{\mathcal J}=0$ with
${\mathcal J}(0)=0$, ${\mathcal J}'(0)=\mathrm{Id}$, so that every Jacobi field
$J$ along $\gamma$ with $J(0)=0$ is the column
$J(t)={\mathcal J}(t)\,\nabla J(0)$ (read in the frame). Let $r>0$. If
$\gamma(r)$ is {\em not} conjugate to $\gamma(0)$ along $\gamma$, then
${\mathcal J}(r)$ is invertible.
\end{lemma}

\begin{proof}
\sketch
\uses{def:conjugate-point,lem:second-order-linear-ode,lem:parallel-frame}
\leanok
The frame gives a linear isometry between the coefficient space and
$T_{\gamma(t)}M$, so it suffices to show that ${\mathcal J}(r)$ is injective;
the coefficient space is finite-dimensional, so injectivity gives invertibility.

Let $x$ satisfy ${\mathcal J}(r)\,x=0$ and suppose $x\ne0$. Let $w\in
T_{\gamma(0)}M$ be the vector whose frame coordinates at $t=0$ are $x$; since the
frame is orthonormal, $|w|_g=|x|\ne0$. Let $J$ be the Jacobi field along $\gamma$
with
$$J(0)=0,\qquad \nabla J(0)=w,$$
which exists and is unique by Lemma~\ref{lem:second-order-linear-ode} (existence
and uniqueness for the linear Jacobi equation, with the data prescribed at the
{\em interior} time $0$; this is where the two-sided form of the uniqueness
statement is used).

By the column property, the frame coordinates of $J(t)$ are ${\mathcal
J}(t)\,x$; at $t=r$ these vanish, and the frame is an isometry, so $J(r)=0$.
Moreover $J$ does not vanish identically on $[0,r]$: if it did, then $\nabla
J(0)=0$ (a field vanishing on an interval has vanishing covariant derivative
there), i.e.\ $w=0$, contradicting $|w|=|x|\ne0$.

Thus $J$ is a Jacobi field along $\gamma$, not identically zero, with
$J(0)=J(r)=0$ --- that is, $\gamma(r)$ {\em is} conjugate to $\gamma(0)$ along
$\gamma$ (Definition~\ref{def:conjugate-point}), contrary to hypothesis. Hence
$x=0$, ${\mathcal J}(r)$ is injective, and so invertible.
\end{proof}

\begin{lemma}[The radial shape operator and its Riccati equation]
\label{lem:radial-shape-riccati}
\lean{MorganTianLib.RadialJacobi,MorganTianLib.jacobiOpCoeff,MorganTianLib.shapeOp,MorganTianLib.wronskian_eq,MorganTianLib.shapeOp_symm,MorganTianLib.hasDerivAt_shapeOp,MorganTianLib.tendsto_shapeOp_sub_inv_smul_id,MorganTianLib.IsRadialJacobi.isUnit_fst,MorganTianLib.geodesic_polar_frame_form}
\uses{lem:second-order-linear-ode, lem:jacobi-matrix-inverse, lem:jacobi-small-time}
\leanok



Let $E$ be a real Hilbert space, let $b>0$, and let
$\mathcal R\colon[0,b]\to L(E)$ be continuous with
$\|\mathcal R\|\le C$ and each $\mathcal R(t)$ symmetric. Let
$({\mathcal J},{\mathcal J}')$ solve the matrix Jacobi equation
${\mathcal J}''+{\mathcal R}(t){\mathcal J}=0$ on $[0,b]$ in the sense
of Lemma~\ref{lem:second-order-linear-ode}, with ${\mathcal J}(0)=0$
and ${\mathcal J}'(0)=\mathrm{Id}$. Let $0<r_0\le b$ and suppose
${\mathcal J}(r)$ is invertible for every $r\in(0,r_0)$. Define the
\emph{radial shape operator}
$$A(r)={\mathcal J}'(r)\,{\mathcal J}(r)^{-1},\qquad r\in(0,r_0).$$
Then:
\begin{enumerate}
\item (Wronskian identity) $\langle{\mathcal J}'(t)a,{\mathcal
J}(t)c\rangle=\langle{\mathcal J}(t)a,{\mathcal J}'(t)c\rangle$ for all
$t\in[0,b]$ and $a,c\in E$; consequently $A(r)$ is symmetric for every
$r\in(0,r_0)$.
\item (Riccati equation) $A$ is differentiable on $(0,r_0)$ and
$$A'+A^2+{\mathcal R}=0 .$$
\item (Asymptotics) $A(r)-\tfrac1r\,\mathrm{Id}\to0$ in operator norm
as $r\to0^+$.
\end{enumerate}
\end{lemma}

\begin{proof}
\sketch
\uses{lem:jacobi-matrix-inverse,lem:jacobi-small-time}
\leanok
Let $\mathcal J$ be the matrix Jacobi solution and set $A=\mathcal J'\mathcal J^{-1}$. Differentiating and using $\mathcal J''+R\mathcal J=0$ gives $A'+A^2+R=0$. Constancy of the Jacobi Wronskian makes $A$ self-adjoint. The small-time expansions of $\mathcal J$ and its inverse give $A(r)=r^{-1}\mathrm{Id}+O(r)$.
\end{proof}

\begin{lemma}[Radial comparison in a parallel frame]
\label{lem:radial-comparison}
\lean{MorganTianLib.shapeOp_inner_le,MorganTianLib.norm_jacobi_sq_le,MorganTianLib.trace_shapeOp_le,MorganTianLib.tendsto_snK_div_self,MorganTianLib.tendsto_norm_jacobi_div_self}
\uses{lem:radial-shape-riccati, lem:operator-riccati-upper, lem:trace-riccati-comparison, lem:ratio-monotone, lem:jacobi-small-time, def:sinh-comparison-function}
\leanok



In the setting of Lemma~\ref{lem:radial-shape-riccati}, fix $k\ge0$.
\begin{enumerate}
\item If $\langle{\mathcal R}(r)X,X\rangle\ge-k\,|X|^2$ for all
$r\in(0,r_0)$ and $X\in E$, then
$$\langle A(r)X,X\rangle\ \le\ \frac{\sn_k'(r)}{\sn_k(r)}\,|X|^2 .$$
\item Under the same hypothesis, for every $w\in E$,
$$|{\mathcal J}(r)w|^2\ \le\ \sn_k^2(r)\,|w|^2\qquad(0<r<r_0).$$
\item If instead $\Tr{\mathcal R}(r)\ge-m\,k$ for all $r\in(0,r_0)$,
where $m=\dim E<\infty$, then
$$\Tr A(r)\ \le\ m\,\frac{\sn_k'(r)}{\sn_k(r)}\qquad(0<r<r_0).$$
\end{enumerate}
\end{lemma}

\begin{proof}
\sketch
\uses{lem:radial-shape-riccati,lem:operator-riccati-upper,lem:trace-riccati-comparison,lem:ratio-monotone,lem:jacobi-small-time}
\leanok
The radial shape operator satisfies the matrix Riccati equation and has asymptotic form $r^{-1}\mathrm{Id}+O(r)$. The sectional bound is exactly the corresponding quadratic-form bound on the Jacobi curvature operator. Applying the lower operator Riccati comparison gives the stated model bound throughout the nonsingular interval.
\end{proof}

\begin{lemma}[Derivative of the radial volume element]
\label{lem:radial-volume-element}
\lean{MorganTianLib.volumeElement,MorganTianLib.hasDerivAt_volumeElement}
\uses{lem:radial-shape-riccati}
\leanok



In the setting of Lemma~\ref{lem:radial-shape-riccati}, suppose in addition
that $\dim E=m<\infty$, and set
$$\lambda(r)=\det{\mathcal J}(r).$$
Then $\lambda$ is differentiable on $(0,r_0)$ and
$$\lambda'(r)=\lambda(r)\,\Tr A(r),$$
equivalently $\partial_r\log\lambda=\Tr A$ wherever $\lambda\ne0$.
\end{lemma}

\begin{proof}
\sketch
\uses{lem:radial-shape-riccati}
\leanok
Fix a basis of $E$ and identify endomorphisms with their matrices; the
determinant and the trace of an endomorphism are those of its matrix, and
composition becomes matrix multiplication. The determinant is a polynomial
in the matrix entries, hence differentiable, and its Fréchet derivative is
Jacobi's formula
$$d(\det)_M(B)=\det M\cdot\Tr\bigl(M^{-1}B\bigr)\qquad(M\ \text{invertible}).$$
Since ${\mathcal J}(r)$ is invertible for $r\in(0,r_0)$, the chain rule
applied to $r\mapsto\det{\mathcal J}(r)$ gives
$$\lambda'(r)=\det{\mathcal J}(r)\cdot
\Tr\bigl({\mathcal J}(r)^{-1}{\mathcal J}'(r)\bigr)
=\lambda(r)\,\Tr\bigl({\mathcal J}'(r){\mathcal J}(r)^{-1}\bigr)
=\lambda(r)\,\Tr A(r),$$
where the middle equality is the cyclicity of the trace,
$\Tr(XY)=\Tr(YX)$.
\end{proof}

\begin{lemma}[Scalar Riccati comparison]
\label{lem:scalar-riccati-comparison}
\lean{MorganTianLib.scalar_riccati_comparison,MorganTianLib.scalar_riccati_comparison_Ioi}
\uses{def:sinh-comparison-function}
\leanok



Fix $k\ge 0$ and $0<r_0\le\infty$. Suppose $\phi\colon
(0,r_0)\to\mathbb{R}$ is differentiable with
$$\phi'+\phi^2\le k\quad\text{on }(0,r_0),\qquad\text{and}\qquad
\phi(r)-\frac1r\to0\ \text{ as } r\to0^+.$$
Then
$$\phi(r)\le\frac{\sn_k'(r)}{\sn_k(r)}\qquad\text{for all }r\in
(0,r_0).$$
\end{lemma}

\begin{proof}
\sketch
\uses{def:sinh-comparison-function}
\leanok
Let $a(r)=\sn_k'(r)/\sn_k(r)$. From $\sn_k''=k\,\sn_k$ we get
$a'+a^2=k$, and from the Taylor expansion of $\sn_k$ at $0$,
$a(r)-\frac1r\to0$ as $r\to0^+$. Set
$\psi(r)=\sn_k^2(r)\left(\phi(r)-a(r)\right)$. Then
$\psi(r)\to0$ as $r\to0^+$ (as $\sn_k^2(r)\to 0$ and $\phi-a\to0$),
and
$$\psi'=2\,\sn_k\sn_k'\,(\phi-a)+\sn_k^2(\phi'-a')
\le 2\,\sn_k^2\,a\,(\phi-a)-\sn_k^2(\phi-a)(\phi+a)
=-\sn_k^2\,(\phi-a)^2\le0,$$
using $\phi'-a'\le(k-\phi^2)-(k-a^2)=-(\phi-a)(\phi+a)$. Hence
$\psi\le0$ on $(0,r_0)$, and since $\sn_k^2>0$ there,
$\phi\le a$.
\end{proof}

We also record an elementary companion, used to convert the shape
operator bounds of the comparison theorems into bounds on the metric
coefficients.

\begin{lemma}[Ratio monotonicity under a logarithmic-derivative bound]
\label{lem:ratio-monotone}
\lean{MorganTianLib.antitoneOn_div_sq_of_deriv_le,MorganTianLib.monotoneOn_div_sq_of_le_deriv}
\leanok


Fix $r_0>0$ and let $s,h\colon(0,r_0)\to\mathbb R$ be differentiable
with $s>0$. If $h'\le 2\frac{s'}{s}\,h$ on $(0,r_0)$, then $h/s^2$
is non-increasing there; if instead $h'\ge2\frac{s'}{s}\,h$, then
$h/s^2$ is non-decreasing.
\end{lemma}

\begin{proof}
\sketch
\leanok
By the quotient rule, on $(0,r_0)$,
$$\left(\frac{h}{s^2}\right)'
=\frac{h's^2-2ss'\,h}{s^4}
=\frac{h'-2\frac{s'}{s}h}{s^2},$$
which is $\le0$ (resp.\ $\ge0$) under the first (resp.\ second)
hypothesis; a differentiable function with nonpositive (resp.\
nonnegative) derivative on an interval is non-increasing (resp.\
non-decreasing).
\end{proof}

\begin{lemma}[Scalar Sturm comparison]
\label{lem:scalar-sturm-comparison}
\lean{MorganTianLib.scalar_sturm_comparison,MorganTianLib.scalar_sturm_pos}
\leanok


Fix $K\ge0$ and $t_1>0$ with $\sqrt K\,t_1<\pi$ (no condition when
$K=0$), and write $s_K$ for the solution of $s_K''+K\,s_K=0$,
$s_K(0)=0$, $s_K'(0)=1$ (the function $s_\lambda$ with $\lambda=K$
from the subsection on local isometries and constant curvature).
Let $f\colon[0,t_1]\to\mathbb{R}$ be continuous and twice
differentiable on $(0,t_1)$, and suppose that: $f''\ge -K f$ on
$(0,t_1)$; $f(0)=0$ and $f>0$ on $(0,t_1)$; $f'$ is bounded near
$0^+$; and $f(t)/t\to c$ as $t\to0^+$. Then
$$f(t)\ \ge\ c\,s_K(t)\qquad\text{for all }t\in(0,t_1].$$
In particular, if $c>0$ then $f(t_1)>0$.
\end{lemma}

\begin{proof}
\sketch
\leanok
On $[0,t_1]$ we have $\sqrt K\,t\le\sqrt K\,t_1<\pi$, so $s_K\ge0$
on $[0,t_1]$ and $s_K>0$ on $(0,t_1]$. Consider the Wronskian
$W=f'\,s_K-f\,s_K'$ on $(0,t_1)$. Since $s_K''=-K\,s_K$,
$$W'=f''\,s_K-f\,s_K''=(f''+K f)\,s_K\ \ge\ 0\qquad
\text{on }(0,t_1),$$
so $W$ is non-decreasing. As $t\to0^+$: $f(t)=(f(t)/t)\cdot t\to
c\cdot0=0$; $s_K(t)\to0$ while $f'$ is bounded near $0^+$, so
$f'\,s_K\to0$; and $|s_K'|\le1$, so $f\,s_K'\to0$. Hence
$W(t)\to0$ as $t\to0^+$, and monotonicity gives $W\ge0$ on
$(0,t_1)$. Therefore
$$\left(\frac f{s_K}\right)'=\frac W{s_K^2}\ \ge\ 0
\qquad\text{on }(0,t_1),$$
so $f/s_K$ is non-decreasing on $(0,t_1)$. Since
$s_K(t)/t\to s_K'(0)=1$ as $t\to0^+$,
$$\frac{f(t)}{s_K(t)}=\frac{f(t)}{t}\cdot\frac{t}{s_K(t)}
\longrightarrow c\qquad(t\to0^+),$$
and monotonicity gives $f/s_K\ge c$ on $(0,t_1)$, i.e.
$f\ge c\,s_K$ there; the case $t=t_1$ follows by continuity of $f$
and $s_K$. Finally, if $c>0$ then
$f(t_1)\ge c\,s_K(t_1)>0$, since $0<t_1$ and $\sqrt K\,t_1<\pi$
give $s_K(t_1)>0$.
\end{proof}

The conjugate-point argument below applies the scalar comparison to
the norm $|Y|$ of a Jacobi field. We isolate that step in a
frame-independent, vector-valued form: expressed in a parallel
orthonormal frame along the geodesic, a Jacobi field becomes a curve
$V$ in a fixed inner product space, and the Jacobi equation together
with the curvature bound gives exactly the hypothesis
$\langle V'',V\rangle\ge-K\,|V|^2$ below.

\begin{lemma}[Vector-valued Sturm comparison]
\label{lem:vector-sturm-comparison}
\lean{MorganTianLib.vector_sturm_comparison,MorganTianLib.vector_sturm_ne_zero,MorganTianLib.tendsto_norm_div_self_of_hasDerivWithinAt,MorganTianLib.rauch_lower_vector,MorganTianLib.rauch_lower_vector_ne_zero}
\uses{lem:scalar-sturm-comparison}
\leanok



Fix $K\ge0$ and $T$ with $\sqrt K\,T<\pi$, and let $E$ be a real
inner product space. Let $V\colon[0,T]\to E$ be continuous on
$[0,T]$ and twice differentiable on $(0,T)$, and suppose that:
$\langle V''(t),V(t)\rangle\ge-K\,|V(t)|^2$ on $(0,T)$; $|V'|$ is
bounded near $0^+$; and $|V(t)|/t\to c$ as $t\to0^+$. Then
$$|V(t)|\ \ge\ c\,s_K(t)\qquad\text{for all }t\in(0,T].$$
In particular, if $c>0$ then $V$ has no zero on $(0,T]$; and if
$V(0)=0$ and $V$ is differentiable at $0$ with derivative $V'(0)$,
then the slope hypothesis holds automatically with $c=|V'(0)|$.
\end{lemma}

\begin{proof}
\sketch
\uses{lem:scalar-sturm-comparison}
\leanok
Apply the scalar Sturm argument to $f(t)=\|V(t)\|$, using a smooth regularization at possible zeros. The self-adjoint coefficient bound and Cauchy--Schwarz give the required scalar differential inequality. Initial asymptotics identify the comparison constant, and the scalar result yields the norm estimate and nonvanishing conclusion.
\end{proof}

The passage from a Jacobi field along a curve, measured in the
(time-dependent) inner products of the metric, to a Euclidean curve
eligible for the vector-valued Sturm comparison is a frame computation
that we record once and for all. It packages the parallel orthonormal
frame of Lemma~\ref{lem:parallel-frame} together with the Parseval
expansion and the endpoint slope analysis.

\begin{lemma}[Jacobi fields in a parallel orthonormal frame]
\label{lem:jacobi-frame-reduction}
\lean{MorganTianLib.jacobi_frame_sturm_comparison,MorganTianLib.jacobi_ne_zero_of_curvature_le,MorganTianLib.chartMetricInner_parseval,MorganTianLib.exists_chartMetricInner_orthonormal,MorganTianLib.chartMetricInner_pos}
\uses{lem:parallel-frame, lem:jacobi-field-coordinates, lem:vector-sturm-comparison, def:riemannian-metric}
\leanok



Fix $K\ge0$ and $T>0$ with $\sqrt K\,T<\pi$. Let $u\colon[0,T]\to U$ be a
$C^1$ curve in a coordinate chart $U$ (the Gram matrix of the metric being
differentiable along $u$), and let $(J,\nabla J)$ be a Jacobi field along
$u$ (Lemma~\ref{lem:jacobi-field-coordinates}) with $J(0)=0$. Suppose that
along $u$ the curvature term of the Jacobi equation satisfies
$$\langle{\mathcal R}(J,\dot u)\dot u,\,J\rangle\ \le\ K\,\langle J,J\rangle
\qquad\text{on }(0,T),$$
where $\langle\cdot,\cdot\rangle=\langle\cdot,\cdot\rangle_{u(t)}$ denotes
the inner product defined by the metric at $u(t)$. Then
$$\sqrt{\langle\nabla J(0),\nabla J(0)\rangle}\;s_K(t)\ \le\
\sqrt{\langle J(t),J(t)\rangle}\qquad\text{for all }t\in(0,T].$$
In particular, if $\nabla J(0)\not=0$ then $J$ has no zero on $(0,T]$.
\end{lemma}

\begin{proof}
\sketch
\uses{lem:parallel-frame,lem:jacobi-field-coordinates,lem:vector-sturm-comparison}
\leanok
Write the Jacobi field and its covariant derivative in a parallel orthonormal frame. Parallelity removes all connection terms, so the Jacobi equation becomes the ordinary vector equation $v''+R(t)v=0$, where $R(t)$ is the self-adjoint matrix of the Jacobi operator. Initial values, zeros, norms, and inner products are preserved by the frame isometry, allowing the scalar and vector comparison lemmas to apply componentwise.
\end{proof}

\begin{lemma}[No conjugate points under an upper curvature bound]
\label{lem:conjugate-sturm}
\lean{MorganTianLib.not_isConjugatePointAt_of_sectionalCurvatureAt_le,MorganTianLib.IsJacobiFieldAlongOn.sqrt_metricInner_comparison,MorganTianLib.IsJacobiFieldAlongOn.ne_zero_of_sectionalCurvatureAt_le,MorganTianLib.IsJacobiFieldAlongOn.eqOn_zero,MorganTianLib.scalar_sturm_comparison_extend,MorganTianLib.scalar_sturm_comparison_of_start}
\uses{def:conjugate-point, def:sectional-curvature, def:geodesic}
\leanok



Fix $K\ge 0$, with the convention $\pi/\sqrt K=+\infty$ for $K=0$.
Let $\gamma\colon[0,T]\to M$ be a unit-speed geodesic such that
$K(P)\le K$ for every $2$-plane $P$ tangent to $M$ at a point of
$\gamma$. Then there is no conjugate point along $\gamma$ at any
parameter $t_1$ with $0<t_1<\min(T,\pi/\sqrt K)$.
\end{lemma}

\begin{proof}
\sketch
\leanok
\uses{def:conjugate-point,lem:jacobi-field-coordinates,lem:chart-curvature-coordinates,lem:jacobi-frame-reduction,lem:vector-sturm-comparison}
Express a Jacobi field in a parallel orthonormal frame and apply the vector-valued Sturm comparison to its coefficient vector. The sectional upper bound controls the Jacobi operator by $K$. Comparison with the model solution $s_K$ gives a positive lower bound for the norm before $\pi/\sqrt K$. Hence a nonzero Jacobi field vanishing at the initial point cannot vanish again in that interval, so there are no conjugate points there.
\end{proof}

The comparison theorems below bound the radial geometry from above
under a lower sectional curvature bound. The proof of the
volume--injectivity radius estimate at the end of the chapter also
requires the mirror-image bounds from below, under an upper
curvature bound; we record them here. Recall the function
$s_\lambda$ (with $s_\lambda''+\lambda s_\lambda=0$,
$s_\lambda(0)=0$, $s_\lambda'(0)=1$) and $D_\lambda$ from the
subsection on local isometries and constant curvature.

The analytic core of the mirror bound is the following
manifold-free comparison for operator-valued solutions of a Riccati
differential inequality. Unlike the upper bound of
Lemma~\ref{lem:scalar-riccati-comparison}, it cannot be reduced to
a scalar inequality along a fixed direction (the Cauchy--Schwarz
inequality points the wrong way); instead one works with the
smallest eigenvalue of the operator, which is merely locally
Lipschitz, through one-sided upper support functions and a barrier
argument for its left difference quotients.

\begin{lemma}[Operator Riccati comparison, lower form]
\label{lem:operator-riccati-lower}
\lean{MorganTianLib.operator_riccati_nonneg,MorganTianLib.minRayleigh,MorganTianLib.le_of_left_slope_eventually_lt,MorganTianLib.apply_eq_minRayleigh_smul}
\leanok


Let $E$ be a nontrivial finite-dimensional real inner product space
and fix $r_0>0$. Let $s\colon(0,r_0)\to\mathbb R$ be differentiable
and positive, bounded near $0^+$, and write $a=s'/s$. Let $U(r)$,
$r\in(0,r_0)$, be a differentiable family of symmetric
endomorphisms of $E$ with $U(r)\to0$ as $r\to0^+$, satisfying the
Riccati differential inequality
$$\langle U'(r)X,X\rangle\ \ge\
-\langle U(r)^2X,X\rangle-2a(r)\,\langle U(r)X,X\rangle
\qquad\text{for all }X\in E,\ r\in(0,r_0).$$
Then $U(r)\ge0$ for all $r\in(0,r_0)$.
\end{lemma}

\begin{proof}
\sketch
\leanok
Let $m(r)$ be the least eigenvalue of the self-adjoint solution. At a differentiability point choose a unit eigenvector and compare the derivative of its Rayleigh quotient with the scalar solution $s_k'/s_k$. The Riccati equation and the curvature lower bound give the scalar differential inequality for $m$. The common $1/r$ asymptotics at zero and a first-contact argument prevent $m$ from crossing below the model solution.
\end{proof}

In the opposite direction — an upper bound on the operator under a
lower bound on the curvature term — the Cauchy--Schwarz inequality
points the right way, and the operator bound follows from the scalar
comparison along each fixed direction.

\begin{lemma}[Operator Riccati comparison, upper form]
\label{lem:operator-riccati-upper}
\lean{MorganTianLib.operator_riccati_le}
\uses{lem:scalar-riccati-comparison}
\leanok



Let $E$ be a real inner product space and fix $k\ge0$, $r_0>0$. Let
$A(r)$, $r\in(0,r_0)$, be a differentiable family of symmetric
endomorphisms of $E$ with
$A(r)-\tfrac1r\,\mathrm{Id}\to0$ in operator norm as $r\to0^+$,
satisfying
$$\langle A'(r)X,X\rangle\ \le\ k\,|X|^2-\langle A(r)^2X,X\rangle
\qquad\text{for all }X\in E,\ r\in(0,r_0).$$
Then for all $r\in(0,r_0)$ and $X\in E$,
$$\langle A(r)X,X\rangle\ \le\
\frac{\sn_k'(r)}{\sn_k(r)}\,|X|^2 .$$
\end{lemma}

\begin{proof}
\sketch
\uses{lem:scalar-riccati-comparison}
\leanok
By homogeneity it suffices to consider a unit vector $X$. Set
$\phi(r)=\langle A(r)X,X\rangle$. By symmetry
$\langle A^2X,X\rangle=|AX|^2$, and by the Cauchy--Schwarz
inequality $\phi^2\le|AX|^2$, so the hypothesis gives
$$\phi'\ \le\ k-|AX|^2\ \le\ k-\phi^2 .$$
Moreover
$|\phi(r)-\tfrac1r|=|\langle(A(r)-\tfrac1r\,\mathrm{Id})X,X\rangle|
\le\|A(r)-\tfrac1r\,\mathrm{Id}\|\to0$ as $r\to0^+$. The scalar
Riccati comparison (Lemma~\ref{lem:scalar-riccati-comparison})
yields $\phi(r)\le\sn_k'(r)/\sn_k(r)$ on $(0,r_0)$.
\end{proof}

\begin{lemma}[Rauch lower bounds under an upper curvature bound]
\label{lem:rauch-lower}
\uses{lem:geodesic-polar-form, def:sectional-curvature}

Fix $\bar K\ge0$. In the setting of
Lemma~\ref{lem:geodesic-polar-form} (a cone $C$ of radius $r_0$ on
which $\exp_p$ is defined and non-singular), suppose in addition
that $r_0\le D_{\bar K}=\pi/\sqrt{\bar K}$ and that $K(P)\le\bar K$
for every $2$-plane $P$ tangent at points of $C$'s image. Then for
$0<r<r_0$, writing $s=s_{\bar K}$:
\begin{enumerate}
\item $A(r)\ge\frac{s'(r)}{s(r)}\,\mathrm{Id}$ as symmetric
operators on the tangent space of the coordinate spheres;
\item $g_{ij}(r,\theta)\ge s(r)^2\,\hat g_{ij}(\theta)$ as
bilinear forms;
\item the function $\frac12r^2$ satisfies
$$\Hess\left(\tfrac12r^2\right)\ \ge\
\min\left(1,\ r\,\frac{s'(r)}{s(r)}\right)\tilde g .$$
\end{enumerate}
\end{lemma}

\begin{proof}
\sketch
\uses{lem:geodesic-polar-form,lem:hessian-symmetric,lem:operator-riccati-lower,lem:ratio-monotone}
In a parallel frame, compare the squared norm of the Jacobi solution with the model solution. The upper sectional-curvature bound controls the curvature matrix, and the Wronskian derivative has the required sign. Its zero initial limit and integration give the lower norm estimate. Positivity of the model function before its first zero excludes a second zero of the Jacobi field.
\end{proof}

Before passing to polar coordinates we isolate the comparison estimate
itself, as a statement about Jacobi fields along a single geodesic.  This
is the analytic content of the Sectional Curvature Comparison theorem; the
passage to the metric $g_{ij}(r,\theta)$ afterwards is the Gauss Lemma
(Lemma~\ref{lem:geodesic-polar-form}(1)) applied to the polar coordinate
fields, which are exactly Jacobi fields vanishing at $p$.

\begin{lemma}[Jacobi field comparison along a geodesic]
\label{lem:jacobi-metric-comparison}
\lean{MorganTianLib.exists_isRadialJacobi, MorganTianLib.IsJacobiSolOn.apply, MorganTianLib.continuousOn_frameCurv, MorganTianLib.frameCurvOp_symm, MorganTianLib.isJacobiSolOn_frameVec, MorganTianLib.exists_isRadialJacobi_of_geodesic, MorganTianLib.curvatureFormAt_jacobi_le_of_sectionalCurvatureAt_ge, MorganTianLib.metricInner_jacobi_le_snK, MorganTianLib.metricInner_jacobi_le_snK_of_sectionalCurvature}
\uses{def:geodesic, def:conjugate-point, def:sectional-curvature, def:sinh-comparison-function, lem:parallel-frame, lem:jacobi-frame-reduction, lem:jacobi-operator-symmetric, lem:second-order-linear-ode, lem:radial-comparison}
\leanok



Fix $k\ge 0$. Let $\gamma$ be a geodesic of $(M,g)$ parameterized at unit
speed, defined on an interval containing $[0,r_0]$ in its interior, with
$\gamma(0)=p$, and suppose
$$-k\le K(P)\qquad\text{for every $2$-plane }P\subset T_{\gamma(r)}M,\quad
0<r<r_0.$$
Suppose moreover that $\gamma$ has no conjugate point of $p$ on $(0,r_0)$,
i.e.\ that the matrix Jacobi field $\mathcal J$ of $\gamma$ (the solution of
$\mathcal J''+R\mathcal J=0$ with $\mathcal J(0)=0$, $\mathcal J'(0)=\Id$ in a
parallel orthonormal frame) is invertible on $(0,r_0)$. Then every Jacobi
field $J$ along $\gamma$ with $J(0)=0$ satisfies
$$|J(r)|_g\ \le\ \sn_k(r)\,|\nabla J(0)|_g,\qquad 0<r<r_0.$$
\end{lemma}

\begin{proof}
\sketch
\leanok
\uses{lem:parallel-frame, lem:jacobi-frame-reduction, lem:jacobi-operator-symmetric, lem:second-order-linear-ode, lem:radial-comparison, def:sinh-comparison-function}
Pass to a parallel frame, where the Jacobi field solves a vector-valued second-order equation with self-adjoint curvature coefficient. The sectional bound gives the required quadratic-form bound on that coefficient. Vector Sturm comparison against $s_K$ then controls the norm and prevents an additional zero; translating back through the frame gives the geometric statement.
\end{proof}


Now we can compare manifolds of varying sectional curvature with
those of constant curvature.


\begin{lemma}[Sectional curvature comparison in a parallel frame]
\label{lem:frame-sectional-comparison}
\lean{MorganTianLib.sectional_curvature_comparison}
\uses{lem:jacobi-metric-comparison, lem:radial-comparison, lem:radial-shape-riccati, def:sectional-curvature, def:sinh-comparison-function}
\leanok



Fix $k\ge0$. Let $\gamma\colon[a,b]\to M$ be a unit-speed geodesic, with
$a<0<B<b$ and $r_0\le B$, and suppose that at every point $\gamma(r)$ with
$0<r<r_0$ every sectional curvature satisfies $-k\le K(P)$. Let
$e=(e_0,\ldots,e_{n-1})$ be a parallel $g$-orthonormal frame along $\gamma$
with $e_0=\gamma'$, let $\mathcal J$ be the associated matrix Jacobi field,
and let $A=\mathcal J'\mathcal J^{-1}$ be its shape operator. If $\gamma$ has
no conjugate point of $\gamma(0)$ on $(0,r_0)$ --- that is, $\mathcal J(r)$ is
invertible there --- then for all $0<r<r_0$:
\begin{enumerate}
\item (metric half) every Jacobi field $J$ along $\gamma$ with $J(0)=0$
satisfies $|J(r)|_g^2\ \le\ \sn_k^2(r)\,|\nabla_XJ(0)|_g^2$;
\item (shape-operator half)
$\langle A(r)Y,Y\rangle\ \le\ \dfrac{\sn_k'(r)}{\sn_k(r)}\,|Y|^2$ for every $Y$.
\end{enumerate}
\end{lemma}

\begin{proof}
\sketch
\uses{lem:jacobi-metric-comparison,lem:radial-comparison}
\leanok
The hypothesis $-k\le K(P)$ at $\gamma(r)$ says exactly that the Jacobi operator
$\mathcal R(r)=R(\cdot,\gamma')\gamma'$ satisfies
$\langle\mathcal R(r)X,X\rangle\ge-k\,|X|^2$: for $X$ a unit vector orthogonal to
$\gamma'$ the left side is the sectional curvature of the plane spanned by $X$ and
$\gamma'$, both sides are homogeneous of degree $2$ in $X$, and the component of
$X$ along $\gamma'$ contributes $0$ to both. Transporting through the isometry
$x\mapsto\sum_ix^ie_i$ of the frame turns this into the hypothesis
$-k\|x\|^2\le\langle\mathcal R(r)x,x\rangle$ of
Lemma~\ref{lem:radial-comparison}. Part (2) is then
Lemma~\ref{lem:radial-comparison}(1), and part (1) is
Lemma~\ref{lem:jacobi-metric-comparison}.
\end{proof}

\begin{lemma}[Polar metric comparison: the pullback of $g$ under $\exp_p$]
\label{lem:exp-pullback-metric-comparison}
\lean{MorganTianLib.expDifferential_metricInner_le_of_not_conjugate,MorganTianLib.expDifferential_metricInner_le_of_sectionalCurvature}
\uses{def:exponential-map, lem:exponential-differential-jacobi, lem:frame-sectional-comparison, lem:jacobi-matrix-nonsingular, def:conjugate-point, def:sectional-curvature, def:sinh-comparison-function}
\leanok



Fix $k\ge0$. Let $(M,g)$ be complete, let $p\in M$, let $u\in T_pM$ be a unit
vector, and let $0<r<r_0$. Write $\gamma_u(s)=\exp_p(su)$ for the unit-speed
geodesic in the direction $u$. Suppose that
\begin{enumerate}
\item $\gamma_u$ has no conjugate point of $p$ at any $s\in(0,r_0)$, and
\item at every point $\gamma_u(s)$, $s\in(0,r_0)$, every sectional curvature
satisfies $-k\le K(P)$.
\end{enumerate}
Then for every $Z\in T_pM$,
$$\bigl|d(\exp_p)_{ru}(Z)\bigr|_g^2\ \le\
\Bigl(\frac{\sn_{k}(r)}{r}\Bigr)^{2}\,|Z|_g^2,$$
that is, $\exp_p^{*}g\ \le\ \bigl(\sn_{k}(r)/r\bigr)^{2}\,g_p$ as quadratic
forms at the point $ru\in T_pM$.

Moreover hypothesis (1) may be dropped in favour of a two-sided curvature
bound: if in addition $K(P)\le K_{\mathrm{up}}$ at every $\gamma_u(s)$,
$s\in[0,r_0]$, with $K_{\mathrm{up}}\ge0$ and
$\sqrt{K_{\mathrm{up}}}\,r_0\le\pi$, then the same estimate holds.
\end{lemma}

This is exactly part (1) of Lemma~\ref{lem:geodesic-polar-form} combined with
the metric half of the comparison, written without reference to a chart on the
sphere. Indeed, in polar coordinates the coordinate field
$\partial_{\theta^i}$ at the point $(r,\theta)$ is $d(\exp_p)_{ru}(Z_i)$ for
$Z_i$ the corresponding tangent vector to the \emph{unit} sphere of $T_pM$, and
$\hat g_{ij}(\theta)=\langle Z_i,Z_j\rangle$; since the sphere of radius $r$ is
$r$ times the unit sphere, the displayed estimate reads
$g_{ij}(r,\theta)\le\sn_k^2(r)\,\hat g_{ij}(\theta)$, which is the inequality
of Theorem~\ref{SCC}. The factor $1/r$ is precisely the bookkeeping that
converts a vector measured on the sphere of radius $r$ into one measured in
$T_pM$.

\begin{proof}
\sketch
\uses{lem:exponential-differential-jacobi,lem:frame-sectional-comparison,lem:jacobi-matrix-nonsingular,lem:conjugate-sturm,lem:second-order-linear-ode,def:exponential-map,def:geodesic}
\leanok
Decompose each tangent vector into radial and angular parts. The Gauss lemma treats the radial component and eliminates the cross term. The angular component is the terminal value of a radial Jacobi field. Applying the frame sectional comparison to that field gives the model upper bound for the angular metric, hence the asserted quadratic-form inequality for $\exp_p^*g$.
\end{proof}

\begin{lemma}[Polar volume comparison: the Jacobian of $\exp_p$]
\label{lem:exp-pullback-volume-comparison}
\lean{MorganTianLib.expDifferential_det_le_of_not_conjugate,MorganTianLib.expDifferential_det_le_of_sectionalCurvature}
\uses{def:exponential-map, lem:exponential-differential-jacobi, lem:volume-element-comparison, lem:jacobi-matrix-nonsingular, def:conjugate-point, def:ricci-curvature, def:sinh-comparison-function}
\leanok



Fix $k\ge0$ and let $n=\dim M\ge2$. Let $(M,g)$ be complete, let $p\in M$, let
$u\in T_pM$ be a unit vector, and let $0<r<r_0$. Write
$\gamma_u(s)=\exp_p(su)$. Suppose that
\begin{enumerate}
\item $\gamma_u$ has no conjugate point of $p$ at any $s\in(0,r_0)$, and
\item $\Ric(\gamma_u',\gamma_u')\ge-(n-1)k$ at every $\gamma_u(s)$,
$s\in[0,r_0]$.
\end{enumerate}
Then the differential $d(\exp_p)_{ru}\colon T_pM\to T_{\gamma_u(r)}M$, measured
between the inner products $g_p$ and $g_{\gamma_u(r)}$ (equivalently: read in
$g$-orthonormal bases of source and target), has positive Jacobian, and
$$0\ <\ \det d(\exp_p)_{ru}\ \le\ \Bigl(\frac{\sn_{k}(r)}{r}\Bigr)^{n-1}.$$

Moreover hypothesis (1) may be dropped in favour of an upper curvature bound: if
in addition $K(P)\le K_{\mathrm{up}}$ at every $\gamma_u(s)$, $s\in[0,r_0]$, with
$K_{\mathrm{up}}\ge0$ and $\sqrt{K_{\mathrm{up}}}\,r_0\le\pi$, then the same
estimate holds.
\end{lemma}

This is part (4) of Lemma~\ref{lem:geodesic-polar-form} together with the volume
half of Theorem~\ref{thm:ricci-curvature-comparison}, written without reference
to a chart on the sphere --- the volume analogue of
Lemma~\ref{lem:exp-pullback-metric-comparison}, and obtained by the same move.
Indeed Morgan--Tian's polar volume element is
$\lambda(r,\theta)=\sqrt{\det g(r,\theta)/\det\hat g(\theta)}$, and since the
flat polar Jacobian of $T_pM$ is $r^{n-1}$ one has
$\lambda(r,\theta)=r^{n-1}\det d(\exp_p)_{ru}$; the displayed estimate is
therefore exactly $\lambda\le\sn_k^{n-1}(r)$, which is the inequality
$\sqrt{\det g(r,\theta)}\le\sn_k^{n-1}(r)$ of
Theorem~\ref{thm:ricci-curvature-comparison}.

The exponent is $n-1$, not $n$: by the Gauss lemma the radial direction is
undistorted, $d(\exp_p)_{ru}(u)=\gamma_u'(r)$, a unit vector, so only the $n-1$
directions orthogonal to $u$ are stretched, each by at most $\sn_k(r)/r$.

\begin{proof}
\sketch
\uses{lem:exponential-differential-jacobi,lem:volume-element-comparison,lem:jacobi-matrix-nonsingular,lem:radial-volume-element,def:exponential-map,def:geodesic}
\leanok
Represent the angular differential of $\exp_p$ by the matrix of radial Jacobi fields in a parallel frame. The frame-level sectional comparison bounds the associated positive Gram matrix by the model scalar matrix. Taking determinants gives the Jacobian estimate with exponent $n-1$; the radial direction contributes the separate unit factor from the Gauss lemma.
\end{proof}


\begin{theorem}[Sectional curvature comparison theorem]
\label{thm:sectional-curvature-comparison}
\lean{MorganTianLib.expDifferential_metricInner_le_of_minimizing, MorganTianLib.sectional_curvature_comparison_radial_of_minimizing}
\uses{def:sectional-curvature, def:sinh-comparison-function}
\label{SCC}



(Sectional Curvature Comparison) Fix $k\ge 0$. Let $(M, g)$ be a
Riemannian manifold with the property that $-k \leq K(P)$ for every
$2$-plane $P$ in $TM$. Fix a minimizing geodesic $\gamma\colon
[0,r_0)\to M$ parameterized at unit speed with $\gamma(0)=p$. Impose
Gaussian polar coordinates $(r,\theta^1,\ldots,\theta^{n-1})$ on a
neighborhood of $\gamma$ so that $g=dr^2+g_{ij}\theta^i\otimes
\theta^j$. Then for all $0<r<r_0$, at $\gamma(r)$ we have
$$ (g_{ij}(r,\theta))_{1 \leq
i,j \leq n-1} \leq \sn_{k}^{2}(r),$$ and the shape operator
associated to  the distance function from $p$, $f$, satisfies
$$(S_{ij}(r,\theta))_{1 \leq i,j \leq n-1} \leq \sqrt{k}\ct_{k}(r).$$
\end{theorem}

The two displayed inequalities are to be read as inequalities of
symmetric bilinear forms on the tangent space to the coordinate
sphere at the point $\gamma(r)$: the first says $g_{ij}(r,\theta)w^iw^j\le
\sn_k^2(r)\,\hat g_{ij}(\theta)w^iw^j$ for every $w$, where $\hat g$
is the round metric of the unit sphere of $T_pM$ in the coordinates
$\theta$, and the second says $S_{ij}(r,\theta)w^iw^j\le
\sqrt{k}\,\ct_k(r)\,g_{ij}(r,\theta)w^iw^j$. (Choosing the
coordinates $\theta$ so that at the point under consideration
$\hat g_{ij}=\delta_{ij}$, respectively $g_{ij}=\delta_{ij}$, one
recovers the literal matrix inequalities displayed.) Note that
$\sqrt k\,\ct_k(r)=\sn_k'(r)/\sn_k(r)$.

\begin{proof}
\sketch
\uses{prop:minimal-geodesic-no-conjugate,lem:geodesic-polar-form,lem:radial-comparison,lem:jacobi-metric-comparison,lem:operator-riccati-upper,lem:ratio-monotone,lem:hessian-symmetric,def:cut-locus,lem:cut-locus-properties}
Along each minimizing radial geodesic, the absence of conjugate points makes the Jacobi matrix invertible. The sectional-curvature lower bound controls its Riccati operator by the scalar model solution. Translating the resulting quadratic-form estimate through the differential of $\exp_p$ gives the Hessian comparison for the distance function and the corresponding comparison of the angular part of the polar metric. Minimality supplies the required nonconjugacy up to the cut time.
\end{proof}

The analogous comparisons in the opposite direction, under a
positive upper bound for the sectional curvature, were recorded in
Lemma~\ref{lem:rauch-lower} above. From an upper curvature bound we
shall also need the following local diffeomorphism property of the
exponential mapping.



\begin{lemma}[Exponential map local diffeomorphism for bounded curvature]
\label{lem:local-diffeomorphism-bounded-curvature}
\lean{MorganTianLib.expMapGlobal_isLocalDiffeo_curvatureBall}
\uses{def:curvature-operator-norm, def:exponential-map}
\leanok
\label{localdiffeo}



Fix $K\ge 0$, with the convention $\pi/\sqrt K=+\infty$ for $K=0$.
Let $(M,g)$ be complete, so that (Theorem~\ref{thm:hopf-rinow})
$\exp_p$ is defined on all of $T_pM$, and in particular on the ball
$B(0,\pi/\sqrt{K})$. Assume that $|\Rm(x)|\le K$ (in the sense of
Definition~\ref{def:curvature-operator-norm}) for all $x\in
B(p,\pi/\sqrt{K})$. Then
$\exp_p$ is a local diffeomorphism from $B(0,\pi/\sqrt{K})$ to
the  ball $B(p,\pi/\sqrt{K})$ in
$M$.
\end{lemma}

Completeness is used only to guarantee that every geodesic from $p$
extends to length $\pi/\sqrt K$, i.e.\ that $\exp_p$ is defined on
$B(0,\pi/\sqrt K)$; the proof below uses nothing else about it. So
the lemma holds verbatim under that weaker hypothesis --- as it does,
for instance, when $B(p,\pi/\sqrt K)$ has compact closure.

\begin{proof}
\sketch
\uses{lem:conjugate-sturm, def:curvature-operator-norm, def:exponential-map, def:sectional-curvature, lem:exponential-differential-jacobi, lem:second-order-linear-ode, def:conjugate-point, thm:hopf-rinow}
The curvature-operator bound implies the same upper bound for sectional curvature. Every radial geodesic of length below $\pi/\sqrt K$ has no conjugate point by Sturm comparison, so the differential of $\exp_p$ is nonsingular there. The inverse function theorem gives a local diffeomorphism at every nonzero vector; at the origin this is the standard exponential-map differential identity.
\end{proof}


There is a crucial comparison result for volume which involves the
Ricci curvature. Its algebraic-analytic core is the following traced
form of the Riccati comparison, which converts the matrix Riccati
equation of the shape operator into the scalar comparison via the
Cauchy--Schwarz inequality for traces.

\begin{lemma}[Trace Riccati comparison]
\label{lem:trace-riccati-comparison}
\lean{MorganTianLib.trace_riccati_comparison,MorganTianLib.sq_trace_le_finrank_mul_trace_comp_self}
\uses{def:sinh-comparison-function}
\leanok



Fix $k\ge0$ and $r_0>0$. Let $E$ be a real inner product space of
finite dimension $m\ge1$ and let $A(r)$, $0<r<r_0$, be a
differentiable family of symmetric endomorphisms of $E$ such that
$$\Tr A'(r)+\Tr\!\left(A(r)^2\right)\ \le\ m\,k
\qquad(0<r<r_0),$$
and $\Tr A(r)-m/r\to0$ as $r\to0^+$. Then
$$\Tr A(r)\ \le\ m\,\frac{\sn_k'(r)}{\sn_k(r)}\qquad(0<r<r_0).$$
\end{lemma}

\begin{proof}
\sketch
\uses{lem:scalar-riccati-comparison}
\leanok
First, for any symmetric endomorphism $A$ of $E$,
$$(\Tr A)^2\ \le\ m\,\Tr(A^2):$$
in an orthonormal basis $(e_i)_{i=1}^m$ of $E$, symmetry gives
$\Tr(A^2)=\sum_i\langle A^2e_i,e_i\rangle=\sum_i|Ae_i|^2$, the
Cauchy--Schwarz inequality in $E$ gives
$|Ae_i|^2\ge\langle Ae_i,e_i\rangle^2$ for each $i$ (as $|e_i|=1$),
and the Cauchy--Schwarz inequality in $\mathbb{R}^m$ gives
$\left(\sum_i\langle Ae_i,e_i\rangle\right)^2
\le m\sum_i\langle Ae_i,e_i\rangle^2$; combining,
$(\Tr A)^2\le m\,\Tr(A^2)$. Now set $\phi(r)=\Tr A(r)/m$. Then
$\phi$ is differentiable on $(0,r_0)$ with $\phi'=\Tr A'/m$ (the
trace is linear), and
$$\phi'+\phi^2=\frac{\Tr A'}{m}+\frac{(\Tr A)^2}{m^2}
\ \le\ \frac{\Tr A'+\Tr(A^2)}{m}\ \le\ k,$$
using the trace inequality and the hypothesis. Moreover
$\phi(r)-1/r=(\Tr A(r)-m/r)/m\to0$ as $r\to0^+$. By the scalar
Riccati comparison (Lemma~\ref{lem:scalar-riccati-comparison}),
$\phi(r)\le\sn_k'(r)/\sn_k(r)$ on $(0,r_0)$, which is the claim.
\end{proof}

We can now state the monotonicity of the volume element in geodesic
polar coordinates, which we isolate because it will be used again in
the proofs of the Bishop--Gromov inequality and of the
volume--injectivity radius estimate.

\begin{lemma}[Monotonicity of the volume element]
\label{lem:volume-element-comparison}
\uses{lem:geodesic-polar-form, def:ricci-curvature, def:sinh-comparison-function}

Fix $k\ge0$. Let $p\in M$ and let $\gamma(t)=\exp_p(tv)$ be a
unit-speed radial geodesic, defined and free of conjugate points on
$[0,r_0)$, along which $\Ric\ge-(n-1)k$. Let $\lambda(t,\theta)$ be
the volume element of Lemma~\ref{lem:geodesic-polar-form}(4) in
polar coordinates around $p$ (defined near the direction of
$\gamma$ for $t<r_0$). Then along $\gamma$,
$$\Tr(A)(t)\le(n-1)\frac{\sn_k'(t)}{\sn_k(t)},\qquad
\frac{\lambda(t)}{\sn_k^{n-1}(t)}\ \text{is non-increasing on
}(0,r_0),\qquad \lim_{t\to0^+}\frac{\lambda(t)}{\sn_k^{n-1}(t)}=1.$$
In particular $\lambda(t)\le\sn_k^{n-1}(t)$ for $t\in(0,r_0)$.
\end{lemma}

\begin{proof}
\sketch
\uses{lem:geodesic-polar-form,lem:radial-comparison,lem:radial-volume-element,lem:trace-riccati-comparison,def:ricci-curvature}
Differentiate the logarithm of the radial Jacobian and identify it with the trace of the radial shape operator. Trace Riccati comparison bounds this derivative by the logarithmic derivative of $s_k^{n-1}$. Hence their quotient is nonincreasing. The Jacobi small-time expansion gives limit one and therefore the pointwise upper bound.
\end{proof}

The next lemma is the whole of the preceding one except for the
polar-coordinate reading of $\lambda$: it is stated for the matrix Jacobi
field $\mathcal J$ of the radial geodesic in a parallel orthonormal frame,
which is where the argument actually takes place, and it is what
Theorem~\ref{thm:bishop-gromov} integrates. Here $\mathcal J$ is the
$E$-valued solution of $\mathcal J''+\mathcal R\mathcal J=0$,
$\mathcal J(0)=0$, $\mathcal J'(0)=\mathrm{Id}$ on the \emph{full}
$n$-dimensional frame, $A=\mathcal J'\mathcal J^{-1}$ is the shape operator,
and $u$ is the radial direction (the frame coordinate of $\gamma'$). The
factor $1/r$ removes the radial column $\mathcal J(r)u=r\,u$, which the polar
coordinate $dr$ already accounts for; $\nu=\det\mathcal J/r$ is thus the
honest polar density, and $\Tr A-1/r$ the honest shape operator of the
geodesic sphere.

\begin{lemma}[Volume element comparison in a parallel frame]
\label{lem:frame-volume-comparison}
\lean{MorganTianLib.sq_trace_sub_le,MorganTianLib.trace_riccati_comparison_perp,MorganTianLib.IsRadialJacobi.apply_radial,MorganTianLib.shapeOp_apply_radial,MorganTianLib.trace_shapeOp_le_perp,MorganTianLib.antitoneOn_div_pow_of_deriv_le,MorganTianLib.tendsto_volumeElement_div_pow,MorganTianLib.volumeElement_pos,MorganTianLib.polarDensity,MorganTianLib.hasDerivAt_polarDensity,MorganTianLib.antitoneOn_polarDensity_div_snK_pow,MorganTianLib.tendsto_polarDensity_div_snK_pow,MorganTianLib.polarDensity_le_snK_pow,MorganTianLib.trace_frameCurvOp_eq_ricciAt,MorganTianLib.frameCurvOp_radial_eq_zero,MorganTianLib.exists_isRadialJacobi_of_geodesic_velocity,MorganTianLib.ricci_curvature_comparison,MorganTianLib.volume_element_comparison_frame}
\uses{lem:radial-shape-riccati, lem:radial-volume-element, lem:trace-riccati-comparison, lem:ratio-monotone, def:ricci-curvature, def:sinh-comparison-function}
\leanok



Fix $k\ge0$ and let $\mathcal J$ be the matrix Jacobi field of a unit-speed
geodesic $\gamma$ on a parallel orthonormal frame of dimension $n\ge2$ whose
$0$-th vector is $\gamma'$, so that the radial direction $u$ satisfies
$\mathcal R(t)u=0$. Suppose $\mathcal J(r)$ is invertible for $r\in(0,r_0)$
($\gamma$ has no conjugate point there) and $\Tr\mathcal R(r)\ge-(n-1)k$
(equivalently, by $\Tr\mathcal R=\Ric(\gamma',\gamma')$, that
$\Ric\ge-(n-1)k$ along $\gamma$). Write $\nu(r)=\det\mathcal J(r)/r$. Then,
for $0<r<r_0$:
\begin{enumerate}
\item $\mathcal J(r)u=r\,u$ and $A(r)u=\tfrac1r\,u$;
\item $\Tr A(r)-\dfrac1r\ \le\ (n-1)\dfrac{\sn_k'(r)}{\sn_k(r)}$;
\item $\nu(r)/\sn_k^{n-1}(r)$ is non-increasing on $(0,r_0)$ and tends to $1$
as $r\to0^+$;
\item $\nu(r)\ \le\ \sn_k^{n-1}(r)$.
\end{enumerate}
\end{lemma}

\begin{proof}
\sketch
\uses{lem:radial-shape-riccati,lem:radial-volume-element,lem:trace-riccati-comparison,lem:ratio-monotone}
\leanok
The logarithmic derivative of the radial Jacobian is the trace of the Riccati operator. Trace Riccati comparison bounds it by $(n-1)s_k'/s_k$. The ratio-monotonicity lemma then shows that the Jacobian divided by $s_k^{n-1}$ is nonincreasing, and its small-time Jacobi asymptotics normalize the ratio to one at the origin.
\end{proof}

\begin{theorem}[Ricci curvature comparison theorem]
\label{thm:ricci-curvature-comparison}
\lean{MorganTianLib.ricci_curvature_comparison, MorganTianLib.expDifferential_det_le_of_minimizing, MorganTianLib.ricci_curvature_comparison_radial_of_minimizing}
\uses{def:ricci-curvature, def:sinh-comparison-function}
\label{riccurvcomp}



{\bf (Ricci curvature comparison)} Fix $k\ge 0$. Assume that $(M,
g)$ satisfies $\Ric \geq -(n-1)k$. Let $\gamma\colon [0,r_0)\to
M$ be a minimal geodesic of unit speed. Then for any $r<r_0$ at
$\gamma(r)$ we have
$$\sqrt{\det\, g(r,\theta)} \leq \sn^{n-1}_{k}(r)$$ and
$$\Tr(S)(r,\theta) \leq (n-1)\frac{\sn_{k}^{'}(r)}{\sn_{k}(r)}.$$
\end{theorem}

As in Theorem~\ref{thm:sectional-curvature-comparison}, the first
inequality is normalized against the round metric: it reads
$\sqrt{\det g(r,\theta)/\det\hat g(\theta)}\le\sn_k^{n-1}(r)$, and
takes the displayed form when the coordinates $\theta$ are chosen
so that $\hat g_{ij}=\delta_{ij}$ at the point in question.

\begin{proof}
\sketch
\uses{prop:minimal-geodesic-no-conjugate,lem:volume-element-comparison,lem:geodesic-polar-form}
Since $\gamma$ is minimal on $[0,r_0)$, there are no conjugate
points along $\gamma|_{[0,r_0)}$
(Proposition~\ref{prop:minimal-geodesic-no-conjugate}, applied as in
the proof of Theorem~\ref{thm:sectional-curvature-comparison}), so
Lemma~\ref{lem:volume-element-comparison} applies to $\gamma$. Its
first assertion is precisely
$$\Tr(S)(r,\theta)=\Tr(A)(r)\le(n-1)\frac{\sn_k'(r)}{\sn_k(r)},$$
and its final assertion is
$\lambda(r)=\sqrt{\det g(r,\theta)/\det\hat g(\theta)}
\le\sn_k^{n-1}(r)$, which is the first displayed inequality in the
normalization explained above.
\end{proof}

Note that the pointwise inequality $\triangle f\le(n-1)/f$ away
from the cut locus, which underlies
Example~\ref{ex:distance-function-laplacian}, follows from this
theorem with $k=0$.

The comparison result in Theorem~\ref{thm:ricci-curvature-comparison} holds out to
every radius, a fact that will be used repeatedly in our arguments.
This result evolved over the period 1964-1980 and now is referred to
as the Bishop-Gromov inequality; see
Proposition 4.1 of \cite{CheegerGromovTaylor}.

Since the Bishop--Gromov inequality is stated for a ball with
compact closure in a manifold that need not be complete, we first
record that the cut-locus apparatus localizes to such a ball.

\begin{lemma}[Localized cut-locus apparatus]
\label{lem:localized-cut-locus}
\uses{def:exponential-map, def:geodesic, def:conjugate-point}

Let $(M,g)$ be a Riemannian manifold, not assumed complete, let
$p\in M$ and $R>0$, and suppose that $B(p,R)$ has compact closure
in $M$. Then:
\begin{enumerate}
\item every unit-speed geodesic from $p$ extends to every length
less than $R$; in particular
$B(0,R)\subset T_pM$ lies in the maximal domain of $\exp_p$;
\item every $q\in B(p,R)$ is joined to $p$ by a minimizing geodesic
of length $d(p,q)$;
\item the set $U_p^{<R}$, consisting of $0$ together with all $v\in
T_pM$ with $0<|v|<R$ such that $\gamma_v$ is the unique minimizing
geodesic from $p$ to $\exp_p(v)$ and $\exp_p$ is a local
diffeomorphism at $v$, is open and star-shaped about $0$: for each
unit vector $v$ there is $c^R(v)\in(0,R]$ with $\{0<t<R\,|\,tv\in
U_p^{<R}\}=(0,c^R(v))$; moreover $d(p,\exp_p(w))=|w|$ for $w\in
U_p^{<R}$, and along each ray the geodesic $\gamma_v$ is minimizing
and free of conjugate points on $[0,c^R(v))$;
\item for every $r<R$ the set $B(p,r)\setminus
\exp_p\left(U_p^{<R}\cap B(0,r)\right)$ has measure zero, and
$\exp_p$ restricted to $U_p^{<R}\cap B(0,r)$ is a
diffeomorphism onto its image in $B(p,r)$.
\end{enumerate}
\end{lemma}

\begin{proof}
\sketch
\uses{prop:minimal-geodesic-no-conjugate,def:conjugate-point,lem:geodesic-polar-form,lem:normal-ball,lem:exponential-differential-jacobi,lem:minimizing-geodesic-exists}
Restrict the cut-time and exponential-map construction to the compact set of directions under consideration. The star-shaped segment domain gives injectivity before the cut time, while upper semicontinuity makes the radial boundary measurable. Local inverse-function arguments away from conjugate points show that the exponential map is a diffeomorphism onto the complement of the localized cut locus. The radial graph is null and smooth maps preserve null sets, so the exceptional image is null.
\end{proof}

\begin{lemma}[Monotonicity of a ratio of integrals]
\label{lem:integral-ratio-monotone}
\lean{MorganTianLib.antitoneOn_integral_ratio}
\leanok


Let $R>0$ and let $f,g\colon(0,R)\to\mathbb R$ be integrable on $(0,r)$ for
every $r<R$, with $g>0$ on $(0,R)$. If $f/g$ is non-increasing on $(0,R)$,
then
$$r\ \longmapsto\ \frac{\int_0^rf(t)\,dt}{\int_0^rg(t)\,dt}$$
is non-increasing on $(0,R)$.
\end{lemma}

\begin{proof}
\sketch
\leanok
Let $0<r_1\le r_2<R$; we may assume $r_1<r_2$. Put $L=f(r_1)/g(r_1)$,
$\Phi(r)=\int_0^rf$ and $\Psi(r)=\int_0^rg$. Since $f/g$ is non-increasing and
$g>0$, we have $f\ge Lg$ pointwise on $(0,r_1)$ and $f\le Lg$ pointwise on
$(r_1,r_2)$. Integrating the first over $(0,r_1)$ gives $\Phi(r_1)\ge L\Psi(r_1)$,
and integrating the second over $(r_1,r_2)$ gives
$\int_{r_1}^{r_2}f\le L\int_{r_1}^{r_2}g$. Also $\Psi(r_1)>0$ and
$\Psi(r_2)>0$ because $g>0$, and $\int_{r_1}^{r_2}g\ge0$. Therefore
$$\Psi(r_1)\int_{r_1}^{r_2}f\ \le\ \bigl(L\,\Psi(r_1)\bigr)\int_{r_1}^{r_2}g
\ \le\ \Phi(r_1)\int_{r_1}^{r_2}g .$$
Adding $\Phi(r_1)\Psi(r_1)$ to both sides and using the additivity relations
$\Phi(r_2)=\Phi(r_1)+\int_{r_1}^{r_2}f$ and
$\Psi(r_2)=\Psi(r_1)+\int_{r_1}^{r_2}g$, this reads
$\Phi(r_2)\Psi(r_1)\le\Phi(r_1)\Psi(r_2)$, i.e.
$\Phi(r_2)/\Psi(r_2)\le\Phi(r_1)/\Psi(r_1)$, as required. (Note that no
differentiability of $\Phi$ or $\Psi$ is used.)
\end{proof}

\begin{lemma}[Normalisation of a ratio of integrals]
\label{lem:integral-ratio-normalized}
\lean{MorganTianLib.tendsto_integral_ratio}
\leanok


Let $R>0$ and let $f,g\colon(0,R)\to\mathbb R$ be integrable on $(0,r)$ for
every $r<R$, with $g>0$ on $(0,R)$. If the pointwise ratio $f/g$ tends to $1$
as $t\to0^+$, then
$$r\ \longmapsto\ \frac{\int_0^rf(t)\,dt}{\int_0^rg(t)\,dt}\ \longrightarrow\ 1
\qquad\text{as }r\to0^+.$$
\end{lemma}

\begin{proof}
\sketch
\leanok
Fix $\varepsilon>0$. Since $f/g\to1$ at $0^+$ there is $\delta\in(0,R)$ with
$|f(t)/g(t)-1|<\varepsilon/2$ for all $t\in(0,\delta)$; as $g(t)>0$ this reads
$(1-\tfrac\varepsilon2)g(t)\le f(t)\le(1+\tfrac\varepsilon2)g(t)$. For
$r\in(0,\delta)$, integrate over $(0,r)$ and divide by $\int_0^rg>0$ (positive
because $g>0$): the primitive ratio $\int_0^rf/\int_0^rg$ lies in
$[1-\tfrac\varepsilon2,1+\tfrac\varepsilon2]$, hence within $\varepsilon$ of $1$.
As in Lemma~\ref{lem:integral-ratio-monotone}, no differentiability and no
fundamental theorem of calculus is used.
\end{proof}

\begin{lemma}[Bishop-Gromov along a radial direction]
\label{lem:frame-bishop-gromov}
\lean{MorganTianLib.bishop_gromov_radial,MorganTianLib.tendsto_bishop_gromov_radial}
\uses{lem:frame-volume-comparison, lem:integral-ratio-monotone, lem:integral-ratio-normalized, def:ricci-curvature, def:sinh-comparison-function}
\leanok



In the situation of Lemma~\ref{lem:frame-volume-comparison} --- a unit-speed
geodesic $\gamma$ on a parallel orthonormal frame of dimension $n\ge2$ whose
$0$-th vector is $\gamma'$, free of conjugate points on $(0,r_0)$, along which
$\Tr\mathcal R\ge-(n-1)k$ (equivalently $\Ric\ge-(n-1)k$), with $k\ge0$ ---
write $\nu(t)=\det\mathcal J(t)/t$ for the polar volume density. Then the ratio
of the radial integrals of $\nu$ and of the model density,
$$r\ \longmapsto\ \frac{\int_0^r\nu(t)\,dt}{\int_0^r\sn_k^{n-1}(t)\,dt},$$
is non-increasing on $(0,r_0)$, and tends to $1$ as $r\to0^+$.
\end{lemma}

\begin{proof}
\sketch
\uses{lem:frame-volume-comparison,lem:integral-ratio-monotone,lem:integral-ratio-normalized}
\leanok
By Lemma~\ref{lem:frame-volume-comparison}(3) the ratio $\nu/\sn_k^{n-1}$ is
non-increasing on $(0,r_0)$, and $\sn_k^{n-1}>0$ there. Both functions are
integrable on $(0,r)$ for each $r<r_0$: $\sn_k^{n-1}$ is continuous, while $\nu$
is continuous on $(0,r_0)$ (the entries of $\mathcal J$ are differentiable and
$\det$ is continuous) and satisfies $0\le\nu\le\sn_k^{n-1}$ there by
Lemma~\ref{lem:frame-volume-comparison}(4), hence is bounded and measurable on
$(0,r)$. Lemma~\ref{lem:integral-ratio-monotone}, applied with $f=\nu$ and
$g=\sn_k^{n-1}$, gives the monotonicity. For the limit: by
Lemma~\ref{lem:frame-volume-comparison}(3) the pointwise ratio
$\nu(t)/\sn_k^{n-1}(t)\to1$ as $t\to0^+$, so
Lemma~\ref{lem:integral-ratio-normalized}, with the same $f,g$, gives that the
ratio of primitives tends to $1$.
\end{proof}

\begin{definition}[The volume density in a chart]
\label{def:chart-volume-density}
\lean{MorganTianLib.chartVolumeDensity}
\uses{def:riemannian-metric}
\leanok



Let $(M,g)$ be an $n$-dimensional Riemannian manifold without boundary and let
$\varphi_\alpha\colon U_\alpha\to E$ be a chart. Write
$$g^\alpha_{ij}(x)=g_x(\partial^\alpha_i,\partial^\alpha_j),\qquad x\in U_\alpha,$$
for the Gram matrix of the coordinate frame $\partial^\alpha_1,\dots,\partial^\alpha_n$
of the chart. Since $g_x$ is an inner product and the $\partial^\alpha_i$ are a basis of
$T_xM$, the matrix $g^\alpha(x)$ is symmetric positive definite; in particular
$\det g^\alpha(x)>0$. The \emph{volume density of $g$ in the chart $\alpha$} is the
function on chart coordinates
$$\lambda_\alpha(y)=\sqrt{\det g^\alpha\bigl(\varphi_\alpha^{-1}(y)\bigr)},
\qquad y\in\varphi_\alpha(U_\alpha).$$
\end{definition}

\begin{lemma}[Transformation law of the volume density]
\label{lem:volume-density-change}
\lean{MorganTianLib.chartGramMatrix_det_change,MorganTianLib.sqrt_chartGramMatrix_det_change}
\uses{def:chart-volume-density}
\leanok



Let $\alpha,\beta$ be charts and $x\in U_\alpha\cap U_\beta$. Let
$A=D(\varphi_\alpha\circ\varphi_\beta^{-1})$ be the derivative at $\varphi_\beta(x)$ of the
transition map, viewed as the change of coordinate frames
$\partial^\beta_i=\sum_a A^a_{\ i}\,\partial^\alpha_a$. Then
$$g^\beta(x)=A^{T}g^\alpha(x)A,\qquad
\det g^\beta(x)=(\det A)^2\det g^\alpha(x),$$
and consequently
$$\sqrt{\det g^\beta(x)}=\lvert\det A\rvert\cdot\sqrt{\det g^\alpha(x)}.$$
\end{lemma}

\begin{proof}
\sketch
\uses{def:chart-volume-density}
\leanok
The metric is a $(0,2)$-tensor, so bilinearity of $g_x$ gives
$$g^\beta_{ij}(x)=g_x(\partial^\beta_i,\partial^\beta_j)
=g_x\Bigl(\sum_a A^a_{\ i}\partial^\alpha_a,\ \sum_b A^b_{\ j}\partial^\alpha_b\Bigr)
=\sum_{a,b}A^a_{\ i}\,g^\alpha_{ab}(x)\,A^b_{\ j},$$
which is the matrix identity $g^\beta=A^{T}g^\alpha A$. Taking determinants and using
multiplicativity of $\det$ together with $\det A^{T}=\det A$ yields
$\det g^\beta=(\det A)^2\det g^\alpha$. Both determinants are positive, so taking square
roots and using $\sqrt{(\det A)^2}=\lvert\det A\rvert$ gives the last identity.
\end{proof}

\begin{lemma}[The chart volume is chart-independent]
\label{lem:riemannian-measure-well-defined}
\lean{MorganTianLib.chartMeasure_apply_eq}
\uses{def:chart-volume-density, lem:volume-density-change}
\leanok



For a chart $\alpha$ and a measurable $S\subseteq U_\alpha$ set
$$\mu_\alpha(S)=\int_{\varphi_\alpha(S)}\lambda_\alpha\,d\mu ,$$
where $\mu$ is Lebesgue (Haar) measure on the model space. If $S$ is a measurable set with
$S\subseteq U_\alpha\cap U_\beta$, then $\mu_\alpha(S)=\mu_\beta(S)$.
\end{lemma}

\begin{proof}
\sketch
\uses{lem:volume-density-change}
\leanok
Let $\tau=\varphi_\alpha\circ\varphi_\beta^{-1}$, a diffeomorphism from
$\varphi_\beta(U_\alpha\cap U_\beta)$ onto $\varphi_\alpha(U_\alpha\cap U_\beta)$, and note
that $\tau$ carries $\varphi_\beta(S)$ bijectively onto $\varphi_\alpha(S)$. Its derivative
at $y=\varphi_\beta(x)$ is the matrix $A$ of Lemma~\ref{lem:volume-density-change}. The
change-of-variables formula for the diffeomorphism $\tau$ gives
$$\int_{\varphi_\alpha(S)}\lambda_\alpha\,d\mu
=\int_{\varphi_\beta(S)}\lvert\det A(y)\rvert\;\lambda_\alpha(\tau(y))\,d\mu(y).$$
For $y\in\varphi_\beta(S)$ with $x=\varphi_\beta^{-1}(y)\in S$ we have
$\varphi_\alpha^{-1}(\tau(y))=x$, hence
$\lambda_\alpha(\tau(y))=\sqrt{\det g^\alpha(x)}$, and Lemma~\ref{lem:volume-density-change}
gives
$$\lvert\det A(y)\rvert\,\lambda_\alpha(\tau(y))
=\lvert\det A\rvert\sqrt{\det g^\alpha(x)}=\sqrt{\det g^\beta(x)}=\lambda_\beta(y).$$
The integrand therefore equals $\lambda_\beta$ pointwise on $\varphi_\beta(S)$, and the
right-hand side is $\mu_\beta(S)$.
\end{proof}

\begin{definition}[The Riemannian measure]
\label{def:riemannian-measure}
\lean{MorganTianLib.riemannianMeasure,MorganTianLib.riemannianMeasure_apply_chart}
\uses{def:chart-volume-density, lem:riemannian-measure-well-defined}
\leanok



Let $(M,g)$ be a second-countable Riemannian manifold without boundary. Choose a countable
family of charts $\{(U_{\alpha_n},\varphi_{\alpha_n})\}_{n\in\N}$ covering $M$, and
disjointify it into the Borel partition
$P_n=U_{\alpha_n}\setminus\bigcup_{k<n}U_{\alpha_k}$. The \emph{Riemannian measure} of
$(M,g)$ is the Borel measure
$$\mu_g(S)=\sum_{n}\mu_{\alpha_n}(S\cap P_n).$$
By Lemma~\ref{lem:riemannian-measure-well-defined} it does not depend on the chosen cover, and
it satisfies the defining property: for \emph{every} chart $\alpha$ and every measurable
$S\subseteq U_\alpha$,
$$\mu_g(S)=\int_{\varphi_\alpha(S)}\sqrt{\det g^\alpha\bigl(\varphi_\alpha^{-1}(y)\bigr)}\;dy,$$
i.e. $\mu_g$ is the measure written classically as $\sqrt{\det g_{ij}}\,dx^1\cdots dx^n$. We
write $\Vol(S)=\mu_g(S)$, and in particular $\Vol B(p,r)=\mu_g(B(p,r))$.

Indeed, since $S=\bigsqcup_n (S\cap P_n)$ and each $S\cap P_n$ lies in both $U_{\alpha_n}$ and
$U_\alpha$, Lemma~\ref{lem:riemannian-measure-well-defined} gives
$\mu_{\alpha_n}(S\cap P_n)=\mu_\alpha(S\cap P_n)$; summing over $n$ and using countable
additivity of $\mu_\alpha$ yields $\mu_g(S)=\mu_\alpha(S)$.
\end{definition}

\begin{lemma}[Change of variables for the Riemannian measure]
\label{lem:riemannian-measure-jacobian}
\lean{MorganTianLib.riemannianMeasure_eq_lintegral_jacobian}
\uses{def:riemannian-measure, def:chart-volume-density}
\leanok



Let $S\subseteq U_\alpha$ be measurable and let $\psi\colon U\to M$ be an injective map from a
measurable $U\subseteq E$ with $\varphi_\alpha\circ\psi$ differentiable on $U$ and
$(\varphi_\alpha\circ\psi)(U)=\varphi_\alpha(S)$. Then
$$\Vol(S)=\int_U \bigl\lvert\det D(\varphi_\alpha\circ\psi)_v\bigr\rvert\;
\lambda_\alpha\bigl(\varphi_\alpha(\psi(v))\bigr)\,d\mu(v).$$
The integrand is the \emph{Riemannian Jacobian} of $\psi$: it is independent of the chart
$\alpha$, although neither of its two factors is.
\end{lemma}

\begin{proof}
\sketch
\uses{def:riemannian-measure,lem:riemannian-measure-well-defined}
\leanok
By Definition~\ref{def:riemannian-measure},
$\Vol(S)=\int_{\varphi_\alpha(S)}\lambda_\alpha\,d\mu$. Apply the change-of-variables formula
to the injective differentiable map $\varphi_\alpha\circ\psi$, whose image is
$\varphi_\alpha(S)$: this converts the integral over $\varphi_\alpha(S)$ into the integral
over $U$ of $\lvert\det D(\varphi_\alpha\circ\psi)_v\rvert$ times the integrand evaluated at
$(\varphi_\alpha\circ\psi)(v)$, which is the stated formula. Chart-independence of the
integrand follows from Lemma~\ref{lem:volume-density-change}: replacing $\alpha$ by $\beta$
multiplies $\lambda_\alpha$ by $\lvert\det A\rvert$ and, by the chain rule, multiplies
$\lvert\det D(\varphi_\alpha\circ\psi)\rvert$ by $\lvert\det A\rvert^{-1}$.
\end{proof}

The hypothesis $S\subseteq U_\alpha$ in
Lemma~\ref{lem:riemannian-measure-jacobian} is a real restriction: the
set we must integrate over is $\exp_p\bigl(U_p\cap B(0,r)\bigr)$, and
for $r$ larger than the injectivity radius this leaves every chart. The
following removes the hypothesis. Since the integrand of
Lemma~\ref{lem:riemannian-measure-jacobian} is chart-independent even
though neither factor is, it is meaningful to demand it globally, and
that is what the notion of a Riemannian Jacobian does.

\begin{definition}[Riemannian Jacobian]
\label{def:riemannian-jacobian}
\lean{MorganTianLib.HasRiemannianJacobianOn}
\uses{def:chart-volume-density}
\leanok



Let $\psi\colon U\to M$ with $U\subseteq E$. A function $\rho\colon
U\to\R$ is a \emph{Riemannian Jacobian} for $\psi$ on $U$ if for every
$v\in U$ and every chart $\alpha$ whose domain contains $\psi(v)$, the
map $\varphi_\alpha\circ\psi$ is differentiable at $v$ with derivative
$D$, and
$$\rho(v)=\lvert\det D\rvert\cdot
\lambda_\alpha\bigl(\varphi_\alpha(\psi(v))\bigr).$$
By Lemma~\ref{lem:volume-density-change} the right-hand side does not
depend on $\alpha$, so this is a condition on $\rho$ and not a choice.
\end{definition}

\begin{lemma}[Change of variables, globally]
\label{lem:riemannian-measure-jacobian-global}
\lean{MorganTianLib.riemannianMeasure_image_eq_lintegral_jacobian}
\uses{def:riemannian-measure, def:riemannian-jacobian, lem:riemannian-measure-jacobian}
\leanok



Let $U\subseteq E$ be measurable, $\psi\colon E\to M$ continuous and
injective on $U$, and let $\rho$ be a Riemannian Jacobian for $\psi$ on
$U$. Then
$$\Vol\bigl(\psi(U)\bigr)=\int_U \rho\,d\mu,$$
with no hypothesis on where $\psi(U)$ lies.
\end{lemma}

\begin{proof}
\sketch
\uses{def:riemannian-measure,lem:riemannian-measure-jacobian}
\leanok
Let $\{P_n\}$ be the measurable partition of $M$ underlying
Definition~\ref{def:riemannian-measure}, with $P_n\subseteq
U_{\alpha_n}$, and set $U_n=U\cap\psi^{-1}(P_n)$: a measurable
partition of $U$. Then $\psi(U_n)\subseteq P_n\subseteq U_{\alpha_n}$
lies in a single chart, so
Lemma~\ref{lem:riemannian-measure-jacobian} applies to it and gives
$\Vol(\psi(U_n))=\int_{U_n}\rho\,d\mu$, the integrand being the
Riemannian Jacobian read in the chart $\alpha_n$ --- which is
legitimate precisely because $\rho$ is chart-independent.

The sets $\psi(U_n)$ are pairwise disjoint, being contained in the
pairwise disjoint $P_n$, and their union is $\psi(U)$. Countable
additivity of $\mu_g$ and of the integral therefore give
$$\Vol\bigl(\psi(U)\bigr)=\sum_n\Vol\bigl(\psi(U_n)\bigr)
=\sum_n\int_{U_n}\rho\,d\mu=\int_U\rho\,d\mu.$$
Injectivity of $\psi$ on $U$ is used only to apply
Lemma~\ref{lem:riemannian-measure-jacobian} on each piece; it restricts
to each $U_n$.
\end{proof}

\begin{lemma}[The Riemannian Jacobian of the exponential map]
\label{lem:exponential-riemannian-jacobian}
\lean{MorganTianLib.expRiemannianJacobian,MorganTianLib.hasRiemannianJacobianOn_expMapGlobal, MorganTianLib.riemannianMeasure_image_segmentDomain_eq_lintegral}
\uses{def:riemannian-jacobian, def:exponential-map, lem:exponential-differential-jacobi, lem:volume-density-change, def:cut-time-locus, prop:exponential-diffeomorphism-cut-locus}



Fix $p\in M$. The exponential map $\exp_p\colon T_pM\to M$ carries a Riemannian
Jacobian on all of $T_pM$, namely the function
$$\rho_p(v)=\bigl\lvert\det D(\varphi_{\exp_p(v)}\circ\exp_p)_v\bigr\rvert\cdot
  \lambda_{\exp_p(v)}\bigl(\varphi_{\exp_p(v)}(\exp_p(v))\bigr),$$
read in the preferred chart centred at $\exp_p(v)$. Consequently, for every measurable
$U\subseteq U_p$ contained in the segment domain (Definition~\ref{def:cut-time-locus}),
$$\Vol\bigl(\exp_p(U)\bigr)=\int_U\rho_p\,d\mu.$$
\end{lemma}

\begin{proof}
\sketch
\uses{def:riemannian-jacobian, lem:exponential-differential-jacobi, lem:volume-density-change, lem:riemannian-measure-jacobian-global, prop:exponential-diffeomorphism-cut-locus}
Choose orthonormal frames at the source and target. In those frames the Riemannian Jacobian is the absolute determinant of the differential. The radial frame identifies the angular block with the Jacobi matrix and the radial block with the unit vector by the Gauss lemma. Thus the determinant equals the polar volume density, with the coordinate Gram factors cancelling by the transformation law.
\end{proof}


\begin{lemma}[Polar coordinates for a ball integral]
\label{lem:polar-coordinates}
\lean{MorganTianLib.lintegral_eq_polar,MorganTianLib.setLIntegral_ball_eq_polar}
\leanok


Let $E$ be a real normed space of finite dimension $n\ge1$, let $\mu$ be a Haar
measure on $E$, let $S\subset E$ be the unit sphere, and let $\mu_S$ be the
induced measure on $S$ (so that $\mu_S(A)=n\cdot\mu\bigl((0,1)\cdot A\bigr)$).
Then for every measurable $f\colon E\to[0,\infty]$ and every $r$,
$$\int_{B(0,r)}f\,d\mu
=\int_{S}\int_0^r t^{n-1}f(t\omega)\,dt\,d\mu_S(\omega).$$
\end{lemma}

\begin{proof}
\sketch
\leanok
The origin is $\mu$-null, so the integral is unchanged by restricting to
$E\setminus\{0\}$. The map $x\mapsto(x/|x|,|x|)$ is a homeomorphism
$E\setminus\{0\}\to S\times(0,\infty)$ carrying $\mu$ to
$\mu_S\otimes t^{n-1}dt$; pushing the integral forward along it and applying
Tonelli's theorem gives the identity over all of $E$. For the ball, apply this
to the function $\mathbf 1_{B(0,r)}f$ and note that for a \emph{unit} $\omega$ and
$t>0$ one has $|t\omega|=t$, so the constraint $t\omega\in B(0,r)$ is exactly
$t<r$; the inner integral therefore runs over $(0,r)$.
\end{proof}

\begin{lemma}[The model polar volume]
\label{lem:model-ball-volume}
\lean{MorganTianLib.modelBallVolume_eq,MorganTianLib.model_polar_volume_identity}
\uses{def:sinh-comparison-function, lem:polar-coordinates}



Write $\rho_k(v)=\bigl(\sn_k(|v|)/|v|\bigr)^{n-1}$ for the volume density of the
exponential map of the model space $H^n_k$ (Lemma~\ref{lem:constant-curvature-jacobi}),
and $\Vol B_{H^n_k}(q_k,r)=\int_{B(0,r)}\rho_k\,d\mu$. Then
$$\Vol B_{H^n_k}(q_k,r)=\omega_{n-1}\int_0^r\sn_k^{n-1}(t)\,dt,
\qquad \omega_{n-1}:=\mu_S(S).$$
\end{lemma}

\begin{proof}
\sketch
\uses{lem:polar-coordinates}
Apply Lemma~\ref{lem:polar-coordinates} to $f=\rho_k$. Since $\rho_k$ depends only
on $|v|$, the inner integrand is $t^{n-1}\rho_k(t\omega)=t^{n-1}(\sn_k(t)/t)^{n-1}
=\sn_k^{n-1}(t)$, independent of $\omega$: the radial weight $t^{n-1}$ of polar
coordinates cancels the $|v|^{-(n-1)}$ in the density. The angular integral is then
the constant $\mu_S(S)=\omega_{n-1}$.
\end{proof}

\begin{lemma}[Bishop-Gromov for the volume of a ball]
\label{lem:bishop-gromov-ball}
\lean{MorganTianLib.bishop_gromov_ball,MorganTianLib.bishop_gromov_ball_ratio}
\uses{lem:polar-coordinates, lem:model-ball-volume, def:sinh-comparison-function}
\leanok



Let $k\ge0$ and $R>0$. Let $\rho\ge0$ be a measurable density on $T_pM$ (thought of
as $\rho(v)=|\det d(\exp_p)_v|$, extended by $0$ past the cut time), and put
$\Vol(r)=\int_{B(0,r)}\rho\,d\mu$. Suppose that in every unit direction
$\omega\in S$ the polar density $\nu_\omega(t)=t^{n-1}\rho(t\omega)$ has
non-increasing ratio to the model density,
$$t\ \longmapsto\ \frac{\nu_\omega(t)}{\sn_k^{n-1}(t)}
\quad\text{non-increasing on }(0,R).$$
Then for $0<r_1\le r_2<R$
$$\Vol(r_2)\cdot\Vol B_{H^n_k}(q_k,r_1)\ \le\
\Vol(r_1)\cdot\Vol B_{H^n_k}(q_k,r_2),$$
equivalently $r\mapsto\Vol(r)/\Vol B_{H^n_k}(q_k,r)$ is non-increasing on $(0,R)$.
\end{lemma}

\begin{proof}
\sketch
\uses{lem:polar-coordinates,lem:model-ball-volume}
\leanok
Apply the radial Bishop--Gromov inequality direction by direction, extending each radial density by zero after its cut time. Integrating over the unit sphere and using polar coordinates converts the radial integrals into ball volumes. The ratio-of-integrals lemma then gives monotonicity of the quotient by the model ball volume.
\end{proof}


\begin{theorem}[Bishop-Gromov relative volume comparison]
\label{thm:bishop-gromov}
\lean{MorganTianLib.bishop_gromov_manifold_ratio,MorganTianLib.bishop_gromov_manifold_with_producers}
\uses{def:ricci-curvature, thm:ricci-curvature-comparison, lem:bishop-gromov-ball, def:riemannian-measure}
\label{BishopGromov}


(Relative Volume Comparison, Bishop-Gromov 1964-1980) Let $(M,g)$ be a
connected Riemannian manifold and fix $p\in M$. Put $n=\dim M$, assume
$n\ge2$, and suppose that $B(p,R)$ has compact closure in $M$ for some
$R>0$. Assume that the metric distance is the Riemannian distance induced
by $g$ and that $(M,g)$ is complete. Assume that
\[
\Ric_x(v,v)\ge -(n-1)k\,g_x(v,v)
\quad(x\in\overline{B(p,R)},\ v\in T_xM),
\]
and that the transported exponential Jacobian is measurable. Then the
genuine manifold ratio
\[
r\longmapsto
\frac{\mu_g(B(p,r))}{\operatorname{ModelVol}_k(r)},
\qquad
\operatorname{ModelVol}_k(r):=\operatorname{modelBallVolume}
 (\operatorname{volume},k,r),
\]
is non-increasing on $(0,R)$, where $\mu_g$ is the Riemannian measure
constructed from the metric and the chosen tangent Haar measure.

For the source normalization and its flat-model corollary, additionally
assume that the small-radius ratio tends to the origin density, that the
flat model volume is a constant multiple of $r^n$, and that the origin
density is $1$. If moreover $\Ric\ge0$
on the closed ball, the ratio has limit $1$ at $0$ and
$r\mapsto\mu_g(B(p,r))/r^n$ is non-increasing.

\end{theorem}

\begin{proof}
\sketch
\uses{lem:volume-element-comparison,lem:geodesic-polar-form,lem:localized-cut-locus,lem:model-polar-isometry,lem:normal-ball}
For almost every radial direction, Ricci comparison makes the ratio of the radial Jacobian to the model density nonincreasing until the cut time. Extend it by zero after the cut time and integrate first in radius and then over the unit sphere. The elementary ratio-of-integrals lemma preserves monotonicity under both integrations. Polar coordinates and nullity of the cut locus identify the resulting integrals with the volumes of the manifold and model balls, and the small-radius normalization gives the limiting value one.
\end{proof}

\section{Local volume and the injectivity radius}

As the following results show, in the presence of bounded curvature the volume
of a ball $B(p,r)$ in $M$ is bounded away from zero if and only if the
injectivity radius of $M$ at $p$ is bounded away from zero.
The arguments in this section repeatedly rescale the metric by a
constant; the effect of such a rescaling on the geometric
quantities in play is recorded in the following lemma.

\begin{lemma}[Rescaling the metric]
\label{lem:metric-rescaling}
\lean{MorganTianLib.metricRescalingLaws,MorganTianLib.rescaledMetric_intrinsicMetricBookInjectivityRadius,MorganTianLib.rescaledMetric_canonicalMetricBookInjectivityRadius,MorganTianLib.canonicalMetricBookInjectivityRadius_eq}
\uses{def:riemannian-metric, def:sectional-curvature, def:curvature-operator-norm, def:ricci-curvature, def:injectivity-radius, def:exponential-map}

Let $(M,g)$ be a Riemannian manifold and let $c>0$ be a constant;
set $g'=c\,g$. Then:
\begin{enumerate}
\item $g$ and $g'$ have the same Christoffel symbols, the same
geodesics (as parameterized curves), the same exponential maps, and
the same $(1,3)$ curvature tensor; the $(0,4)$ tensor satisfies
$R'_{ijkl}=c\,R_{ijkl}$;
\item $K'(P)=K(P)/c$ for every $2$-plane $P$, and
$|\Rm'(x)|=|\Rm(x)|/c$ for the curvature operator norm of
Definition~\ref{def:curvature-operator-norm};
\item $\Ric'=\Ric$ as $(0,2)$-tensors; in particular
$\Ric\ge(n-1)k\,g$ holds if and only if
$\Ric'\ge(n-1)(k/c)\,g'$;
\item $d'=\sqrt{c}\,d$; hence $B_{g'}(p,\sqrt{c}\,r)=B_g(p,r)$,
and $(M,g')$ is complete if and only if $(M,g)$ is;
\item $\Vol_{g'}=c^{n/2}\Vol_g$;
\item The intrinsic exponential-ball radius scales by
$\sqrt{c}$:
\[
\operatorname{inj}^{\rm int}_{(M,g')}(p)=
\sqrt{c}\,\operatorname{inj}^{\rm int}_{(M,g)}(p).
\]
For the canonical metric-ball radius the same scaling is available under
the completeness hypothesis, and that canonical radius is definitionally
the metric-ball radius used above. These ball-based radii are distinguished
from the cut-time radius $\inj_M^{\rm cut}(p)$ of Definition~\ref{def:injectivity-radius}.
\end{enumerate}
\end{lemma}

\begin{proof}
\sketch
\uses{thm:levi-civita-connection,def:riemann-curvature-tensor,def:geodesic,def:sectional-curvature,def:curvature-operator,def:curvature-operator-norm,def:ricci-curvature,def:injectivity-radius,def:exponential-map}
For $\tilde g=cg$, the inverse metric contributes $c^{-1}$ and derivatives of the metric contribute $c$, so the Christoffel symbols and Levi--Civita connection are unchanged. Consequently geodesics and the $(1,3)$ curvature tensor are unchanged, whereas the $(0,4)$ tensor scales by $c$. Tracing and comparing the induced norms gives the stated scaling laws for sectional, Ricci, scalar, and curvature-operator quantities. Curve lengths scale by $\sqrt c$, yielding the distance, ball, volume, and injectivity-radius formulas.
\end{proof}


\begin{proposition}[Bounded curvature implies volume-injectivity relation]
\label{prop:injectivity-radius-volume}
\uses{def:injectivity-radius, def:curvature-operator-norm}
\notready
\label{injvol}

Fix an integer $n>0$. For every $\epsilon>0$ there is $\delta>0$
depending on $n$ and $\epsilon$ such that the following holds.
Suppose that $(M^n,g)$ is a complete connected Riemannian
manifold of dimension $n$ and that $p\in M$. Suppose that $|\Rm(x)|\le
r^{-2}$ for all $x\in B(p,r)$. If the injectivity radius of $M$ at
$p$ is at least $\epsilon r$, then $\Vol(B(p,r))\ge \delta
r^n$.
\end{proposition}

\begin{proof}
\sketch
\uses{lem:rauch-lower,def:curvature-operator-norm,lem:geodesic-polar-form}
Rescale so that the curvature bound and radius are normalized. If the injectivity radius is already comparable to the radius there is nothing to prove. Otherwise Klingenberg's lemma supplies a short geodesic loop. The pulled-back normal ball contains many disjoint translates associated with successive multiples of this loop; metric and volume comparison bound each from below and place them inside a fixed larger ball. Counting these disjoint pieces gives the stated relation between volume and injectivity radius, and scaling restores the general formula.
\end{proof}

 We shall normally work with volume, which behaves nicely under Ricci
flow, but in order to take limits we need to bound the injectivity
radius away from zero. Thus, the more important, indeed crucial,
result for our purposes is the converse to the previous proposition;
see Theorem 4.3, especially Inequality (4.22), on page 46 of
\cite{CheegerGromovTaylor}, or see Theorem 5.8 on page 96 of
\cite{CheegerEbin}.

Our proof follows Section 4 of \cite{CheegerGromovTaylor}
(Theorem 4.3 there, via their Lemmas 4.4--4.6), with all steps
carried out in full; the one point which is asserted without proof
there --- the strict convexity of squared-distance functions in the
ball with the pulled-back metric --- is established below
(Lemma~\ref{lem:pullback-unique-minimizers}) by a homotopy argument
special to the pulled-back ball. The mechanism is a count of
preimages of points under the exponential map: a short geodesic
loop at $p$ forces $p$ (and then every nearby point) to have many
preimages in a ball of the tangent space on which
$\exp_p$ is non-singular, and each preimage carries its own copy
of the local volume; an upper bound for the total volume upstairs
then forces the volume downstairs to be small.

\begin{lemma}[The pulled-back ball and lifts of short curves]
\label{lem:pullback-ball}
\source{cheeger-gromov-taylor:page-0032}
\uses{def:exponential-map, def:geodesic, lem:geodesic-polar-form, def:local-isometry}
\notready


Let $p\in M$, let $0<\rho_*\le\infty$, and suppose that
$\exp_p$ is defined and non-singular on
$B=B(0,\rho_*)\subset T_pM$. Let $\tilde g=\exp_p^*g$ on $B$.
Then:
\begin{enumerate}
\item $\exp_p\colon (B,\tilde g)\to M$ is a local isometry; the
rays $t\mapsto tv$ are unit-speed $\tilde g$-geodesics (for unit
$v$); every piecewise smooth curve $c$ in $B$ satisfies
$L_{\tilde g}(c)\ge\bigl||c(1)|-|c(0)|\bigr|$; and the
$\tilde g$-geodesics starting at $0$ are exactly the rays.
\item (Path lifting) For every continuous curve $c\colon[0,1]\to M$
of finite length $L(c)$ with $c(0)=\exp_p(y_0)$ for some
$y_0\in B$ with $|y_0|+L(c)<\rho_*$, there is a unique continuous
curve $\tilde c$ (the \emph{lift of $c$ at $y_0$}) with $\tilde
c(0)=y_0$ and $\exp_p\circ\tilde c=c$; it satisfies $|\tilde
c(t)|\le|y_0|+L(c|_{[0,t]})$. The lift of the reversed curve at
$\tilde c(1)$ is the reversal of $\tilde c$, and lifts of
concatenations are the concatenations of the successive lifts.
\item (Homotopy invariance) Let $H\colon[0,1]\times[0,1]\to M$ be
continuous, write $c_u=H(\cdot,u)$, and suppose that all $c_u$ have
the same endpoints $q_0=\exp_p(y_0)$ and $q_1$, and finite
lengths $L(c_u)\le\Lambda$ with $|y_0|+\Lambda<\rho_*$. Then the
lift endpoints $\tilde c_u(1)$ coincide for all $u$.
\item (Endpoint map and cancellation) For a continuous curve $c$
from $p$ of finite length $L(c)<\rho_*$ set
$E(c):=\tilde c(1)\in\exp_p^{-1}(c(1))$, the endpoint of the
lift of $c$ at $0$; write $\mathrm{rad}_y$, $y\in B$, for the curve
$t\mapsto\exp_p(ty)$, so that $E(\mathrm{rad}_y)=y$. Then, with
all length sums below assumed $<\rho_*$ and $\cup$ denoting
concatenation (first curve first):
\begin{enumerate}
\item[(a)] $c$ is homotopic to $\mathrm{rad}_{E(c)}$ with fixed
endpoints through curves of length $\le L(c)$;
\item[(b)] if $E(\theta_1\cup c)=E(\theta_2\cup c)$ for curves
$\theta_1,\theta_2$ from $p$ to $q'$ and $c$ from $q'$, then
$E(\theta_1)=E(\theta_2)$;
\item[(c)] $E(\theta\cup c)$ is the endpoint of the lift of $c$ at
$E(\theta)$; in particular it depends on $\theta$ only through
$E(\theta)$;
\item[(d)] if $\sigma$ is a curve from $p$ with
$E(\sigma)=y$, then $E(\theta\cup\sigma)
=E(\theta\cup\mathrm{rad}_y)$ for every loop $\theta$ at $p$.
\end{enumerate}
\item (Deck motions) For a loop $\theta$ at $p$ of length
$\ell_\theta$ and $a>0$ with $a+\ell_\theta<\rho_*$, define
$$\pi_\theta\colon B(0,a)\to B(0,a+\ell_\theta),\qquad
\pi_\theta(y):=E(\theta\cup\mathrm{rad}_y).$$
Then $\exp_p\circ\pi_\theta=\exp_p$; the map
$\pi_\theta$ is smooth with everywhere isometric differential, so
$L_{\tilde g}(\pi_\theta\circ\omega)=L_{\tilde g}(\omega)$ for
every curve $\omega$ in $B(0,a)$; moreover
$\pi_{\theta'}\circ\pi_\theta=\pi_{\theta'\cup\theta}$ and
$\pi_{\bar\theta}\circ\pi_\theta=\mathrm{id}$ on suitable balls
(with $\bar\theta$ the reversed loop), and
$\pi_\theta(E(\sigma))=E(\theta\cup\sigma)$ for curves $\sigma$
from $p$.
\end{enumerate}
\end{lemma}

\begin{proof}
\sketch
\uses{lem:geodesic-polar-form,def:local-isometry,def:exponential-map,lem:local-isometry-exponential,lem:local-isometry-normal-balls,lem:normal-ball}
Use the exponential map as a local isometry on the pulled-back metric ball. Short curves lift successively through normal neighborhoods; uniqueness on overlaps glues the local lifts and preserves their lengths. Radial geodesics supply existence at the center, while the local-isometry exponential identity gives the endpoint and velocity formulas. Reversal and concatenation follow from uniqueness of lifts.
\end{proof}

\begin{lemma}[Unique minimizers and convexity in the pulled-back ball]
\label{lem:pullback-unique-minimizers}
\uses{lem:pullback-ball, def:sectional-curvature, lem:rauch-lower, lem:conjugate-sturm}
\notready

In the setting of Lemma~\ref{lem:pullback-ball} with $\rho_*=1$,
suppose in addition that $|K(P)|\le1$ for every $2$-plane tangent
to $(B,\tilde g)$. Write $\widetilde\exp_z$ for the exponential
map of $(B,\tilde g)$ at $z$. Then for every $z\in
\overline{B(0,1/32)}$:
\begin{enumerate}
\item $\widetilde\exp_z$ is defined and non-singular on
$B(0,1/2)\subset T_z B$, and injective on $B(0,1/8)$;
\item every $y\in B$ with $d_{\tilde g}(z,y)\le1/10$ lies in
$\Omega_z:=\widetilde\exp_z(B(0,1/8))$; in particular any two
points $z,y\in\overline{B(0,1/32)}$ are joined by a
$\tilde g$-minimizing geodesic, of length $\le1/8$, unique among
geodesics from $z$ to $y$ of length $\le 1/8$; and on
$\Omega_z$ the distance
$d_z:=d_{\tilde g}(z,\cdot)$ satisfies
$d_z=|\widetilde\exp_z^{-1}|$ and is smooth away from $z$;
\item $\Hess_{\tilde g}(d_z^2)\ge\tilde g$ on $\Omega_z$.
\end{enumerate}
(Here $d_{\tilde g}$ denotes the infimum of $\tilde g$-lengths of
curves in $B$.)
\end{lemma}

\begin{proof}
\sketch
\uses{lem:pullback-ball,lem:rauch-lower,lem:conjugate-sturm,lem:geodesic-polar-form,lem:scalar-riccati-comparison,lem:localized-cut-locus,lem:exponential-differential-jacobi}
\emph{Two-sided comparison with the Euclidean metric.} On the cone
over the full unit sphere with $r_0=1/2$,
Lemma~\ref{lem:geodesic-polar-form} applies to $\tilde g$ around
$0$ (non-singularity is our hypothesis). The hypothesis $K\ge-1$
gives, exactly as in the shape-operator estimate in the proof of
Theorem~\ref{thm:sectional-curvature-comparison} (the parallel
eigenvector trick together with
Lemma~\ref{lem:scalar-riccati-comparison}), the bound $A(r)\le
\frac{\sn_1'(r)}{\sn_1(r)}\mathrm{Id}$ and hence, integrating
$\partial_rg_{ij}=2S_{ij}$ as there,
$g_{ij}(r,\theta)\le\sn_1^2(r)\hat g_{ij}(\theta)$; the hypothesis
$K\le1$ gives $g_{ij}(r,\theta)\ge s_1^2(r)\hat g_{ij}(\theta)$ by
Lemma~\ref{lem:rauch-lower}(2). Since
$\sn_1^2(r)\le r^2(1+r^2)$ and $s_1^2(r)\ge r^2(1-r^2/3)$ for
$0\le r\le 1/2$, comparing with
$g_{\mathrm{euc}}=dr^2+r^2g_{S^{n-1}}$ yields
\begin{equation}\label{twosided}
0.9\,g_{\mathrm{euc}}\ \le\ \tilde g\ \le\ 1.3\,g_{\mathrm{euc}}
\qquad\text{on }B(0,1/2)
\end{equation}
(as quadratic forms; the radial parts agree and
$1-r^2/3\ge0.9$, $1+r^2\le1.3$ there). In particular, for
$z,y\in\overline{B(0,1/32)}$, the straight segment gives
$d_{\tilde g}(z,y)\le1.2\,|z-y|\le1.2\cdot\tfrac{2}{32}\le0.08$.

\emph{Geodesics from $z$.} We apply
Lemma~\ref{lem:localized-cut-locus}(1) to the Riemannian manifold
$(B,\tilde g)$ at the point $z$, with $R=0.6$. Its hypothesis
holds: every point of the metric ball $B_{\tilde g}(z,0.6)$ is
joined to $z$ by a curve of $\tilde g$-length $<0.6$, hence lies
in $\{|x|\le|z|+0.6\}$ by the radial estimate of
Lemma~\ref{lem:pullback-ball}(1); as $|z|+0.6\le1/32+0.6<1$, the
closure of $B_{\tilde g}(z,0.6)$ is a closed subset of this
compact subset of $B$, hence compact. Therefore every unit-speed
$\tilde g$-geodesic from $z$ extends to every length less than
$0.6$; in particular $\widetilde\exp_z$ is defined on $B(0,1/2)$.
It is non-singular there: by Lemma~\ref{lem:conjugate-sturm}
($K\le1$ along all these geodesics and $1/2<\pi$) there is no
conjugate point along any of these geodesics at parameters
below $1/2$, so $d(\widetilde\exp_z)_v$ is non-singular for
$|v|<1/2$ by Lemma~\ref{lem:exponential-differential-jacobi}.

\emph{Injectivity of $\widetilde\exp_z$ on $B(0,1/8)$.} Suppose
$\sigma_1,\sigma_2$ are distinct $\tilde g$-geodesics from $z$ to a
common point $y'$, of lengths $L_1,L_2\le1/8$. Both lie in
$\{|x|\le5/32\}$. Consider the loop $\ell=\sigma_1\cup\bar\sigma_2$
at $z$, of length $\le1/4$. Using the linear structure of $B\subset
T_pM$, the map $(u,t)\mapsto u\,\ell(t)$ is a free homotopy from
$\ell$ to the constant loop at $0$, through loops based at $uz$
contained in $\{|x|\le5/32\}$ and of $\tilde g$-length at most
$1.2\cdot L_{\mathrm{euc}}(u\ell)=1.2\,u\,L_{\mathrm{euc}}(\ell)
\le1.2\cdot\left(1.1\cdot\tfrac14\right)\le0.33$
by (\ref{twosided}). Converting to a based homotopy at $z$ by
conjugating with the track $u\mapsto uz$ (of $\tilde g$-length
$\le1.2\,|z|\le0.04$) exhibits $\ell$ as fixed-endpoint homotopic,
through loops at $z$ of length $\le0.33+2\cdot0.04\le0.45$, to the
loop $\bar\tau\cup\tau$ ($\tau$ the track, run from $0$ to $z$),
which retracts to the constant loop through loops of length
$\le0.08$. Now apply
Lemma~\ref{lem:pullback-ball} \emph{to the manifold $(B,\tilde g)$
at the point $z$} with $\rho_*=1/2$: its hypotheses hold by the
previous paragraph, all the loops of the homotopy have length
$\le0.45<1/2$, so by parts (3) and (2) of that lemma the lift of
$\ell$ at $0\in T_zB$ is closed:
$$E_z(\sigma_1\cup\bar\sigma_2)=E_z(\text{constant})=0.$$
But $\sigma_1,\sigma_2$ are geodesics \emph{starting at $z$}, so
their lifts at $0$ are the rays to $w_i:=L_i\,\sigma_i'(0)$
(Lemma~\ref{lem:pullback-ball}(1) applied to $(B,\tilde g)$ at
$z$); and the lift of $\ell$ is the lift of $\sigma_1$ (ending at
$w_1$) followed by the lift of $\bar\sigma_2$ at $w_1$. The lift of
$\ell$ being closed means that the lift of $\bar\sigma_2$ at $w_1$
ends at $0$; reversing (Lemma~\ref{lem:pullback-ball}(2)), the lift
of $\sigma_2$ at $0$, which is the ray to $w_2$, ends at $w_1$;
hence $w_2=w_1$, so $L_1=L_2$ and
$\sigma_1'(0)=\sigma_2'(0)$, and $\sigma_1=\sigma_2$ ---
a contradiction. This proves both the injectivity in (1) and the
uniqueness in (2).

\emph{Existence of minimizers, and $d_z$.} We apply
Lemma~\ref{lem:localized-cut-locus}(2) to $(B,\tilde g)$ at $z$,
now with $R=1/9$: as above, the closure of the metric ball
$B_{\tilde g}(z,1/9)$ is compact (it lies in
$\{|x|\le1/32+1/9\}$, by the radial estimate). Hence every $y\in
B$ with $d_{\tilde g}(z,y)\le1/10<1/9$ is joined to $z$ by a
minimizing $\tilde g$-geodesic $\sigma$, of length $d_{\tilde
g}(z,y)\le1/10<1/8$; writing $w=L(\sigma)\,\sigma'(0)$ (for
unit-speed $\sigma$) we get $y=\widetilde\exp_z(w)$ with
$|w|<1/8$, i.e., $y\in\Omega_z$. For
$z,y\in\overline{B(0,1/32)}$ we have $d_{\tilde
g}(z,y)\le0.08\le1/10$, so such $z,y$ are joined by a minimizing
geodesic of length $\le0.08<1/8$; its uniqueness among geodesics
of length $\le1/8$ was proved above.

Next, $d_z=|\widetilde\exp_z^{-1}|$ on all of $\Omega_z$: let
$y=\widetilde\exp_z(w)$ with $|w|<1/8$. The radial geodesic gives
$d_z(y)\le|w|$. Conversely, no curve $c$ from $z$ to $y$ has
$L(c)<|w|$: for such a $c$, apply
Lemma~\ref{lem:pullback-ball}(2) to $(B,\tilde g)$ at $z$ with
$\rho_*=1/2$ (its hypotheses --- $\widetilde\exp_z$ defined and
non-singular on $B(0,1/2)$ --- hold by the second paragraph, and
$0+L(c)<1/8<1/2$): the lift of $c$ at $0\in T_zB$ ends at a
preimage of $y$ in $\overline{B(0,L(c))}\subset B(0,1/8)$, which
by the injectivity of $\widetilde\exp_z$ on $B(0,1/8)$ must be
$w$; but then the radial estimate
(Lemma~\ref{lem:pullback-ball}(1) for $(B,\tilde g)$ at $z$)
gives $L(c)\ge|w|$, a contradiction. Therefore $d_z(y)=|w|$,
i.e., $d_z=|\widetilde\exp_z^{-1}|$ on $\Omega_z$, which is
smooth away from $z$ since $\widetilde\exp_z$ is a diffeomorphism
of $B(0,1/8)$ onto $\Omega_z$ (non-singular and injective there).

(3) In geodesic polar coordinates around $z$ in $(B,\tilde g)$
(furnished on $\Omega_z\setminus\{z\}$ by the diffeomorphism just
established), $d_z$ is the radial coordinate, so
Lemma~\ref{lem:rauch-lower}(3) with $\bar K=1$ gives
$$\Hess_{\tilde g}\left(\tfrac12d_z^2\right)
\ge\min\left(1,\ d_z\,\frac{s_1'(d_z)}{s_1(d_z)}\right)\tilde g
=\min\left(1,\,d_z\cot d_z\right)\tilde g
\ge\left(1-\tfrac{d_z^2}{2}\right)\tilde g\ge\tfrac12\,\tilde g$$
on $\Omega_z\setminus\{z\}$ (where $d_z<1/8$). Since
$d_z^2=|\widetilde\exp_z^{-1}|^2$ is smooth on all of $\Omega_z$
(the squared norm is smooth at the origin) and both sides of the
inequality $\Hess_{\tilde g}(d_z^2)\ge\tilde g$ are continuous,
the bound extends from the dense subset
$\Omega_z\setminus\{z\}$ to all of $\Omega_z$, which is assertion
(3).
\end{proof}

\begin{lemma}[Distinctness of the lifted multiples of a geodesic loop]
\label{lem:loop-multiples-distinct}
\source{cheeger-gromov-taylor:page-0033}
\uses{lem:pullback-ball, lem:pullback-unique-minimizers}


In the setting of Lemma~\ref{lem:pullback-unique-minimizers}
($\rho_*=1$, $|K(P)|\le1$ on $(B,\tilde g)$), let $\gamma$ be a
geodesic loop at $p$ of length $2\ell>0$, smooth on the open
parameter interval, and let $N\ge1$ be an integer with
$2\ell N\le1/256$. Then the points
$$E(\gamma^0)=0,\ E(\gamma^1),\ \ldots,\ E(\gamma^N)$$
(where $\gamma^m$ denotes the $m$-fold concatenation) are pairwise
distinct points of $\overline{B(0,1/256)}$, all mapping to $p$
under $\exp_p$.
\end{lemma}

\begin{proof}
\sketch
\uses{lem:pullback-ball,lem:pullback-unique-minimizers}
Each $\gamma^m$, $m\le N$, has length $2\ell m\le1/256<1$, so
$x_m:=E(\gamma^m)$ is defined, $|x_m|\le2\ell m\le1/256$
(Lemma~\ref{lem:pullback-ball}(2)) and
$\exp_p(x_m)=\gamma^m(\text{end})=p$.

First note that $E(\gamma)\not=0$: since $\gamma$ is smooth on the
open interval, its lift at $0$ is an unbroken $\tilde g$-geodesic
starting at $0$, hence a ray
(Lemma~\ref{lem:pullback-ball}(1)), so
$E(\gamma)=2\ell\,\gamma'(0)\not=0$.

Suppose now that $x_m=x_{m'}$ with $0\le m<m'\le N$. Writing
$\gamma^{m'}=\gamma^{i}\cup\gamma^{m}$ with $i:=m'-m\ge 1$ and
$\gamma^m$ as the common tail, cancellation
(Lemma~\ref{lem:pullback-ball}(4b)) gives $E(\gamma^i)=E(\text{constant
loop})=0$.

Let $\pi:=\pi_\gamma$ be the deck motion of
Lemma~\ref{lem:pullback-ball}(5); we shall use it, together with
its powers and inverse, on the ball $\overline{B(0,0.45)}$ (all
defined, since $0.45+2\ell N\le0.45+1/256<1$). We derive a
contradiction from $E(\gamma^i)=0$ in three steps. Throughout,
$d=d_{\tilde g}$, and we use the radial estimate of
Lemma~\ref{lem:pullback-ball}(1): every curve $c$ in $B$ satisfies
$L_{\tilde g}(c)\ge\bigl||c(1)|-|c(0)|\bigr|$, so in particular
$d(y_1,y_2)\ge\bigl||y_1|-|y_2|\bigr|$.

\emph{$\pi$ has no fixed point.} If $\pi(y)=y$ then
$E(\gamma\cup\mathrm{rad}_y)=y=E(\mathrm{rad}_y)$, so cancellation
with the common tail $\mathrm{rad}_y$ gives $E(\gamma)=0$,
contradicting the first paragraph.

\emph{$\pi$ permutes a finite orbit.} By
Lemma~\ref{lem:pullback-ball}(5) and (4d),
$$\pi(x_j)=E\left(\gamma\cup\mathrm{rad}_{E(\gamma^j)}\right)
=E\left(\gamma\cup\gamma^j\right)=x_{j+1},$$
and $x_i=E(\gamma^i)=0=x_0$; so $\pi$ permutes the finite set
$O:=\{x_0,\ldots,x_{i-1}\}\subset\overline{B(0,1/256)}$
cyclically. Moreover, by the composition rule,
$\pi^{j}=\pi_{\gamma^{j}}$ for $1\le j\le i$, and for
$y\in\overline{B(0,0.4)}$,
$$\pi_{\gamma^i}(y)=E\left(\gamma^i\cup\mathrm{rad}_y\right)$$
is the endpoint of the lift of $\mathrm{rad}_y$'s image at
$E(\gamma^i)=0$ (Lemma~\ref{lem:pullback-ball}(4c)), i.e., equals
$y$: thus $\pi^i=\mathrm{id}$ on $\overline{B(0,0.4)}$. The same
applies to the reversed loop: $\pi^{-1}=\pi_{\bar\gamma}$ and
$\pi_{\bar\gamma^{\,j}}=\pi^{-j}$, with
$\pi^{-i}=\mathrm{id}$ there as well.

In particular the powers of $\pi$ do not increase distances at the
scales we use: for $y_1,y_2\in\overline{B(0,0.2)}$ with
$d(y_1,y_2)\le0.2$ and $|j|\le i$,
$$d\left(\pi^{j}(y_1),\pi^{j}(y_2)\right)\le d(y_1,y_2).$$
Indeed, for $\epsilon\in(0,0.01)$ take a curve
$\omega$ from $y_1$ to $y_2$ of $\tilde g$-length $\le
d(y_1,y_2)+\epsilon\le0.21$; by the radial estimate $\omega$
stays in $B(0,0.2+0.22)\subset B(0,0.45)$, where $\pi^{j}$ is
defined with everywhere isometric differential
(Lemma~\ref{lem:pullback-ball}(5), applied to the loop
$\gamma^{|j|}$ or $\bar\gamma^{\,|j|}$, of length
$2\ell|j|\le1/256$). Hence $\pi^{j}\circ\omega$ joins
$\pi^{j}(y_1)$ to $\pi^{j}(y_2)$ with the same length, so
$d(\pi^{j}y_1,\pi^{j}y_2)\le d(y_1,y_2)+\epsilon$, and letting
$\epsilon\to0$ the claim follows. Combining the two directions
gives the equality
\begin{equation}\label{deckisometry}
d\left(\pi^{j}(y_1),\pi^{j}(y_2)\right)
=d(y_1,y_2),\qquad
y_1,y_2\in\overline{B(0,1/7)},\ d(y_1,y_2)\le0.2,\
|j|\le i:
\end{equation}
here $\le$ is the contraction inequality; and since
$|\pi^{j}(y_k)|\le2\ell|j|+|y_k|\le1/256+1/7\le0.2$
(Lemma~\ref{lem:pullback-ball}(2)) and
$d(\pi^{j}y_1,\pi^{j}y_2)\le0.2$, the contraction inequality
applied to the points $\pi^{j}(y_k)$ with the power $-j$, together
with $\pi^{-j}\circ\pi^{j}=\pi_{\bar\gamma^{\,j}\cup\gamma^{j}}
=\mathrm{id}$ on $\overline{B(0,1/7)}$
(the composition rule and the reverse--retrace homotopy of
Lemma~\ref{lem:pullback-ball}(5)), gives the reverse inequality
$d(y_1,y_2)\le d(\pi^{j}y_1,\pi^{j}y_2)$.

\emph{The center of mass.} Set
$$F(y):=\sum_{j=0}^{i-1}d\left(x_j,y\right)^2,
\qquad y\in\overline{B(0,1/16)},$$
a continuous function on a compact set; let $y^*$ be a minimum
point, and write $m=F(y^*)$. Since all
$x_j\in\overline{B(0,1/256)}$, and since for any $y\in B$ the ray
$t\mapsto ty$ is a curve of $\tilde g$-length $|y|$
(Lemma~\ref{lem:pullback-ball}(1)), so that $d(0,y)\le|y|$, we
have $d(0,x_j)\le|x_j|\le1/256\le0.005$, so $F(0)\le
i\,(0.005)^2$; while for $|y|=1/16$ the radial estimate gives
$d(x_j,y)\ge|y|-|x_j|\ge1/16-1/256\ge 0.058$, so
$F(y)\ge i\,(0.058)^2>F(0)$: the minimum is attained in the open
ball $B(0,1/16)$. The minimum point lies close to the orbit: since
$d(x_j,y^*)\ge d(0,y^*)-d(0,x_j)\ge d(0,y^*)-0.005$ for every
$j$, the inequality
$i\left(d(0,y^*)-0.005\right)^2\le F(y^*)\le F(0)\le
i\,(0.005)^2$ (valid whenever $d(0,y^*)\ge0.005$) forces
$d(0,y^*)\le0.01$; in particular $|y^*|\le0.01$ by the radial
estimate. By (\ref{deckisometry}) (applicable: the points
$x_{j}$ and any $y\in\overline{B(0,1/16)}$ lie in
$\overline{B(0,1/7)}$ at mutual distances
$d(x_j,y)\le d(x_j,0)+d(0,y)\le|x_j|+|y|
\le1/256+1/16\le0.2$, using $d(0,\cdot)\le|\cdot|$ as above) and
the permutation property,
$$F(\pi(y))=\sum_jd\left(x_j,\pi(y)\right)^2
=\sum_jd\left(\pi\left(\pi^{-1}(x_j)\right),\pi(y)\right)^2
=\sum_jd\left(\pi^{-1}(x_j),y\right)^2=F(y)$$
for $y\in\overline{B(0,1/16)}$ with $\pi(y)\in
\overline{B(0,1/16)}$ (the inverse $\pi^{-1}=\pi_{\bar\gamma}$
permutes $O$ as well). Now
$d(0,\pi(y^*))\le d(0,x_1)+d(x_1,\pi(y^*))
=d(0,x_1)+d(x_0,y^*)\le0.005+0.01\le0.015$ (using
(\ref{deckisometry}) and $\pi(x_0)=x_1$), so
$|\pi(y^*)|\le0.015$ and $\pi(y^*)$ lies in the open ball
$B(0,1/16)$ as well, with $F(\pi(y^*))=F(y^*)=m$: it is also a
minimum point.

If $y^*\not=\pi(y^*)$, note first that both points lie in
$\overline{B(0,1/32)}$ (their norms are $\le0.015<1/32$), so by
Lemma~\ref{lem:pullback-unique-minimizers}(2) there is a
minimizing geodesic $\sigma\colon[0,1]\to B$ from $y^*$ to
$\pi(y^*)$, of length $d(y^*,\pi(y^*))\le
d(y^*,0)+d(0,\pi(y^*))\le0.01+0.015=0.025$. Every point of
$\sigma$ satisfies $d(y^*,\sigma(t))\le0.025$, hence
$d(0,\sigma(t))\le0.035$ and, by the radial estimate,
$|\sigma(t)|\le0.035<1/16$: the whole of $\sigma$ lies in the open
ball $B(0,1/16)$, the interior of the domain over which $F$ was
minimized. Moreover, for each $j$,
$d\left(x_j,\sigma(t)\right)\le d(x_j,0)+d(0,\sigma(t))
\le0.005+0.035\le1/10$, so by the first clause of
Lemma~\ref{lem:pullback-unique-minimizers}(2) (applicable, as
$x_j\in\overline{B(0,1/256)}\subset\overline{B(0,1/32)}$) we get
$\sigma(t)\in\Omega_{x_j}$, on which $d_{x_j}^2$ is smooth with
$\Hess(d_{x_j}^2)\ge\tilde g$
(Lemma~\ref{lem:pullback-unique-minimizers}(2),(3)). Hence $t\mapsto
F(\sigma(t))$ is smooth with second derivative
$\sum_j\Hess(d_{x_j}^2)(\sigma',\sigma')\ge
i\,|\sigma'|^2>0$ ($\sigma$ is non-constant), so it is strictly
convex on $[0,1]$; consequently, for $0<t<1$,
$F(\sigma(t))<\max\left(F(\sigma(0)),F(\sigma(1))\right)=m$. But
$\sigma(t)\in B(0,1/16)$, where $F\ge m$ --- a contradiction.
Therefore $\pi(y^*)=y^*$, contradicting fixed-point freeness.
Hence the $x_m$, $0\le m\le N$, are pairwise distinct.
\end{proof}

\begin{theorem}[Volume lower bound implies injectivity radius lower bound]
\label{thm:volume-injectivity-radius}
\uses{def:injectivity-radius, def:curvature-operator-norm}
\notready
\label{volinj}

Fix an integer $n>0$. For every $\epsilon>0$ there is $\delta>0$
depending on $n$ and $\epsilon$ such that the following holds.
Suppose that $(M^n,g)$ is a complete connected Riemannian
manifold of dimension $n$ and that $p\in M$. Suppose that $|\Rm(x)|\le
r^{-2}$ for all $x\in B(p,r)$. If $Vol(B(p,r))\ge \epsilon r^n$ then
the injectivity radius of $M$ at $p$ is at least $\delta r$.
\end{theorem}

\begin{proof}
\sketch
\uses{def:curvature-operator-norm,thm:hopf-rinow,lem:conjugate-sturm,lem:klingenberg-loop,lem:pullback-ball,lem:pullback-unique-minimizers,lem:loop-multiples-distinct,lem:volume-element-comparison,lem:geodesic-polar-form,thm:bishop-gromov,def:injectivity-radius,def:ricci-curvature,lem:metric-rescaling,lem:exponential-differential-jacobi,lem:model-polar-isometry}
\source{cheeger-gromov-taylor:page-0032,cheeger-gromov-taylor:page-0033}
Replacing $g$ by $r^{-2}g$, which by
Lemma~\ref{lem:metric-rescaling} (with $c=r^{-2}$) multiplies
$|\Rm|$ by $r^{2}$,
distances by $r^{-1}$, volumes by $r^{-n}$ and the injectivity
radius by $r^{-1}$, while preserving completeness, we may assume
$r=1$: thus $M$ is
complete, $|\Rm|\le1$ on $B(p,1)$, and $\Vol B(p,1)\ge\epsilon$; we
must bound $\inj_M(p)$ below by a constant $\delta(n,\epsilon)>0$.

Throughout, by Definition~\ref{def:curvature-operator-norm} the
hypothesis gives $|K(P)|\le1$ for all $2$-planes tangent at points
of $B(p,1)$, hence also $\Ric\ge-(n-1)$ there (expressing
$\Ric(X,X)$ as a sum of $n-1$ sectional curvatures in an
orthonormal basis). Every unit-speed geodesic from $p$
restricted to $[0,T]$, $T<1$, lies in
$\overline{B(p,T)}\subset B(p,1)$, where $K(P)\le1$; applying
Lemma~\ref{lem:conjugate-sturm} to it (with $K=1$, and
$T<1<\pi=\pi/\sqrt{1}$)
shows: \emph{no geodesic from $p$ has a conjugate point at
distance less than $1$}. By
completeness and the Hopf--Rinow theorem
(Theorem~\ref{thm:hopf-rinow}(1)), $\exp_p$ is defined
on all of $T_pM$, and by
Lemma~\ref{lem:exponential-differential-jacobi} it is
non-singular on $B(0,1)$ (a singular point $v$, $|v|<1$, would
produce a conjugate point at distance $|v|<1$ along
$\gamma_v$). Thus
Lemma~\ref{lem:pullback-ball} applies with $\rho_*=1$; the
pulled-back metric $\tilde g=\exp_p^*g$ on $B=B(0,1)$
satisfies $|K_{\tilde g}(P)|\le1$ (the curvature tensor of a
pullback under a local isometry is the pullback of the curvature
tensor), so Lemmas~\ref{lem:pullback-unique-minimizers} and
\ref{lem:loop-multiples-distinct} apply as well.

If $\inj_M(p)\ge\tfrac{1}{1024}$, we are done with
$\delta=\tfrac1{1024}$. So suppose
$\ell:=\inj_M(p)<\tfrac1{1024}$. Since there are no conjugate
points at distance $\le\ell<1$ along geodesics from $p$,
Lemma~\ref{lem:klingenberg-loop} provides a geodesic loop $\gamma$
at $p$ of length exactly $2\ell>0$, smooth on the open parameter
interval.

\emph{Counting preimages.} Let
$$N:=\left\lfloor\frac{1}{512\,\ell}\right\rfloor\ \ge\ 2,
\qquad\text{so that}\qquad 2\ell N\le\frac{1}{256}.$$
By Lemma~\ref{lem:loop-multiples-distinct} the points
$x_m=E(\gamma^m)$, $0\le m\le N$, are $N+1$ distinct points of
$\overline{B(0,1/256)}$ with $\exp_p(x_m)=p$. Now fix any
$q\in B(p,s)$ with $s:=\tfrac1{256}$, and let $\sigma$ be a
minimizing geodesic from $p$ to $q$, which exists by completeness
and the Hopf--Rinow theorem (Theorem~\ref{thm:hopf-rinow}(2)),
of length $<s$. The points
$$y_m:=E\left(\gamma^m\cup\sigma\right),\qquad 0\le m\le N,$$
satisfy $\exp_p(y_m)=q$ and $|y_m|\le2\ell m+s\le\tfrac1{128}$
(Lemma~\ref{lem:pullback-ball}(2)), and they are pairwise
distinct: if $y_m=y_{m'}$ then cancellation of the common tail
$\sigma$ (Lemma~\ref{lem:pullback-ball}(4b)) gives
$x_m=x_{m'}$, so $m=m'$. Hence \emph{every point of $B(p,s)$ has
at least $N+1$ preimages under $\exp_p$ in
$\overline{B(0,1/128)}$}.

\emph{Volume transfer.} Since $\exp_p$ is a local
diffeomorphism on $B(0,1)$, the open set $B(0,1/64)$ can be written
as a countable disjoint union of Borel sets on each of which
$\exp_p$ is injective; summing the change of variables formula
over the pieces gives
$$\int_M\#\left(\exp_p^{-1}(q)\cap
B(0,1/64)\right)d\vol(q)
=\Vol_{\tilde g}\left(B(0,1/64)\right),$$
where $\Vol_{\tilde g}$ denotes the volume for the pulled-back
metric. The left side is at least $(N+1)\Vol B(p,s)$ by the
previous paragraph (note $\overline{B(0,1/128)}\subset B(0,1/64)$).
For the right side, along every ray the hypotheses of
Lemma~\ref{lem:volume-element-comparison} hold with $k=1$ and
$r_0=1$ (no conjugate points, $\Ric\ge-(n-1)$), so the volume
element satisfies $\lambda(t,v)\le\sn_1^{n-1}(t)$ and, integrating
in polar coordinates (Lemma~\ref{lem:geodesic-polar-form}(4)),
$$\Vol_{\tilde g}\left(B(0,1/64)\right)
\le\omega_{n-1}\int_0^{1/64}\sn_1^{n-1}(t)\,dt=:V_1(1/64),$$
a constant depending only on $n$. Hence
\begin{equation}\label{volinjcount}
(N+1)\,\Vol B(p,s)\ \le\ V_1(1/64).
\end{equation}

\emph{The volume lower bound.} Closed balls in the complete
manifold $M$ are compact, by the Hopf--Rinow theorem
(Theorem~\ref{thm:hopf-rinow}(3)). Thus
Theorem~\ref{thm:bishop-gromov} applies on $B(p,R)$ for every
$R\le1$, with $k=1$: for $s<r'<1$,
$$\frac{\Vol B(p,s)}{V_1(s)}\ \ge\ \frac{\Vol B(p,r')}{V_1(r')},$$
where $V_1(t)=\omega_{n-1}\int_0^t\sn_1^{n-1}$ is the volume of
the $t$-ball in the model $H^n_1$ (computed via
Lemma~\ref{lem:model-polar-isometry}, as in the proof of
Theorem~\ref{thm:bishop-gromov}). Letting $r'\to1$ and using
$\Vol B(p,r')\to\Vol B(p,1)\ge\epsilon$,
$$\Vol B(p,s)\ \ge\ \frac{V_1(s)}{V_1(1)}\,\epsilon
\ =:\ c_1(n)\,\epsilon.$$

\emph{Conclusion.} Combining this with (\ref{volinjcount}),
$$\frac{1}{512\,\ell}\ \le\ N+1\ \le\
\frac{V_1(1/64)}{c_1(n)\,\epsilon},$$
so $\ell\ge\delta_0(n,\epsilon):=
\dfrac{c_1(n)\,\epsilon}{512\,V_1(1/64)}$. Taking
$\delta=\min\left(\tfrac1{1024},\ \delta_0(n,\epsilon)\right)$
proves the theorem (after undoing the initial rescaling, which
multiplies both $\inj_M(p)$ and the required bound by $r$).
\end{proof}
